\chapter{耶利米書总结}

\section{耶利米简介}
\subsection{耶利米出身}
耶利米名字的意思是「我为耶和华所高举」。\\
故乡亚拿突，属便雅悯地的祭司城(书 21:18、代上 6:60)，靠近耶路撒冷。\\
父亲名叫希勒家(耶 1:1)，叔叔叫沙龙(耶 32:7)。他有可能是被所罗门革职的祭司亚比亚他的后代(王上 2:26-27)

\subsection{耶利米时代背景}






\subsubsection{约西亚王（640－609 BC，31年）}
\begin{table}[H]
\centering
\begin{tabular}{p{1.4cm}|p{13.40cm}|p{2.20cm}}\hline\hline
\multicolumn{1}{c|}{时间}&\multicolumn{1}{c|}{内容}&经文\\\hline
8岁&约西亚登基的时候年八岁，在耶路撒冷做王三十一年&王下22:1 \\\hline
16岁&他执政第八年，虽然还年幼，却开始寻求他祖先大卫的上帝。&代下34:3\\\hline
20岁&他执政第十二年，开始清除犹大和耶路撒冷的丘坛、亚舍拉神像、雕刻和铸造的偶像。&代下34:3\\\hline
21岁&猶大王亞們的兒子約西亞在位十三年，耶和華的話臨到耶利米。&\textcolor{red}{耶1:2}\\\hline
26岁&约西亚王十八年，王差遣米书兰的孙子、亚萨利雅的儿子书记沙番上耶和华殿去，吩咐他说： 4 “你去见大祭司希勒家，使他将奉到耶和华殿的银子，就是守门的从民中收聚的银子，数算数算， 5 交给耶和华殿里办事的人。使他们转交耶和华殿里做工的人，好修理殿的破坏之处， 6 就是转交木匠和工人并瓦匠，又买木料和凿成的石头修理殿宇。 ...”&王下22:3-6 /代下34\\\hline
&大祭司希勒家对书记沙番说：“我在耶和华殿里得了律法书。”...王听见律法书上的话，便撕裂衣服，吩咐祭司希勒家与沙番的儿子亚希甘、米该亚的儿子亚革波、书记沙番和王的臣仆亚撒雅说： “你们去，为我，为民，为犹大众人，以这书上的话求问耶和华。&王下22:8-13\\\hline
      &约西亚王十八年在耶路撒冷向耶和华守这逾越节。&王下22:23\\\hline
不详&在位的时候,责备犹大假意归向上帝&\textcolor{red}{耶3:6-10}\\\hline
39岁&西亚王在战场上中箭受伤...送到耶路撒冷。约西亚死在那里&代下35:20-24\\\hline
&耶利米还为此特地作了一首哀歌，追悼约西亚王&代下35:25\\\hline
\end{tabular}
\end{table}



\subsubsection{约哈斯王（609，3个月）}
约西亚王一死，约哈斯就被国民拥立为王(王下 23:30)\\
但是法老尼哥将他锁禁 ，不许他作王。后来他又被带到埃及，死在埃及，不得再回故土(王下 23:33-34)。

\subsubsection{约雅敬王（608－597，11年）}
约雅敬原名以利亚敬，法老王尼哥废约哈斯王之后，就立约哈斯的哥哥以利亚敬为王，并且把他改名叫约雅敬。约雅敬效忠尼哥，向国民索取金银给尼哥 。宗教上，他回复玛拿西时代的混合主义，背叛耶和华。

\subsubsection{约雅斤王（597，3个月）}
约雅敬的儿子约雅斤王在巴比伦大军围城之际就任，三个月后他就出城投降巴比伦。尼布甲尼撒王进城掠夺圣殿、宫殿的财宝，又将王、官员、国中的上流领导阶级、工艺技术人员都掳到巴比伦去。祭司以西结就在这次被掳。不过被掳的人在巴比伦仍享有相当程度的自由，而约雅斤王也受到相当的礼遇(王下25:27-30)。

\subsubsection{西底家王（597－586 BC，11年）}
尼布甲尼撒立约雅斤的叔叔玛探雅作犹大王，给他改名叫西底家。此时，耶路撒冷尚未被毁灭，祖国仍然存在，被掳的人心存希望能够归回。但是西底家背叛巴比伦，犹大亡在他的手中，他自己也被巴比伦军所辱后被掳。

\begin{table}[H]
\centering
\begin{tabular}{p{3cm}|p{1.40cm}|p{10.40cm}}\hline\hline
\multicolumn{1}{c|}{时间}&\multicolumn{1}{c|}{内容}&经文\\\hline
%(一) 约西亚王年间
约西亚在位13年&1:2&耶和华的话临到耶利米\\\hline
约西亚在位时&3:6-10&责备犹大假意归向上帝\\\hline\hline
%(二)约哈斯年间
约哈斯不详&22:10-12&论约哈斯王的结局\\\hline\hline
%(三) 约雅敬年间
约雅敬登基&26 章&耶利米在圣殿劝民悔改，差点遇害\\\hline
第四年&25 章&预言犹大被掳七十年\\\hline
第四年&36:1-8&耶利米吩咐巴录写下书卷\\\hline
第四年&45:1-5&耶利米吩咐巴录写下书卷并且安慰巴录\\\hline
第四年&46:2&巴比伦王在迦基米施打败法老尼哥的军队\\\hline
第五年&36:9-28&约雅敬王焚书卷，耶利米再写\\\hline
第五年&36:28-31&论约雅敬王的灾祸\\\hline
不　详&35 章&用利甲族之例勉以色列人\\\hline
不　详&22:18-19&论约雅敬王的结局\\\hline\hline
%(四) 约雅斤(哥尼雅)
约雅斤不详&22:24-30&论到哥尼雅的遭遇\\\hline
不　详&24 章&以两筐无花果形容被掳及留下的犹大人\\\hline\hline
%(五) 西底家年间
西底家登基&27 章&预言列国要服事巴比伦\\\hline
登基时&37:1-2&王不听上帝藉耶利米所说的话\\\hline
登基时&49:34&论以拦必受惩罚\\\hline
登基时&52:1-3&西底家王行恶\\\hline
第四年&28 章&假先知哈拿尼雅说巴比伦的轭要折断\\\hline
第四年&51:59-64&论巴比伦灭亡的预言\\\hline
第十年&32:1-5&论西底家王的结局\\\hline
第十年&32:6-9&耶利米买地\\\hline
第九到&39:1-10&巴比伦进攻，犹大覆亡\\\hline
十一年&52:4-11&\\\hline
未注明&21:1-14&王差巴施户珥和西番雅去求问耶利米\\\hline
未注明&24:8-10&论到西底家王的结局\\\hline
未注明&34:1-7&巴比伦进攻时论耶路撒冷和王的结局\\\hline
未注明&34:8-22&责备君民立约宣告自由却又反悔\\\hline
未注明&37:1-21&王差犹甲和西番雅去求问耶利米，王求问先知、保护先知\\\hline
未注明&38:1-28&王暗询耶利米当如何行\\\hline\hline
\end{tabular}
\end{table}

\subsubsection{基大利统治犹大，约哈难40-44}

\subsection{犹大状况}
上帝责备犹大说，背道的以色列比奸诈的犹大还显为义(耶 3:11)。\\
他们假意归 向上帝(耶 3:10)，\\
即使指着耶和华起誓，也是假的(耶 5:2)。\\
实际的状况是，祭司与传讲律法的都不 认识上帝(耶 2:8)，\\
百姓也离弃上帝(耶2:13)，\\
他们所拜偶像的数目与犹大城的数目 相等(耶 2:28)，\\
耶路撒冷没有一个人行公义、求诚实(耶 5:1)。\\

\subsection{蒙召目的}
耶利米蒙召传讲上帝的话，目的是要施行拔出、拆毁、毁坏、倾覆，又要建立、 栽植(耶 1:10、12:16-17、18:7-9、 31:28、45:4)。

\subsection{传讲的信息}
他看见犹大所行的一切，丝毫 不客气地加以指责。他内心最希望的是百姓能够悔改，如果百姓悔改，上帝就不会降 下所说的灾(耶 4:1、7:3、25:5、36:3)。\\
他其实也明白，犹大一定会灭亡，因为他 想要为百姓代求，都受到制止(耶 7:16、11:14、14:10-12)；即使有摩西、撒母耳的 代求，也没有用(耶 15:1)。\\
因此，当巴比伦大军来进攻耶路撒冷时，耶利米一再劝 犹大的君王和百姓，要投降巴比伦，以免招来更大的灾祸。\\
但是耶利米知道上帝的带领，时候将到，就是七十年之后，被掳的人要归回，那时 他们要敬拜上帝、归向上帝。\\
有一天，上帝要与百姓立新约。耶利米书 31:33，「耶和华说：『那些日子 以后，我与以色列家所立的约乃是这样：我要将我的律法放在他们裏面，写在他们心 上。我要作他们的上帝，他们要作我的子民。』」\\
这就把我们引进耶稣基督所带来的新约。在新约里，有形的圣殿不再扮演重要的 角色。上帝的律法是写在人的心版上的。约柜不再被提说(耶 3:16)，圣殿可能被毁 (耶 7:14)。我们的身子就是上帝的殿(林前 6:19)，每一个人都当献上自己的身体当 作活祭(罗 12:1)。\\


\subsection{个人经历}
\subsubsection{蒙召和个性}
未出母胎就已经被上帝分别为圣、受上帝之派作列国的先知(耶 1:4-5)。\\
他曾经以「年幼」来向上帝推托，但上帝应许要与他同在、要拯救他(耶 1:6-10)，\\
同时上帝也预告他要遭到犹大君王、首领、祭司并众民的反对(耶 1:18-19)\\
他爱上帝、爱君王、爱百姓。但是上帝给他的使 命却是向他所爱的人报告祸音，以致他的心常常难过，几乎心碎(耶 4:19-22、20:7-18)。 \\

\subsubsection{遭受逼迫九死一生}
\begin{table}[H]
\centering
\begin{tabular}{p{1.5cm}|p{1.40cm}|p{14cm}}\hline\hline
\multicolumn{1}{c|}{时间}&\multicolumn{1}{c|}{内容}&经文\\\hline
&11:18-23&耶利米的故乡也有人反对他。他们威胁耶利米说，如果再不停止说预言， 他们可能会杀掉他\\\hline
陶人製器為喻之后&18:18-23&他們就說:「來吧!我們可以設計謀害耶利米,因為我們有祭司講律法,智慧人設謀略,先知說預言,都不能斷絕。來吧!我們可以用舌頭擊打他,不要理會他的一切話。」\\\hline
毀窯匠瓦瓶為喻&19/20:1-6&[猶大王西底家在位時(主前597-586)?]祭司音麥的兒子巴施戶珥做耶和華殿的總管，聽見耶利米預言灾祸就打先知耶利米，用耶和華殿裡便雅憫高門內的枷將他枷在那裡。\\\hline
每逢講論的時候&20:7-18&我終日成為笑話，人人都戲弄我。耶和華的話終日成了我的凌辱、譏刺。我聽見了許多人的讒謗，四圍都是驚嚇。就是我知己的朋友，也都窺探我，願我跌倒\\\hline
约雅敬王登基&26&耶利米在圣殿呼吁众民悔改，否则圣殿要像示罗一样变成荒场。祭司、先知和众民要把耶利米置之于死地。犹大的首领引用先知弥迦对希西家王的忠告，救了耶利米。乌利亚被约雅敬王处死，沙番的儿子亚希甘保护耶利米\\\hline
约雅敬4-5/9&36:9-31&耶利米吩咐巴录把上帝的话写下来，并且在圣殿里念给众人听。约雅敬王听书卷时把书卷划破丢到火中焚毁，并且还要抓耶利米和巴录。\\\hline
西底家作王&27-28&耶利米在颈项上挂着轭，表示犹大及列国都要被巴比伦所管束。 假先知哈拿尼雅却折断耶利米颈项上的轭，擅自预言两年之内巴比伦的轭被除去。耶利米宣告折断木轭却换来铁轭，哈拿尼亚当年就去世。\\\hline
写信给被掳百姓&29&耶利米写信给被掳的百姓 ，劝他们安居在被掳之地，直等到七十年满。在巴比伦有亚哈和西底家扰乱百姓，叫被掳的人不要听耶利米在信上所说的话。还有示玛雅寄信给耶路撒冷的百姓和祭司西番雅，要求他们处置耶利米，不许耶利米乱说话\\\hline
西底家王&37&西底家屡次差人去问先知耶利米。被告知要服事巴比伦王，他却不听。当埃及军队前来协助，而使耶路撒冷被围攻的危机稍微解除时，耶利米被关在文士约拿单家里的囚房中。西底家王提耶利米出监移到护卫兵的院中。\\\hline
西底家王&38&示法提雅、基大利 、犹甲、巴示户珥等人不喜欢耶利米所说：耶路撒冷将被巴比伦所灭的信息，就把耶利米下在一个只有淤泥的牢狱中。古实人以伯米勒报告给西底家王，并设法营救\\\hline
基大利死后&42-43&留在犹大的百姓遭遇危机求教于耶利米。耶利米劝他们要留在犹大地，不要下到埃及。首领和百姓挟持耶利米住在埃及答比匿，耶利米预言埃及也要遭受巴比伦的攻击(43:8-13/44:29-30/46:1-12)。百姓在埃及仍然拜偶像，招来先知耶利米的责备(耶44)。\\\hline
\end{tabular}
\end{table}



\subsection{耶利米用过的比喻异象}
\begin{table}[H]
\centering
\begin{tabular}{p{3cm}|p{3.0cm}|p{11cm}}\hline\hline
\multicolumn{1}{c|}{时间}&\multicolumn{1}{c|}{内容}&经文\\\hline
杏树枝的异象&1:11-12&杏树枝在希伯来原文和「留意保守」同音，上帝用杏树枝这个双关语来让耶利米知道，上帝必留意保守，使他所说的话成就\\\hline
烧开的锅从北而倾& 1:13-14&这是表示 灾祸从北方而来\\\hline
无用的腰帶&13:1-11&敗壞猶大的驕傲和耶路撒冷的大驕傲必像這腰帶變為無用\\\hline
酒盛满坛&13:12-14&君王祭司先知並耶路撒冷的一切居民，都酩酊大醉不蒙憐憫，以致滅絕\\\hline
终身不娶&16:1～18:18&象征犹大将遭严厉的审判\\\hline
窑匠和泥土& 18:1-14&如果离恶就不降祸；如果行恶就不降福气\\\hline
打碎窑匠瓦瓶&19&用以表示上帝要击打他的百姓。\\\hline

好坏无花果&24:1-10/29:17-18&顺神者生，背神者亡\\\hline
忿怒的杯&25:15～38&神的审判不只及于犹大，外邦列国也要受祂的审判\\\hline
示罗被弃&26&若不悔改耶路撒冷也将被弃\\\hline
绳索与轭&27-28&服事巴比伦王\\\hline
先知购田&32&他日重归故土\\\hline
给利甲族喝酒&35&神使利甲族蒙福\\\hline
北国亡国&&若不悔改犹大也将灭亡\\\hline
\end{tabular}
\end{table}


\section{約西亞 1-20}
\subsection{给耶利米的启示 1/16:1-9}
\subsubsection{約西亞在位十三年奉召為先知40+年,施行拔出/拆毀/毀壞/傾覆/又建立/栽植 1:1-10}
\begin{table}[H]
\centering
\begin{tabular}{p{9cm}|p{8.40cm}}\hline\hline
\multicolumn{1}{c|}{耶利米奉召為先知}&\multicolumn{1}{c}{為施行拔出、拆毀、毀壞、傾覆，又建立、栽植}\\\hline
便雅憫地亞拿突城的祭司中,希勒家的兒子耶利米的話記在下面。猶大王亞們的兒子約西亞在位十三年,耶和華的話臨到耶利米。從猶大王約西亞的兒子約雅敬在位的時候,直到猶大王約西亞的兒子西底家在位的末年,就是十一年五月間耶路撒冷人被擄的時候,耶和華的話也常臨到耶利米。
&耶利米說：耶和華的話臨到我說： 「我未將你造在腹中，我已曉得你；你未出母胎，我已分別你為聖。我已派你做列國的先知。」 \\\hline
我就說：「主耶和華啊！我不知怎樣說，因為我是年幼的。」 
&耶和華對我說：「你不要說『我是年幼的』,因為我差遣你到誰那裡去,你都要去,我吩咐你說什麼話,你都要說。你不要懼怕他們,因為我與你同在,要拯救你。」這是耶和華說的。於是耶和華伸手按我的口對我說：「我已將當說的話傳給你。看哪,我今日立你在列邦列國之上,為要施行拔出、拆毀、毀壞、傾覆,又要建立、栽植。」
\\\hline
\end{tabular}
\end{table}

\subsubsection{杏枝/燒開的鍋從北而傾异象-神的話必成就 1:11-19}
\begin{table}[H]
\centering
\begin{tabular}{p{11.25cm}|p{5.50cm}}\hline\hline
\multicolumn{1}{c|}{杏枝,燒開的鍋從北而傾11-16}&\multicolumn{1}{c}{勉其不要驚惶，慰必與其同在17-19}\\\hline
耶和華的話又臨到我說:「耶利米你看見什麼?」我說:「我看見一根杏樹枝。」 耶和華對我說:「你看得不錯,因為我留意保守我的話使得成就。」
&\multirow{2}{5.86cm}{「所以你當束腰，起來將我所吩咐你的一切話告訴他們。不要因他們驚惶，免得我使你在他們面前驚惶。看哪，我今日使你成為堅城、鐵柱、銅牆，與全地，和猶大的君王、首領、祭司，並地上的眾民反對。他們要攻擊你，卻不能勝你，因為我與你同在，要拯救你。」這是耶和華說的。}\\\cline{1-1}
耶和華的話第二次臨到我說：「你看見什麼？」我說：「我看見一個燒開的鍋，從北而傾。」 耶和華對我說：「必有災禍從北方發出，臨到這地的一切居民。」耶和華說：「看哪，我要召北方列國的眾族。他們要來，各安座位在耶路撒冷的城門口，周圍攻擊城牆，又要攻擊猶大的一切城邑。至於這民的一切惡，就是離棄我，向別神燒香，跪拜自己手所造的，我要發出我的判語，攻擊他們。&\\\hline
\end{tabular}
\end{table}

\subsubsection{亞拿突人設計謀害先知，必因刀劍饑荒滅亡 11:18-23}
耶和華指示我，我就知道，你將他們所行的給我指明。 19 我卻像柔順的羊羔被牽到宰殺之地，我並不知道他們設計謀害我說：「我們把樹連果子都滅了吧！將他從活人之地剪除，使他的名不再被記念。」

按公義判斷，察驗人肺腑心腸的萬軍之耶和華啊，我卻要見你在他們身上報仇，因我將我的案件向你稟明了。 21 所以，耶和華論到尋索你命的亞拿突人如此說：「他們說『你不要奉耶和華的名說預言，免得你死在我們手中』， 22 所以萬軍之耶和華如此說：『看哪，我必刑罰他們，他們的少年人必被刀劍殺死，他們的兒女必因饑荒滅亡。 23 並且沒有餘剩的人留給他們，因為在追討之年，我必使災禍臨到亞拿突人。』」

%\subsection{先知與主理論，為何惡人亨通  12}
\subsubsection{先知惑惡人行詭詐得亨通；主示其被父兄用奸詐对待 12:1-4}
耶和華啊，我與你爭辯的時候，你顯為義，但有一件，我還要與你理論：惡人的道路為何亨通呢？大行詭詐的為何得安逸呢？ 2 你栽培了他們，他們也扎了根，長大，而且結果。他們的口是與你相近，心卻與你遠離。 3 耶和華啊，你曉得我，看見我，察驗我向你是怎樣的心。求你將他們拉出來，好像將宰的羊，叫他們等候殺戮的日子。 4 這地悲哀，通國的青草枯乾，要到幾時呢？因其上居民的惡行，牲畜和飛鳥都滅絕了，他們曾說：「他看不見我們的結局。」

耶和華說：「你若與步行的人同跑，尚且覺累，怎能與馬賽跑呢？你在平安之地雖然安穩，在約旦河邊的叢林要怎樣行呢？ 6 因為連你弟兄和你父家都用奸詐待你，他們也在你後邊大聲喊叫。雖向你說好話，你也不要信他們。

\subsubsection{先知自言見惡於人求眷顧,蒙主慰藉答應與其同在 15:10-21}
\begin{table}[H]
\centering
\begin{tabular}{p{8.5cm}|p{8.750cm}}\hline\hline
\multicolumn{1}{c|}{先知自言見惡於人求眷顧}&\multicolumn{1}{c}{蒙主慰藉答應與其同在}\\\hline
我的母親哪，我有禍了！因你生我作為遍地相爭相競的人。我素來沒有借貸於人，人也沒有借貸於我，人人卻都咒罵我。 
&耶和華說:「我必要堅固你,使你得好處。災禍苦難臨到的時候,我必要使仇敵央求你。「人豈能將銅與鐵,就是北方的鐵折斷呢?我必因你在四境之內所犯的一切罪,把你的貨物財寶當掠物,白白地交給仇敵。我也必使仇敵帶這掠物到你所不認識的地去,因我怒中起的火要將你們焚燒。」\\\hline
耶和華啊，你是知道的！求你記念我，眷顧我，向逼迫我的人為我報仇。不要向他們忍怒取我的命，要知道我為你的緣故受了凌辱。耶和華萬軍之神啊，我得著你的言語就當食物吃了，你的言語是我心中的歡喜快樂，因我是稱為你名下的人。我沒有坐在宴樂人的會中，也沒有歡樂。我因你的感動(手)獨自靜坐，因你使我滿心憤恨。我的痛苦為何長久不止呢？我的傷痕為何無法醫治，不能痊癒呢？難道你待我有詭詐，像流乾的河道嗎？
&耶和華如此說：「你若歸回，我就將你再帶來，使你站在我面前。你若將寶貴的和下賤的分別出來，你就可以當做我的口。他們必歸向你，你卻不可歸向他們。我必使你向這百姓成為堅固的銅牆，他們必攻擊你，卻不能勝你。因我與你同在，要拯救你，搭救你。」這是耶和華說的。「我必搭救你脫離惡人的手，救贖你脫離強暴人的手。」\\\hline
\end{tabular}
\end{table}



\subsubsection{因斯土必滅,示耶利米不可娶妻，不要為百姓悲傷 16:1-9}
\begin{table}[H]
\centering
\begin{tabular}{p{6.2cm}|p{11.cm}}\hline\hline
\multicolumn{1}{c|}{禁娶妻生兒養女}&\multicolumn{1}{c}{禁入喪家哀哭}\\\hline
耶和華的話又臨到我說：「你在這地方不可娶妻，生兒養女。因為論到在這地方所生的兒女，又論到在這國中生養他們的父母，耶和華如此說：他們必死得甚苦，無人哀哭，必不得葬埋，必在地上像糞土，必被刀劍和饑荒滅絕。他們的屍首必給空中的飛鳥和地上的野獸做食物。」
&耶和華如此說:不要進入喪家,不要去哀哭,也不要為他們悲傷,因我已將我的平安、慈愛、憐憫從這百姓奪去了。這是耶和華說的。連大帶小,都必在這地死亡,不得葬埋。人必不為他們哀哭,不用刀劃身,也不使頭光禿。他們有喪事,人必不為他們擘餅,因死人安慰他們。他們喪父喪母,人也不給他們一杯酒安慰他們。你不可進入宴樂的家,與他們同坐吃喝。因為萬軍之耶和華以色列的神如此說:你們還活著的日子,在你們眼前,我必使歡喜和快樂的聲音、新郎和新婦的聲音,從這地方止息了。\\\hline
\end{tabular}
\end{table}




\subsubsection{祈主拯救做避難所 17:14-18}
耶和華說：「離開我的，他們的名字必寫在土裡，因為他們離棄我這活水的泉源。」 

14 耶和華啊，求你醫治我，我便痊癒，拯救我，我便得救，因你是我所讚美的。
他們對我說：「耶和華的話在哪裡呢？叫這話應驗吧！」 16 至於我，那跟從你做牧人的職分，我並沒有急忙離棄，也沒有想那災殃的日子，這是你知道的。我口中所出的言語都在你面前。 17 不要使我因你驚恐，當災禍的日子，你是我的避難所。 18 願那些逼迫我的蒙羞，卻不要使我蒙羞；使他們驚惶，卻不要使我驚惶；使災禍的日子臨到他們，以加倍的毀壞毀壞他們！

\subsubsection{耶利米求主重懲陷害己命者 18:18-23}
他們就說：「來吧！我們可以設計謀害耶利米，因為我們有祭司講律法，智慧人設謀略，先知說預言，都不能斷絕。來吧！我們可以用舌頭擊打他，不要理會他的一切話。」

19 耶和華啊，求你理會我，且聽那些與我爭競之人的話。 20 豈可以惡報善呢？他們竟挖坑要害我的性命！求你記念我怎樣站在你面前為他們代求，要使你的憤怒向他們轉消。 21 故此，願你將他們的兒女交於饑荒和刀劍，願他們的妻無子且做寡婦；又願他們的男人被死亡所滅，他們的少年人在陣上被刀擊殺。 22 你使敵軍忽然臨到他們的時候，願人聽見哀聲從他們的屋內發出，因他們挖坑要捉拿我，暗設網羅要絆我的腳。 23 耶和華啊，他們要殺我的那一切計謀，你都知道。不要赦免他們的罪孽，也不要從你面前塗抹他們的罪惡，要叫他們在你面前跌倒，願你發怒的時候罰辦他們。
\subsubsection{因說預言受撻於巴施戶珥,神更其名預兆驚惶 19:14-20:6}
\begin{table}[H]
\centering
\begin{tabular}{p{7.cm}|p{10.50cm}}\hline\hline
\multicolumn{1}{c|}{因說預言受撻於巴施戶珥}&\multicolumn{1}{c}{神更其名預兆驚惶}\\\hline
耶利米從陀斐特，就是耶和華差他去說預言的地方回來，站在耶和華殿的院中，對眾人說：「萬軍之耶和華以色列的神如此說：我必使我所說的一切災禍臨到這城和屬城的一切城邑，因為他們硬著頸項不聽我的話。」
&\multirow{2}{10.76cm}{次日巴施戶珥將耶利米開枷釋放。於是耶利米對他說:「耶和華不是叫你的名為『巴施戶珥』,乃是叫你『瑪歌珥米撒畢（四面驚嚇）』。因耶和華如此說:我必使你自覺驚嚇,你也必使眾朋友驚嚇。他們必倒在仇敵的刀下,你也必親眼看見。我必將猶大人全交在巴比倫王的手中,他要將他們擄到巴比倫去,也要用刀將他們殺戮。並且我要將這城中的一切貨財和勞碌得來的,並一切珍寶,以及猶大君王所有的寶物,都交在他們仇敵的手中,仇敵要當做掠物帶到巴比倫去。你這巴施戶珥和一切住在你家中的人,都必被擄去。你和你的眾朋友,就是你向他們說假預言的,都必到巴比倫去,要死在那裡,葬在那裡。」}\\\cline{1-1}
祭司音麥的兒子巴施戶珥做耶和華殿的總管，聽見耶利米預言這些事，他就打先知耶利米，用耶和華殿裡便雅憫高門內的枷將他枷在那裡。&\\\hline
\end{tabular}
\end{table}


\subsubsection{被人譏辱却忍不住講論,密友亦加謀害主與其同在，先知自詛生辰 20:7-18}
\begin{table}[H]
\centering
\begin{tabular}{p{5.25cm}|p{6.75cm}|p{4.75cm}}\hline\hline
\multicolumn{1}{c|}{被人譏辱却忍不住講論}&\multicolumn{1}{c|}{密友亦加謀害主與其同在}&\multicolumn{1}{c}{先知自詛生辰}\\\hline
耶和華啊,你曾勸導我,我也聽了你的勸導，你比我有力量，且勝了我。我終日成為笑話，人人都戲弄我。我每逢講論的時候就發出哀聲，我喊叫說：「有強暴和毀滅！」因為耶和華的話終日成了我的凌辱、譏刺。我若說我不再提耶和華，也不再奉他的名講論，我便心裡覺得似乎有燒著的火閉塞在我骨中，我就含忍不住不能自禁。我聽見了許多人的讒謗，四圍都是驚嚇。
&就是我知己的朋友，也都窺探我，願我跌倒，說：「告他吧！我們也要告他！或者他被引誘，我們就能勝他，在他身上報仇。」然而，耶和華與我同在，好像甚可怕的勇士。因此，逼迫我的必都絆跌，不能得勝。他們必大大蒙羞，就是受永不忘記的羞辱，因為他們行事沒有智慧。試驗義人，察看人肺腑心腸的萬軍之耶和華啊，求你容我見你在他們身上報仇，因我將我的案件向你稟明了。你們要向耶和華唱歌，讚美耶和華！因他救了窮人的性命脫離惡人的手。
&願我生的那日受咒詛，願我母親產我的那日不蒙福！給我父親報信說「你得了兒子」，使我父親甚歡喜的，願那人受咒詛！願那人像耶和華所傾覆而不後悔的城邑，願他早晨聽見哀聲，晌午聽見呐喊。因他在我未出胎的時候不殺我，使我母親成了我的墳墓，胎就時常重大。我為何出胎見勞碌愁苦，使我的年日因羞愧消滅呢？\\\hline
\end{tabular}
\end{table}

\subsection{以色列/猶大}



%\subsection{跟随虚无的神，就成为虚妄的人 2}
\subsubsection{忆雅各幼年的忠诚,責列祖祭司官長先知負恩隨虛無的神 2:1-8}
耶和華的話臨到我說： 2 「你去向耶路撒冷人的耳中喊叫說：『耶和華如此說：你幼年的恩愛，婚姻的愛情，你怎樣在曠野，在未曾耕種之地跟隨我，我都記得。 3 那時以色列歸耶和華為聖，作為土產初熟的果子，凡吞吃他的必算為有罪，災禍必臨到他們。』這是耶和華說的。」

4 雅各家，以色列家的各族啊，你們當聽耶和華的話！ 5 耶和華如此說：「你們的\textcolor{red}{列祖}見我有什麼不義，竟遠離我，隨從虛無的神，自己成為虛妄的呢？ 6 他們也不說：『那領我們從埃及地上來，引導我們經過曠野，沙漠有深坑之地，和乾旱死蔭，無人經過、無人居住之地的耶和華，在哪裡呢？』 7 我領你們進入肥美之地，使你們得吃其中的果子和美物，但你們進入的時候就玷汙我的地，使我的產業成為可憎的。 8 \textcolor{red}{祭司}都不說：『耶和華在哪裡呢？』傳講律法的都不認識我，\textcolor{red}{官長}違背我，\textcolor{red}{先知}藉巴力說預言，隨從無益的神。」

\subsubsection{以色列換神并做了兩件惡事,遭災受禍因離棄神自招 2:9-19}
耶和華說：「我因此必與你們爭辯，也必與你們的子孫爭辯。 10 你們且過到基提海島去察看，打發人往基達去留心查考，看曾有這樣的事沒有！ 11 豈有一國換了他的神嗎？其實這不是神！但我的百姓將他們的榮耀換了那無益的神。 12 諸天哪，要因此驚奇，極其恐慌，甚為淒涼！」這是耶和華說的。 13 「因為我的百姓做了兩件惡事，就是離棄我這活水的泉源，為自己鑿出池子，是破裂不能存水的池子。


%\subsubsection{遭災受禍因離棄神自招 2:14-19}
以色列是僕人嗎？是家中生的奴僕嗎？為何成為掠物呢？ 15 少壯獅子向他咆哮，大聲吼叫，使他的地荒涼，城邑也都焚燒，無人居住。 16 挪弗人和答比匿人也打破你的頭頂。 17 這事臨到你身上，不是你自招的嗎？不是因耶和華你神引你行路的時候，你離棄他嗎？ 18 現今你為何在埃及路上要喝西曷的水呢？你為何在亞述路上要喝大河的水呢？ 19 你自己的惡必懲治你，你背道的事必責備你。由此可知可見，你離棄耶和華你的神，不存敬畏我的心，乃為惡事，為苦事。」這是主萬軍之耶和華說的。

\subsubsection{歷數猶大隨從眾巴力罪惡 2:20-28}
%\subsubsection{在各高岡上青翠樹下屈身行淫，用鹼洗濯罪孽仍顯 2:20-22}
「我在古時折斷你的軛，解開你的繩索，你說：『我必不侍奉耶和華！』因為你\textcolor{red}{在各高岡上、各青翠樹下屈身行淫}。 (我在古時折斷你的軛，解開你的繩索，你就說：『我必不侍奉別神！』誰知，你在各高岡上、各青翠樹下仍屈身行淫。) 21 然而我栽你是上等的葡萄樹，全然是真種子，你怎麼向我變為外邦葡萄樹的壞枝子呢？ 22 你雖用鹼，多用肥皂洗濯，你罪孽的痕跡仍然在我面前顯出。」這是主耶和華說的。 

%\subsubsection{喜愛別神的心，如同快行的獨峰駝和慾心發動的野驢 2:23-25}
23 「你怎能說『我沒有玷汙，沒有隨從眾巴力』？你看你谷中的路，就知道你所行的如何。你是快行的獨峰駝，狂奔亂走。 24 你是野驢，慣在曠野，慾心發動就吸風；起性的時候，誰能使牠轉去呢？凡尋找牠的必不至疲乏，在牠的月份必能尋見。 25 我說：『你不要使腳上無鞋，喉嚨乾渴。』你倒說：『這是枉然！我\textcolor{red}{喜愛別神}，我必隨從他們。』

%\subsubsection{遭遇患難始求救援，必因與城的數目相等的假神羞愧 2:26-37}
「賊被捉拿，怎樣羞愧，以色列家和他們的君王、首領、祭司、先知也都照樣羞愧。 27 他們向木頭說『你是我的父』，向石頭說『你是生我的』。他們以背向我，不以面向我，及至遭遇患難的時候，卻說：『起來，拯救我們！』 28 你為自己做的神在哪裡呢？你遭遇患難的時候，叫他們起來拯救你吧！猶大啊，你\textcolor{red}{神的數目與你城的數目相等}！」

\subsubsection{受神責打是徒然,必因埃及蒙羞 2:29-37}
29 耶和華說：「你們為何與我爭辯呢？你們都違背了我。 30 我責打你們的兒女是徒然的，他們不受懲治。你們自己的\textcolor{red}{刀吞滅你們的先知}，好像殘害的獅子。 


31 這世代的人哪，你們要看明耶和華的話！我豈向以色列做曠野呢？或做幽暗之地呢？我的百姓為何說『我們脫離約束，再不歸向你了』？ 32 處女豈能忘記她的裝飾呢？新婦豈能忘記她的美衣呢？我的百姓\textcolor{red}{卻忘記了我}無數的日子。 

33 你怎麼修飾你的道路要求愛情呢？就是惡劣的婦人，你也叫她們行你的路。 34 並且你的\textcolor{red}{衣襟上有無辜窮人的血}，你殺他們並不是遇見他們挖窟窿，乃是因這一切的事。 35 你還說：『我無辜，耶和華的怒氣必定向我消了。』看哪，我必審問你，因你自說『我沒有犯罪』。 

36 你為何東跑西奔，要更換你的路呢？你必\textcolor{red}{因埃及蒙羞}，像從前因亞述蒙羞一樣。 37 你也必兩手抱頭從埃及出來，因為耶和華已經棄絕你所倚靠的，你必不因他們得順利。

\subsubsection{慰犹大雖行淫作惡還可以歸向神, 但其行惡還不懼怕較以色列尤甚3:1-5}
「有話說，人若休妻，妻離他而去，做了別人的妻，前夫豈能再收回她來？若收回她來，那地豈不是大大玷汙了嗎？但你和許多親愛的行邪淫，還可以歸向我。」這是耶和華說的。

2 「你向淨光的高處舉目觀看，你在何處沒有淫行呢？你坐在道旁等候，好像阿拉伯人在曠野埋伏一樣，並且你的淫行、邪惡玷汙了全地。 3 因此甘霖停止，春(晚)雨不降，你還是有娼妓之臉，不顧羞恥。 4 從今以後，你豈不向我呼叫說：『我父啊，你是我幼年的恩主！ 5 耶和華豈永遠懷怒，存留到底嗎？』看哪，你又發惡言又行壞事，隨自己的私意而行(你雖這樣說，還是行惡，放縱慾心。)。」

\subsection{約西亞王在位的時候}
\subsubsection{吩咐先知向北方宣告，以色列悔罪歸誠仍蒙矜憫 3:6-13}
約西亞王在位的時候，耶和華又對我說：「背道的以色列所行的，你看見沒有？她上各高山，在各青翠樹下行淫。 7 她行這些事以後，我說她必歸向我，她卻不歸向我。她奸詐的妹妹猶大也看見了。 8 背道的以色列行淫，我為這緣故給她休書休她。我看見她奸詐的妹妹猶大還不懼怕，也去行淫。 9 因以色列輕忽了她的淫亂，和石頭木頭行淫，地就被玷汙了。 10 雖有這一切的事，她奸詐的妹妹猶大還不一心歸向我，不過是假意歸我。」這是耶和華說的。


耶和華對我說：「背道的以色列比奸詐的猶大還顯為義。 12 你去向北方宣告說：『耶和華說：背道的以色列啊，回來吧！我必不怒目看你們，因為我是慈愛的，我必不永遠存怒。這是耶和華說的。 13 只要承認你的罪孽，就是你違背耶和華你的神，在各青翠樹下向別神東奔西跑，沒有聽從我的話。這是耶和華說的。 

%\subsubsection{许賜牧者，猶大和以色列從北方回耶路撒冷 3:14-18}
14 耶和華說：背道的兒女啊，回來吧！因為我做你們的丈夫，並且我必將你們從一城取一人，從一族取兩人，帶到錫安。 15 我也必將合我心的牧者賜給你們，他們必以知識和智慧牧養你們。 16 耶和華說：你們在國中生養眾多，當那些日子，人必不再提說耶和華的約櫃，不追想，不記念，不覺缺少，也不再製造。 17 那時，人必稱耶路撒冷為耶和華的寶座，萬國必到耶路撒冷，在耶和華立名的地方聚集。他們必不再隨從自己頑梗的惡心行事。 18 當那些日子，猶大家要和以色列家同行，從北方之地一同來到我賜給你們列祖為業之地。


『我說，我怎樣將你安置在兒女之中，賜給你美地，就是萬國中肥美的產業！我又說，你們必稱我為父，也不再轉去不跟從我。 20 以色列家，你們向我行詭詐，真像妻子行詭詐離開她丈夫一樣。這是耶和華說的。 21 在淨光的高處聽見人聲，就是以色列人哭泣懇求之聲，乃因他們走彎曲之道，忘記耶和華他們的神。 22 你們這背道的兒女啊，回來吧！我要醫治你們背道的病。』」

\subsubsection{人的幫助是枉然，偶像吞吃勞碌所得 3:22b-25+4:2}
「看哪！我們來到你這裡，因你是耶和華我們的神。 23 仰望從小山或從大山的喧嚷中得幫助，真是枉然的，以色列得救誠然在乎耶和華我們的神。 24 從我們幼年以來，那可恥的偶像將我們列祖所勞碌得來的羊群、牛群和他們的兒女都吞吃了。 25 我們在羞恥中躺臥吧，願慚愧將我們遮蓋！因為從立國 (幼年)以來，我們和我們的列祖常常得罪耶和華我們的神，沒有聽從耶和華我們神的話。」

耶和華說：「以色列啊，你若回來歸向我，若從我眼前除掉你可憎的偶像，你就不被遷移。 2 你必憑誠實、公平、公義指著永生的耶和華起誓，列國必因耶和華稱自己為有福，也必因他誇耀。」

\subsubsection{勸猶大悔改，大毀滅從北方已經動身 4:3-9}
耶和華對猶大和耶路撒冷人如此說：「要開墾你們的荒地，不要撒種在荊棘中。 4 猶大人和耶路撒冷的居民哪，你們當自行割禮，歸耶和華，將心裡的汙穢除掉，恐怕我的憤怒因你們的惡行發作，如火著起，甚至無人能以熄滅。

5 「你們當傳揚在猶大，宣告在耶路撒冷說：『你們當在國中吹角，高聲呼叫說：你們當聚集，我們好進入堅固城！』 6 應當向錫安豎立大旗！要逃避，不要遲延！因我必使災禍與大毀滅從北方來到。」 7 有獅子從密林中上來，是毀壞列國的。牠已經動身出離本處，要使你的地荒涼，使你的城邑變為荒場無人居住。 8 因此你們當腰束麻布，大聲哀號，因為耶和華的烈怒沒有向我們轉消。 9 耶和華說：「到那時，君王和首領的心都要消滅，祭司都要驚奇，先知都要詫異。」

\subsubsection{大聲傳給列國，猶大罪惡的結果是仇敵必如雲上來 4:10-18}
10 我說：「哀哉！主耶和華啊，你真是大大地欺哄這百姓和耶路撒冷，說『你們必得平安』，其實刀劍害及性命了！」

那時，必有話對這百姓和耶路撒冷說：「有一陣熱風從曠野淨光的高處向我的眾民(民女)颳來，不是為簸揚，也不是為揚淨， 12 必有一陣更大的風從這些地方為我颳來。現在我又必發出判語，攻擊他們。」 13 看哪！仇敵必如雲上來，他的戰車如旋風，他的馬匹比鷹更快。我們有禍了！我們敗落了！ 14 耶路撒冷啊，你當洗去心中的惡，使你可以得救。惡念存在你心裡要到幾時呢？ 15 有聲音從但傳揚，從以法蓮山報禍患。 16 「你們當傳給列國，報告攻擊耶路撒冷的事說：有探望的人從遠方來到，向猶大的城邑大聲呐喊。 17 他們周圍攻擊耶路撒冷好像看守田園的，因為她背叛了我。這是耶和華說的。 18 你的行動、你的作為招惹這事，這是你罪惡的結果，實在是苦，是害及你心了！」

\subsubsection{先知哀嘆猶大遭災地必荒涼，人因馬兵惊慌,神說百姓愚頑 4:19-31}
我的肺腑啊，我的肺腑啊，我心疼痛！我心在我裡面煩躁不安，我不能靜默不言，因為我已經聽見角聲和打仗的喊聲。 20 毀壞的信息連絡不絕，因為全地荒廢。我的帳篷忽然毀壞，我的幔子頃刻破裂。 21 我看見大旗、聽見角聲要到幾時呢？ 22 耶和華說：「我的百姓愚頑，不認識我。他們是愚昧無知的兒女，有智慧行惡，沒有知識行善。」

%\subsubsection{全地必然荒涼，人因馬兵惊慌 4:23-31}
23 先知說：我觀看地，不料地是空虛混沌。我觀看天，天也無光。 24 我觀看大山，不料盡都震動，小山也都搖來搖去。 25 我觀看，不料無人，空中的飛鳥也都躲避。 26 我觀看，不料肥田變為荒地，一切城邑在耶和華面前因他的烈怒都被拆毀。

27 耶和華如此說：「全地必然荒涼，我卻不毀滅淨盡。 28 因此地要悲哀，在上的天也必黑暗。因為我言已出，我意已定，必不後悔，也不轉意不做。」 29 各城的人因馬兵和弓箭手的響聲就都逃跑，進入密林，爬上磐石。各城被撇下，無人住在其中。 30 你淒涼的時候要怎樣行呢？你雖穿上朱紅衣服，佩戴黃金裝飾，用顏料修飾眼目，這樣標緻是枉然的！戀愛你的藐視你，並且尋索你的性命。 31 我聽見有聲音，彷彿婦人產難的聲音，好像生頭胎疼痛的聲音，是錫安女子(民)的聲音。她喘著氣，挓挲手，說：「我有禍了！在殺人的跟前，我的心發昏了！」

%\subsection{\textcolor{red}{三愚}民的結局——无知的人怎樣侍奉外邦神，必在外邦地侍奉外邦人 5}
\subsubsection{錫安貧窮民尊大人\textcolor{red}{愚昧}不曉得神的法則 5:1-6}
「你們當在耶路撒冷的街上跑來跑去，在寬闊處尋找，看看有一人行公義、求誠實沒有！若有，我就赦免這城。 2 其中的人雖然指著永生的耶和華起誓，所起的誓實在是假的。」 3 耶和華啊，你的眼目不是看顧誠實嗎？你擊打他們，他們卻不傷慟；你毀滅他們，他們仍不受懲治。他們使臉剛硬過於磐石，不肯回頭。

4 我說：「這些人實在是貧窮的，是愚昧的，因為不曉得耶和華的作為和他們神的法則。 5 我要去見尊大的人，對他們說話，因為他們曉得耶和華的作為和他們神的法則。」哪知，這些人齊心將軛折斷，掙開繩索。 6 因此，林中的獅子必害死他們，晚上(野地)的豺狼必滅絕他們，豹子要在城外窺伺他們。凡出城的必被撕碎，因為他們的罪過極多，背道的事也加增了。

\subsubsection{棄真奉假神必被討罪,\textcolor{red}{民愚}不認耶和華不信先知降災警告 5:7-13}
「我怎能赦免你呢？你的兒女離棄我，又指著那不是神的起誓。我使他們飽足，他們就行姦淫，成群地聚集在娼妓家裡。 8 他們像餵飽的馬，到處亂跑，各向他鄰舍的妻發嘶聲。」 9 耶和華說：「我豈不因這些事討罪呢？豈不報復這樣的國民呢？」


「你們要上她葡萄園的牆施行毀壞，但不可毀壞淨盡，只可除掉她的枝子，因為不屬耶和華。 11 原來以色列家和猶大家大行詭詐攻擊我。」這是耶和華說的。 
12 他們不認耶和華，說：「這並不是他，災禍必不臨到我們，刀劍和饑荒我們也看不見。 13 先知的話必成為風，道也不在他們裡面，這災必臨到他們身上。」

\subsubsection{神差遠方國來攻擊却不毀滅淨盡，\textcolor{red}{民愚}有耳不聽也不棄惡得福 5:14-31}
14 所以耶和華萬軍之神如此說：「因為百姓說這話，我必使我的話在你口中為火，使他們為柴，這火便將他們燒滅。」 15 耶和華說：「以色列家啊，我必使一國的民從遠方來攻擊你，是強盛的國，是從古而有的國。他們的言語你不曉得，他們的話你不明白。 16 他們的箭袋是敞開的墳墓，他們都是勇士。 17 他們必吃盡你的莊稼和你的糧食，是你兒女該吃的；必吃盡你的牛羊，吃盡你的葡萄和無花果。又必用刀毀壞你所倚靠的堅固城。」 18 耶和華說：「就是到那時，我也不將你們毀滅淨盡。」


%\subsubsection{怎樣侍奉外邦神，必在外邦地侍奉外邦人 5:19}
百姓若說：『耶和華我們的神為什麼向我們行這一切事呢？』你就對他們說：『你們怎樣離棄耶和華(我)，在你們的地上侍奉外邦神，也必照樣在不屬你們的地上侍奉外邦人。』

%\subsubsection{神定沙為海的界，百姓行不義罪孽充盈却不戰兢 5:20-29}
「當傳揚在雅各家，報告在猶大說： 21 愚昧無知的百姓啊，你們有眼不看，有耳不聽，現在當聽這話！ 22 耶和華說：你們怎麼不懼怕我呢？我以永遠的定例，用沙為海的界限，水不得越過，因此你們在我面前還不戰兢嗎？波浪雖然翻騰，卻不能逾越；雖然砰訇，卻不能過去。 23 但這百姓有背叛忤逆的心，他們叛我而去。 24 心內也不說：『我們應當敬畏耶和華我們的神，他按時賜雨，就是秋雨春雨，又為我們定收割的節令，永存不廢。』 25 你們的罪孽使這些事轉離你們，你們的罪惡使你們不能得福。 26 因為在我民中有惡人，他們埋伏窺探，好像捕鳥的人，他們設立圈套陷害人。 27 籠內怎樣滿了雀鳥，他們的房中也照樣充滿詭詐，所以他們得成為大，而且富足。 28 他們肥胖光潤，作惡過甚，不為人申冤，就是不為孤兒申冤，不使他亨通，也不為窮人辨屈。 29 耶和華說：我豈不因這些事討罪呢？豈不報復這樣的國民呢？」

%\subsubsection{假先知祭司偽言欺民得民喜愛 5:30-31}
「國中有可驚駭、可憎惡的事， 31 就是先知說假預言，祭司藉他們把持權柄，我的百姓也喜愛這些事。到了結局，你們怎樣行呢？

%\subsection{悖逆的人\textcolor{red}{不聽}訓誨，北方之民\textcolor{red}{聽旨}来讨罪 6}
\subsubsection{預言猶大之災禍從北方來，民耳未受割禮不能聽，以神的話為羞辱 6:1-15}
便雅憫人哪，你們要逃出耶路撒冷！在提哥亞吹角，在伯哈基琳立號旗，因為有災禍與大毀滅從北方張望。 2 那秀美嬌嫩的錫安女子(民)，我必剪除。 3 牧人必引他們的羊群到她那裡，在她周圍支搭帳篷，各在自己所占之地使羊吃草。

%\subsubsection{命敵來攻乃因其罪 6:4-7}
你們要準備攻擊她！」「起來吧，我們可以趁午時上去！哀哉，日已漸斜，晚影拖長了！ 5 起來吧，我們夜間上去，毀壞她的宮殿！」 6 因為萬軍之耶和華如此說：「你們要砍伐樹木，築壘攻打耶路撒冷。這就是那該罰的城，其中盡是欺壓。 7 井怎樣湧出水來，這城也照樣湧出惡來，在其間常聽見有強暴毀滅的事，病患損傷也常在我面前。

 8 耶路撒冷啊，你當受教，免得我心與你生疏，免得我使你荒涼，成為無人居住之地。
萬軍之耶和華曾如此說：「敵人必擄盡以色列剩下的民，如同摘淨葡萄一樣。你要像摘葡萄的人摘了又摘，回手放在筐子裡。」

 10 現在我可以向誰說話作見證，使他們聽呢？他們的耳朵未受割禮，不能聽見。看哪，耶和華的話他們以為羞辱，不以為喜悅。 11 因此我被耶和華的憤怒充滿，難以含忍。

「我要傾在街中的孩童和聚會的少年人身上，連夫帶妻，並年老的與日子滿足的都必被擒拿。 12 他們的房屋、田地和妻子都必轉歸別人，我要伸手攻擊這地的居民。」這是耶和華說的。 



13 「因為他們從最小的到至大的都一味地貪婪，從先知到祭司都行事虛謊。 14 他們輕輕忽忽地醫治我百姓的損傷，說『平安了！平安了！』，其實沒有平安。 15 他們行可憎的事知道慚愧嗎？不然！他們毫不慚愧，也不知羞恥。因此，他們必在仆倒的人中仆倒，我向他們討罪的時候，他們必致跌倒。」這是耶和華說的。


\subsubsection{百姓不聽守望人角聲，必一同跌在絆腳石上 6:16-30}
耶和華如此說：「你們當站在路上察看，訪問古道，哪是善道，便行在其間，這樣你們心裡必得安息。他們卻說：『我們不行在其間。』 17 我設立守望的人照管你們，說：『要聽角聲。』他們卻說：『我們不聽。』 18 列國啊，因此你們當聽！會眾啊，要知道他們必遭遇的事！ 19 地啊，當聽！我必使災禍臨到這百姓，就是他們意念所結的果子。因為他們不聽從我的言語，至於我的訓誨(律法)，他們也厭棄了。 20 從示巴出的乳香，從遠方出的菖蒲(甘蔗)，奉來給我有何益呢？你們的燔祭不蒙悅納，你們的平安祭我也不喜悅。」 21 所以耶和華如此說：「我要將絆腳石放在這百姓前面，父親和兒子要一同跌在其上，鄰舍與朋友也都滅亡。」


%\subsubsection{北方大國從地極來，勸民因受災哀痛己罪 6:22-26}
耶和華如此說：「看哪，有一種民從北方而來，並有一大國被激動，從地極來到。 23 他們拿弓和槍，性情殘忍，不施憐憫，他們的聲音像海浪砰訇。錫安城(女子)啊，他們騎馬都擺隊伍，如上戰場的人要攻擊你。」 24 我們聽見他們的風聲，手就發軟；痛苦將我們抓住，疼痛彷彿產難的婦人。 25 你們不要往田野去，也不要行在路上，因四圍有仇敵的刀劍和驚嚇。 26 我民(民女)哪，應當腰束麻布，滾在灰中！你要悲傷，如喪獨生子痛痛哭號，因為滅命的要忽然臨到我們。

%\subsubsection{被煉而又煉終是徒然,稱為被棄的銀渣 6:27-30}
27 「我使你在我民中為高臺(試驗人的)，為保障，使你知道試驗他們的行動。 28 他們都是極悖逆的，往來讒謗人。他們是銅是鐵，都行壞事。 29 風箱吹火，鉛被燒毀，他們煉而又煉，終是徒然，因為惡劣的還未除掉。 30 人必稱他們為被棄的銀渣，因為耶和華已經棄掉他們。」



%\subsection{代求声止息，悔改声响起 7}
\subsubsection{勸猶大悔罪免擄至異邦,细数猶大眾罪 7:1-11}
\begin{table}[H]
\centering
\begin{tabular}{p{10.75cm}|p{6.50cm}}\hline\hline
\multicolumn{1}{c|}{勸猶大悔罪免擄至異邦1-7}&\multicolumn{1}{c}{细数猶大眾罪8-11}\\\hline
耶和華的話臨到耶利米說： 「你當站在耶和華殿的門口，在那裡宣傳這話說：『你們進這些門敬拜耶和華的一切猶大人，當聽耶和華的話！萬軍之耶和華以色列的神如此說：你們改正行動作為，我就使你們在這地方仍然居住。你們\textcolor{red}{不要倚靠}虛謊的話說：「這些是耶和華的殿，是耶和華的殿，是耶和華的殿。」你們若實在改正行動作為，在人和鄰舍中間誠然施行公平，\textcolor{red}{不欺壓}寄居的和孤兒寡婦，在這地方\textcolor{red}{不流}無辜人的血，也\textcolor{red}{不隨從}別神陷害自己，我就使你們在這地方仍然居住，就是我古時所賜給你們列祖的地，直到永遠。
&『看哪，你們倚靠虛謊無益的話！你們偷盜、殺害、姦淫、起假誓、向巴力燒香，並隨從素不認識的別神，且來到這稱為我名下的殿，在我面前敬拜，又說：「我們可以自由了！」你們這樣的舉動是要行那些可憎的事嗎？這稱為我名下的殿，在你們眼中豈可看為賊窩嗎？我都看見了！』」這是耶和華說的。\\\hline
\end{tabular}
\end{table}

 



\subsubsection{從出埃及吩咐的一件事:聽話,違逆者之祭不蒙悅納 7:21-28}
萬軍之耶和華以色列的神如此說：「你們將燔祭加在平安祭上，吃肉吧！ 22 因為我將你們列祖從埃及地領出來的那日，燔祭平安祭的事我並沒有提說，也沒有吩咐他們。 23 我只吩咐他們這一件，說：『你們當聽從我的話，我就做你們的神，你們也做我的子民。你們行我所吩咐的一切道，就可以得福。』 24 他們卻不聽從，不側耳而聽，竟隨從自己的計謀和頑梗的惡心，向後不向前。 25 自從你們列祖出埃及地的那日直到今日，我差遣我的僕人眾先知到你們那裡去，每日從早起來差遣他們。 26 你們卻不聽從，不側耳而聽，竟硬著頸項行惡，比你們列祖更甚。」

27 「你要將這一切的話告訴他們，他們卻不聽從；呼喚他們，他們卻不答應。 28 你要對他們說：『這就是不聽從耶和華他們神的話，不受教訓的國民！從他們的口中，誠實滅絕了。』」

\subsubsection{陀斐特稱為殺戮谷,百姓屍首為飛鳥和野獸食物 7:29-34}
耶路撒冷啊，要剪髮拋棄，在淨光的高處舉哀，因為耶和華丟掉離棄了惹他憤怒的世代。 30 耶和華說：「猶大人行我眼中看為惡的事，將可憎之物設立在稱為我名下的殿中，汙穢這殿。 31 他們在欣嫩子谷建築陀斐特的丘壇，好在火中焚燒自己的兒女。這並不是我所吩咐的，也不是我心所起的意。」 32 耶和華說：「因此，日子將到，這地方不再稱為陀斐特和欣嫩子谷，反倒稱為殺戮谷。因為要在陀斐特葬埋屍首，甚至無處可葬。 33 並且這百姓的屍首必給空中的飛鳥和地上的野獸做食物，並無人鬨趕。 34 那時，我必使猶大城邑中和耶路撒冷街上歡喜和快樂的聲音、新郎和新婦的聲音，都止息了，因為地必成為荒場。」



%\subsection{百姓为何不得痊愈呢 8}
\subsubsection{死人骸骨不再收殮,恆久背道不肯回頭 8:1-7}
\begin{table}[H]
\centering
\begin{tabular}{p{6.75cm}|p{10.50cm}}\hline\hline
\multicolumn{1}{c|}{死人骸骨不再收殮1-2}&\multicolumn{1}{c}{斥其百姓恆久背道不肯回頭 8:3-7}\\\hline
耶和華說：「到那時，人必將猶大王的骸骨和他首領的骸骨、祭司的骸骨、先知的骸骨並耶路撒冷居民的骸骨，都從墳墓中取出來， 2 拋散在日頭、月亮和天上眾星之下，就是他們從前所喜愛、所侍奉、所隨從、所求問、所敬拜的。這些骸骨不再收殮，不再葬埋，必在地面上成為糞土。 
&並且這惡族所剩下的民在我所趕他們到的各處，寧可揀死不揀生。」這是萬軍之耶和華說的。你要對他們說：『耶和華如此說：人跌倒，不再起來嗎？人轉去，不再轉來嗎？ 5 這耶路撒冷的民為何恆久背道呢？他們守定詭詐，不肯回頭。 6 我留心聽，聽見他們說不正直的話，無人悔改惡行說：「我做的是什麼呢？」他們各人轉奔己路，如馬直闖戰場。 7 空中的鸛鳥知道來去的定期，斑鳩燕子與白鶴也守候當來的時令，我的百姓卻不知道耶和華的法則。\\\hline
\end{tabular}
\end{table}


\subsubsection{文士虛假,先知祭司輕忽醫治 8:8-13}
8 『你們怎麼說「我們有智慧，耶和華的律法在我們這裡」？看哪，\textcolor{red}{文士}的假筆舞弄虛假。 9 智慧人慚愧，驚惶，被擒拿，他們棄掉耶和華的話，心裡還有什麼智慧呢？ 10 所以我必將他們的妻子給別人，將他們的田地給別人為業。因為他們從最小的到至大的都一味地貪婪，從\textcolor{red}{先知}到\textcolor{red}{祭司}都行事虛謊。 11 他們輕輕忽忽地醫治我百姓的損傷，說「平安了！平安了！」，其實沒有平安。 12 他們行可憎的事知道慚愧嗎？不然！他們毫不慚愧，也不知羞恥。因此，他們必在仆倒的人中仆倒，我向他們討罪的時候，他們必致跌倒。』」

\subsubsection{上帝警百姓必受重災,先知憂損傷不得痊癒 8:14-22}
\begin{table}[H]
\centering
\begin{tabular}{p{7.75cm}|p{9.50cm}}\hline\hline
\multicolumn{1}{c|}{上帝警百姓必受重災}&\multicolumn{1}{c}{先知憂損傷不得痊癒}\\\hline
這是耶和華說的。耶和華說：「我必使他們全然滅絕。葡萄樹上必沒有葡萄，無花果樹上必沒有果子，葉子也必枯乾。我所賜給他們的，必離開他們過去。」 
&我們為何靜坐不動呢？我們當聚集，進入堅固城，在那裡靜默不言。因為耶和華我們的神使我們靜默不言，又將苦膽水給我們喝，都因我們得罪了耶和華。我們指望平安，卻得不著好處；指望痊癒的時候，不料受了驚惶！\\\hline
聽見從但那裡敵人的馬噴鼻氣,他的壯馬發嘶聲,全地就都震動。因為他們來吞滅這地和其上所有的,吞滅這城與其中的居民。「看哪,我必使毒蛇到你們中間,是不服法術的,必咬你們。」這是耶和華說的。
&我有憂愁，願能自慰，我心在我裡面發昏。 聽啊，是我百姓的哀聲從極遠之地而來，說：「耶和華不在錫安嗎？錫安的王不在其中嗎？ \\\hline
耶和華說：「他們為什麼以雕刻的偶像和外邦虛無的神惹我發怒呢？」 &麥秋已過，夏令已完，我們還未得救！」先知說：因我百姓的損傷，我也受了損傷；我哀痛，驚惶將我抓住。在基列豈沒有乳香呢？在那裡豈沒有醫生呢？我百姓為何不得痊癒呢？\\\hline
\end{tabular}
\end{table}

%\subsection{以認識耶和华夸口 9}
\subsubsection{欲到曠野哭泣猶大因說謊，\textcolor{red}{地變為荒場}無智慧人明白 9:1-12}
但願我的頭為水，我的眼為淚的泉源，我好為我百姓(民女)中被殺的人晝夜哭泣！ 2 唯願我在曠野有行路人住宿之處，使我可以離開我的民出去，因他們都是行姦淫的，是行詭詐的一黨。 3 「他們彎起舌頭像弓一樣，為要說謊話。他們在國中增長勢力，不是為行誠實，乃是惡上加惡，並不認識我。」這是耶和華說的。 4 「你們各人當謹防鄰舍，不可信靠弟兄，因為弟兄盡行欺騙，鄰舍都往來讒謗人。 5 他們各人欺哄鄰舍，不說真話；他們教舌頭學習說謊，勞勞碌碌地作孽。 6 你的住處在詭詐的人中，他們因行詭詐，不肯認識我。」這是耶和華說的。

7 所以萬軍之耶和華如此說：「看哪，我要將他們熔化熬煉，不然我因我百姓的罪該怎樣行呢？ 8 他們的舌頭是毒箭，說話詭詐，人與鄰舍口說和平話，心卻謀害他。」 9 耶和華說：「我豈不因這些事討他們的罪呢？豈不報復這樣的國民呢？」

%\subsubsection{耶路撒冷變為荒場 9:10-11}
我要為山嶺哭泣悲哀，為曠野的草場揚聲哀號，因為都已乾焦，甚至無人經過。人也聽不見牲畜鳴叫，空中的飛鳥和地上的野獸都已逃去。 11 「我必使耶路撒冷變為亂堆，為野狗的住處；也必使猶大的城邑變為荒場，無人居住。」誰是智慧人，可以明白這事？耶和華的口向誰說過，使他可以傳說？遍地為何滅亡，乾焦好像曠野，甚至無人經過呢？

\subsubsection{因棄法違命被逐\textcolor{red}{人屍首不得收取}，心未受割禮与異邦受罰 9:12-26}
 13 耶和華說：「因為這百姓離棄我在他們面前所設立的律法，沒有遵行，也沒有聽從我的話， 14 只隨從自己頑梗的心行事，照他們列祖所教訓的隨從眾巴力。」 15 所以萬軍之耶和華以色列的神如此說：「看哪，我必將茵陳給這百姓吃，又將苦膽水給他們喝。 16 我要把他們散在列邦中，就是他們和他們列祖素不認識的列邦。我也要使刀劍追殺他們，直到將他們滅盡。」

%\subsubsection{唱哀歌因為死亡上來 9:17-22}
萬軍之耶和華如此說：「你們應當思想，將善唱哀歌的婦女召來，又打發人召善哭的婦女來。 18 叫她們速速為我們舉哀，使我們眼淚汪汪，使我們的眼皮湧出水來。 19 因為聽見哀聲出於錫安說：『我們怎樣敗落了！我們大大地慚愧，我們撇下地土，人也拆毀了我們的房屋。』」 20 婦女們哪，你們當聽耶和華的話，領受他口中的言語；又當教導你們的兒女舉哀，各人教導鄰舍唱哀歌。 21 因為死亡上來，進了我們的窗戶，入了我們的宮殿，要從外邊剪除孩童，從街上剪除少年人。 22 你當說：「耶和華如此說：人的屍首必倒在田野像糞土，又像收割的人遺落的一把禾稼，無人收取。」

%\subsubsection{人毋自恃財物宜恃神 9:23-24}
耶和華如此說：「智慧人不要因他的智慧誇口，勇士不要因他的勇力誇口，財主不要因他的財物誇口。 24 誇口的卻因他有聰明，認識我是耶和華，又知道我喜悅在世上施行慈愛、公平和公義，以此誇口。」這是耶和華說的。

%\subsubsection{警告猶大心中受割禮.異邦必受罰 9:25-26}
耶和華說：「看哪，日子將到，我要刑罰一切受過割禮，心卻未受割禮的， 26 就是埃及、猶大、以東、亞捫人、摩押人和一切住在曠野剃周圍頭髮的。因為列國人都沒有受割禮，以色列人心中也沒有受割禮。」




%\subsection{認識神就当降服在祂面前 10}
\subsubsection{偽神斷難比主,唯耶和華是真神 10:1-16}
\begin{table}[H]
\centering
\begin{tabular}{p{10.8cm}|p{6.50cm}}\hline\hline
\multicolumn{1}{c|}{偽神斷難比主}&\multicolumn{1}{c}{唯耶和華是真神}\\\hline
以色列家啊要聽耶和華對你們所說的話!耶和華如此說"你們不要效法列國的行為也不要為天象驚惶,因列國為此事驚惶。眾民的風俗是虛空的,他們在樹林中用斧子砍伐一棵樹,匠人用手工造成偶像。他們用金銀裝飾它,用釘子和錘子釘穩使它不動搖。它好像棕樹是鏇成的,不能說話不能行走,必須有人抬著。你們不要怕它,它不能降禍也無力降福"
&耶和華啊，沒有能比你的。你本為大，有大能大力的名。萬國的王啊，誰不敬畏你？敬畏你本是合宜的，因為在列國的智慧人中，雖有政權的尊榮，也不能比你。\\\hline
他們盡都是畜類，是愚昧的。偶像的訓誨算什麼呢？偶像不過是木頭！有銀子打成片，是從他施帶來的，並有從烏法來的金子，都是匠人和銀匠的手工；又有藍色紫色料的衣服，都是巧匠的工作。 
&唯耶和華是真神，是活神，是永遠的王。他一發怒，大地震動；他一惱恨，列國都擔當不起。\\\hline
你們要對他們如此說：「不是那創造天地的神，必從地上從天下被除滅。」
&耶和華用能力創造大地，用智慧建立世界，用聰明鋪張穹蒼。他一發聲，空中便有多水激動，他使雲霧從地極上騰。他造電隨雨而閃，從他府庫中帶出風來。\\\hline
各人都成了畜類，毫無知識，各銀匠都因他雕刻的偶像羞愧。他所鑄的偶像本是虛假的，其中並無氣息，都是虛無的，是迷惑人的工作，到追討的時候必被除滅。&雅各的份不像這些，因他是造做萬有的主，以色列也是他產業的支派，萬軍之耶和華是他的名。\\\hline
\end{tabular}
\end{table}


\subsubsection{勸民備物被掳，嘆牧者愚昧帳篷傾毀，切禱主從寬懲治 10:17-25}
\begin{table}[H]
\centering
\begin{tabular}{p{5.cm}|p{12.750cm}}\hline\hline
\multicolumn{1}{c|}{勸民備物被掳}&\multicolumn{1}{c}{嘆牧者愚昧帳篷傾毀，切禱主從寬懲治}\\\hline
受圍困的人哪，當收拾你的財物，從國中帶出去。因為耶和華如此說：「這時候，我必將此地的居民好像用機弦甩出去，又必加害在他們身上，使他們覺悟。」
&\multirow{2}{12.8cm}{我卻說:「這真是我的痛苦,必須忍受。」我的帳篷毀壞,我的繩索折斷。我的兒女離我出去,沒有了,無人再支搭我的帳篷,掛起我的幔子。因為牧人都成為畜類,沒有求問耶和華,所以不得順利,他們的羊群也都分散。有風聲,看哪,敵人來了！有大擾亂從北方出來,要使猶大城邑變為荒涼,成為野狗的住處。耶和華啊,我曉得人的道路不由自己,行路的人也不能定自己的腳步。耶和華啊,求你從寬懲治我,不要在你的怒中懲治我,恐怕使我歸於無有。願你將憤怒傾在不認識你的列國中,和不求告你名的各族上。因為他們吞了雅各,不但吞了,而且滅絕,把他的住處變為荒場。}\\\cline{1-1}
民說:「禍哉！我受損傷,我的傷痕極其重大。」&\\\hline
\end{tabular}
\end{table}

%\subsection{猶大悔改，就要真正回到与神的关系中 11}
\subsubsection{宣告遵行主約,民頑梗不聽且隨從與城數目相等別神,咒詛必臨 11:1-13}
\begin{table}[H]
\centering
\begin{tabular}{p{5.75cm}|p{11.50cm}}\hline\hline
\multicolumn{1}{c|}{宣告遵行主約}&\multicolumn{1}{c}{民頑梗不聽且隨從與城數目相等別神,咒詛必臨}\\\hline
\multirow{2}{5.86cm}{耶和華的話臨到耶利米說：「當聽這約的話，告訴猶大人和耶路撒冷的居民，對他們說：『耶和華以色列的神如此說：不聽從這約之話的人必受咒詛。這約是我將你們列祖從埃及地領出來，脫離鐵爐的那日所吩咐他們的，說：「你們要聽從我的話，照我一切所吩咐的去行。這樣你們就做我的子民，我也做你們的神。」 5 我好堅定向你們列祖所起的誓，給他們流奶與蜜之地，正如今日一樣。』」我就回答說：「耶和華啊，阿們！」}
&耶和華對我說：「你要在猶大城邑中和耶路撒冷街市上，宣告這一切話說：你們當聽從遵行這約的話。 7 因為我將你們列祖從埃及地領出來的那日，直到今日，都是從早起來，切切告誡他們說：『你們當聽從我的話。』 8 他們卻不聽從，不側耳而聽，竟隨從自己頑梗的惡心去行。所以我使這約中一切咒詛的話臨到他們身上，這約是我吩咐他們行的，他們卻不去行。」\\\cline{2-2}
&耶和華對我說：「在猶大人和耶路撒冷居民中有同謀背叛的事。 10 他們轉去效法他們的先祖，不肯聽我的話，犯罪作孽，又隨從別神，侍奉它。以色列家和猶大家背了我與他們列祖所立的約。」 11 所以耶和華如此說：「我必使災禍臨到他們，是他們不能逃脫的。他們必向我哀求，我卻不聽。 12 那時，猶大城邑的人和耶路撒冷的居民要去哀求他們燒香所供奉的神，只是遭難的時候，這些神毫不拯救他們。 13 猶大啊，你神的數目與你城的數目相等，你為那可恥的巴力所築燒香的壇，也與耶路撒冷街道的數目相等！\\\hline
\end{tabular}
\end{table}

\subsubsection{禁耶利米為民禱，因主已定意降禍 11:14-17}
「所以你不要為這百姓祈禱，不要為他們呼求禱告，因為他們遭難向我哀求的時候，我必不應允。 15 我所親愛的既行許多淫亂，聖肉也離了你，你在我殿中做什麼呢？你作惡就喜樂。 16 從前耶和華給你起名叫青橄欖樹，又華美又結好果子，如今他用鬨嚷之聲，點火在其上，枝子也被折斷。 17 原來栽培你的萬軍之耶和華已經說要降禍攻擊你，是因以色列家和猶大家行惡，向巴力燒香，惹我發怒，是自作自受。」



\subsubsection{田野百獸攻擊使地荒涼，許拔出占據以色列的惡鄰 12:7-17}
7 「我離了我的殿宇，撇棄我的產業，將我心裡所親愛的交在她仇敵的手中。 8 我的產業向我如林中的獅子，她發聲攻擊我，因此我恨惡她。 9 我的產業向我豈如斑點的鷙鳥呢？鷙鳥豈在她四圍攻擊她呢？你們要去聚集田野的百獸，帶來吞吃吧！ 10 許多牧人毀壞我的葡萄園，踐踏我的份，使我美好的份變為荒涼的曠野。 11 他們使地荒涼，地既荒涼，便向我悲哀。全地荒涼，因無人介意。 12 滅命的都來到曠野中一切淨光的高處，耶和華的刀從地這邊直到地那邊盡行殺滅，凡有血氣的都不得平安。 13 他們種的是麥子，收的是荊棘，勞勞苦苦卻毫無益處。因耶和華的烈怒，你們必為自己的土產羞愧。」

%\subsubsection{許拔出占據以色列的惡鄰，歸正者得返故土 12:14-17}
耶和華如此說：「一切惡鄰，就是占據我使百姓以色列所承受產業的，我要將他們拔出本地，又要將猶大家從他們中間拔出來。 15 我拔出他們以後，我必轉過來憐憫他們，把他們再帶回來，各歸本業，各歸故土。 16 他們若殷勤學習我百姓的道，指著我的名起誓說『我指著永生的耶和華起誓』，正如他們從前教我百姓指著巴力起誓，他們就必建立在我百姓中間。 17 他們若是不聽，我必拔出那國，拔出而且毀滅。」這是耶和華說的。

%\subsection{以麻帶糜爛，酒滿罈，北國滅，警猶大勿驕傲當自卑免得全被擄掠 13}



%\subsection{三次呼求声永不止息——荒旱，偽先知，棄掉猶大 14}
\subsubsection{為猶大荒旱認罪祈禱，神定意降灾禁先知祈禱 14:1-12}
\begin{table}[H]
\centering
\begin{tabular}{p{6.25cm}|p{5.50cm}|p{5.2cm}}\hline\hline
\multicolumn{1}{c|}{預言猶大荒旱}&\multicolumn{1}{c|}{耶利米認罪祈禱}&\multicolumn{1}{c}{耶利米禁先知為百姓祈禱}\\\hline
耶和華論到乾旱之災的話臨到耶利米：
「猶大悲哀，城門衰敗，眾人披上黑衣坐在地上，耶路撒冷的哀聲上達。他們的貴胄打發家僮打水，他們來到水池，見沒有水，就拿著空器皿，蒙羞慚愧，抱頭而回。耕地的也蒙羞抱頭，因為無雨降在地上，地都乾裂。田野的母鹿生下小鹿，就撇棄，因為無草。野驢站在淨光的高處，喘氣好像野狗，因為無草，眼目失明。」
&耶和華啊，我們的罪孽雖然作見證告我們，還求你為你名的緣故行事。我們本是多次背道，得罪了你。以色列所盼望，在患難時做他救主的啊，你為何在這地像寄居的？又像行路的只住一宵呢？你為何像受驚的人，像不能救人的勇士呢？耶和華啊，你仍在我們中間，我們也稱為你名下的人，求你不要離開我們！
&耶和華對這百姓如此說：「這百姓喜愛妄行(漂流)，不禁止腳步，所以耶和華不悅納他們。現今要記念他們的罪孽，追討他們的罪惡。」耶和華又對我說：「不要為這百姓祈禱求好處。他們禁食的時候，我不聽他們的呼求；他們獻燔祭和素祭，我也不悅納。我卻要用刀劍、饑荒、瘟疫滅絕他們。」\\\hline
\end{tabular}
\end{table}
%\subsubsection{耶利米認罪祈禱 14:7-9}
%\subsubsection{禁先知為百姓祈禱，要降刀劍、饑荒、瘟疫 14:10-12}


\subsubsection{為“平安”的假預言求问，神责偽先知妄言，聽其言者遭報 14:13-16}
我就說：「唉，主耶和華啊！那些先知常對他們說：『你們必不看見刀劍，也不遭遇饑荒，耶和華要在這地方賜你們長久的平安。』」 
14 耶和華對我說：「那些先知託我的名說假預言，我並沒有打發他們，沒有吩咐他們，也沒有對他們說話。他們向你們預言的，乃是虛假的異象和占卜，並虛無的事以及本心的詭詐。 
15 所以耶和華如此說：論到託我名說預言的那些先知，我並沒有打發他們，他們還說『這地不能有刀劍饑荒』，其實那些先知必被刀劍、饑荒滅絕。 16 聽他們說預言的百姓，必因饑荒、刀劍拋在耶路撒冷的街道上，無人葬埋，他們連妻子帶兒女都是如此。我必將他們的惡倒在他們身上。(我必使他們罪惡的報應臨到他們身上)

\subsubsection{為猶大被棄己民遭災悲慘哀禱 14:17-22}
「你要將這話對他們說：『願我眼淚汪汪，晝夜不息，因為我百姓(民的處女)受了裂口破壞的大傷。 18 我若出往田間，就見有被刀殺的；我若進入城內，就見有因饑荒患病的。連先知帶祭司在國中往來，也是毫無知識 (不知怎樣才好)。』」

19 你全然棄掉猶大嗎？你心厭惡錫安嗎？為何擊打我們，以致無法醫治呢？我們指望平安，卻得不著好處；指望痊癒，不料受了驚惶！ 20 耶和華啊，我們承認自己的罪惡和我們列祖的罪孽，因我們得罪了你。 21 求你為你名的緣故，不厭惡我們，不辱沒你榮耀的寶座。求你追念，不要背了與我們所立的約。 22 外邦人虛無的神中有能降雨的嗎？天能自降甘霖嗎？耶和華我們的神啊，能如此的不是你嗎？所以，我們仍要等候你，因為這一切都是你所造的。


\subsubsection{耶和華棄絕猶大,命定加以諸災 15:1-9}
\begin{table}[H]
\centering
\begin{tabular}{p{14.5cm}|p{2.750cm}}\hline\hline
\multicolumn{1}{c|}{耶和華}&\multicolumn{1}{c}{百姓、先知}\\\hline
耶和華對我說：「雖有摩西和撒母耳站在我面前代求，我的心也不顧惜這百姓。你將他們從我眼前趕出，叫他們去吧！ 
&他們問你說：『我們往哪裡去呢？』\\\hline
你便告訴他們：『耶和華如此說：定為死亡的，必致死亡；定為刀殺的，必交刀殺；定為饑荒的，必遭饑荒；定為擄掠的，必被擄掠。耶和華說：我命定四樣害他們，就是刀劍殺戮，狗類撕裂，空中的飛鳥和地上的野獸吞吃毀滅。又必使他們在天下萬國中拋來拋去，都因猶大王希西家的兒子瑪拿西在耶路撒冷所行的事。』」
&「耶路撒冷啊，誰可憐你呢？誰為你悲傷呢？誰轉身問你的安呢？」\\\hline
耶和華說：「你棄絕了我，轉身退後，因此我伸手攻擊你、毀壞你，我後悔甚不耐煩。我在境內各城門口(我在這地邊界的關口)，用簸箕簸了我的百姓，使他們喪掉兒女。我毀滅他們，他們仍不轉離所行的道。他們的寡婦在我面前比海沙更多。我使滅命的午間來，攻擊少年人的母親，使痛苦驚嚇忽然臨到她身上。生過七子的婦人力衰氣絕，尚在白晝，日頭忽落，她抱愧蒙羞。其餘的人，我必在他們敵人跟前交於刀劍。」這是耶和華說的。&\\\hline
\end{tabular}
\end{table}




%\subsection{示耶利米不可娶妻}



\subsubsection{示百姓降大災因其行惡比列祖更甚，許以色列仍歸故土 16:10-15}
你將這一切的話指示這百姓，他們問你說：『耶和華為什麼說要降這大災禍攻擊我們呢？我們有什麼罪孽呢？我們向耶和華我們的神犯了什麼罪呢？』

你就對他們說：『耶和華說：因為你們列祖離棄我，隨從別神，侍奉敬拜，不遵守我的律法。 12 而且你們行惡比你們列祖更甚，因為各人隨從自己頑梗的惡心行事，甚至不聽從我。 13 所以我必將你們從這地趕出，直趕到你們和你們列祖素不認識的地，你們在那裡必晝夜侍奉別神，因為我必不向你們施恩。』
%\subsubsection{許以色列人仍歸故土 16:14-15}
耶和華說：「日子將到，人必不再指著那領以色列人從埃及地上來之永生的耶和華起誓， 15 卻要指著那領以色列人從北方之地，並趕他們到的各國上來之永生的耶和華起誓。並且我要領他們再入我從前賜給他們列祖之地。」

\subsubsection{以色列的罪孽不能再隱藏，先知以耶和華為力量保障 16:16-21}
16 耶和華說：「我要召許多打魚的把以色列人打上來，然後我要召許多打獵的，從各山上、各岡上、各石穴中獵取他們。 17 因我的眼目察看他們的一切行為，他們不能在我面前遮掩，他們的罪孽也不能在我眼前隱藏。 18 我先要加倍報應他們的罪孽和罪惡，因為他們用可憎之屍玷汙我的地土，又用可厭之物充滿我的產業。」

%\subsubsection{偶像無益，唯耶和華有能力 16:19-21}
耶和華啊，你是我的力量，是我的保障，在苦難之日是我的避難所。列國人必從地極來到你這裡，說：「我們列祖所承受的不過是虛假，是虛空無益之物。 20 人豈可為自己製造神呢？其實這不是神！」

21 耶和華說：「我要使他們知道，就是這一次使他們知道我的手和我的能力，他們就知道我的名是耶和華了。」













\subsection{比喻}
\subsubsection{引示羅被弃，警猶大勿靠殿;禁先知禱告，因神定意傾怒 7:12-20}
\begin{table}[H]
\centering
\begin{tabular}{p{8.25cm}|p{8.750cm}}\hline\hline
\multicolumn{1}{c|}{引示羅被弃，警猶大勿靠殿12-15}&\multicolumn{1}{c}{禁先知禱告，因神定意傾怒16-20}\\\hline
你們且往示羅去，就是我先前立為我名的居所，察看我因這百姓以色列的罪惡向那地所行的如何。」耶和華說：「現在因你們行了這一切的事，我也從早起來警戒你們，你們卻不聽從，呼喚你們，你們卻不答應。所以，我要向這稱為我名下、你們所倚靠的殿，與我所賜給你們和你們列祖的地施行，照我從前向示羅所行的一樣。我必將你們從我眼前趕出，正如趕出你們的眾弟兄，就是以法蓮的一切後裔。
&所以你不要為這百姓祈禱，不要為他們呼求禱告，也不要向我為他們祈求，因我不聽允你。他們在猶大城邑中和耶路撒冷街上所行的，你沒有看見嗎？孩子撿柴，父親燒火，婦女摶麵做餅，獻給天后，又向別神澆奠祭，惹我發怒。」耶和華說：「他們豈是惹我發怒呢？不是自己惹禍，以致臉上慚愧嗎？」所以主耶和華如此說：「看哪，我必將我的怒氣和憤怒傾在這地方的人和牲畜身上，並田野的樹木和地裡的出產上，必如火著起，不能熄滅。\\\hline
\end{tabular}
\end{table}

\subsubsection{以麻帶糜爛喻必照樣敗壞猶大的驕傲，使其變為無用 13:1-11}
耶和華對我如此說：「你去買一根麻布帶子束腰，不可放在水中。」 2 我就照著耶和華的話，買了一根帶子束腰。 3 耶和華的話第二次臨到我說： 4 「要拿著你所買的腰帶，就是你腰上的帶子，起來往幼發拉底河去，將腰帶藏在那裡的磐石穴中。」 5 我就去，照著耶和華所吩咐我的，將腰帶藏在幼發拉底河邊。 6 過了多日，耶和華對我說：「你起來往幼發拉底河去，將我吩咐你藏在那裡的腰帶取出來。」 7 我就往幼發拉底河去，將腰帶從我所藏的地方刨出來，見腰帶已經變壞，毫無用了。

8 耶和華的話臨到我說： 9 「耶和華如此說：我必照樣敗壞猶大的驕傲和耶路撒冷的大驕傲。 10 這惡民不肯聽我的話，按自己頑梗的心而行，隨從別神，侍奉敬拜，他們也必像這腰帶變為無用。 11 耶和華說：腰帶怎樣緊貼人腰，照樣，我也使以色列全家和猶大全家緊貼我，好叫他們屬我為子民，使我得名聲，得頌讚，得榮耀，他們卻不肯聽。

\subsubsection{以酒滿罈為喻，君王百姓都酩酊大醉不蒙憐憫被擄掠 13:12-19}
「所以你要對他們說：『耶和華以色列的神如此說：各壇都要盛滿了酒。』他們必對你說：『我們豈不確知各壇都要盛滿了酒呢？』 13 你就要對他們說：『耶和華如此說：我必使這地的一切居民，就是坐大衛寶座的君王，和祭司與先知，並耶路撒冷的一切居民，都酩酊大醉。 14 耶和華說：我要使他們彼此相碰，就是父與子彼此相碰。我必不可憐，不顧惜，不憐憫，以致滅絕他們。』」

你們當\textcolor{red}{聽}，當側耳而聽，不要驕傲，因為耶和華已經說了。 16 耶和華你們的神未使黑暗來到，你們的腳未在昏暗山上絆跌之先，當將榮耀歸給他，免得你們盼望光明，他使光明變為死蔭成為幽暗。 17 你們若不聽這話，我必因你們的驕傲在暗地哭泣，我眼必痛哭流淚，因為耶和華的群眾被擄去了。

你要對君王和太后說：『你們當自卑，坐在下邊，因你們的頭巾，就是你們的華冠，已經脫落了。』 19 南方的城盡都關閉，無人開放；猶大全被擄掠，且擄掠淨盡。

\subsubsection{以北方滅國為喻，警猶大弃惡行离偶像 13:20-27}
你們要舉目\textcolor{red}{觀看}從北方來的人！先前賜給你的群眾，就是你佳美的群眾，如今在哪裡呢？ 21 耶和華立你自己所交的朋友為首，轄制你，那時你還有什麼話說呢？痛苦豈不將你抓住像產難的婦人嗎？ 22 你若心裡說：『這一切事為何臨到我呢？』你的衣襟揭起，你的腳跟受傷，是因你的罪孽甚多。 23 古實人豈能改變皮膚呢？豹豈能改變斑點呢？若能，你們這習慣行惡的便能行善了。 24 所以我必用曠野的風吹散他們，像吹過的碎秸一樣。」 
25 耶和華說：「這是你所當得的，是我量給你的份，因為你忘記我，倚靠虛假(偶像)。 26 所以我要揭起你的衣襟蒙在你臉上，顯出你的醜陋。 27 你那些可憎惡之事，就是在田野的山上行姦淫，發嘶聲，做淫亂的事，我都看見了。耶路撒冷啊，你有禍了！你不肯潔淨，還要到幾時呢？」

\subsubsection{罪用鐵筆、金鋼鑽刻在心版上和壇角为喻 17:1-27}
%\subsubsection{棄主賴人和偶像者受禍,倚恃耶和華者得福 17:1-13}
\begin{table}[H]
\centering
\begin{tabular}{p{11.75cm}|p{5.50cm}}\hline\hline
\multicolumn{1}{c|}{棄主賴人和偶像者受禍}&\multicolumn{1}{c}{倚恃耶和華者得福}\\\hline
猶大的罪是用鐵筆、用金鋼鑽記錄的，銘刻在他們的心版上和壇角上。他們的兒女記念他們高岡上、青翠樹旁的壇和木偶。我田野的山哪，我必因你在四境之內所犯的罪，把你的貨物、財寶並丘壇當掠物交給仇敵。並且你因自己的罪必失去我所賜給你的產業，我也必使你在你所不認識的地上服侍你的仇敵，因為你使我怒中起火，直燒到永遠。」
&\multirow{2}{5.6cm}{倚靠耶和華，以耶和華為可靠的，那人有福了！他必像樹栽於水旁，在河邊扎根。炎熱來到並不懼怕，葉子仍必青翠；在乾旱之年毫無掛慮，而且結果不止。}\\\cline{1-1}
耶和華如此說：「倚靠人血肉的膀臂，心中離棄耶和華的，那人有禍了！因他必像沙漠的杜松，不見福樂來到，卻要住曠野乾旱之處，無人居住的鹼地。&\\\hline

人心比萬物都詭詐，壞到極處，誰能識透呢？我耶和華是鑒察人心、試驗人肺腑的，要照各人所行的和他做事的結果報應他。那不按正道得財的，好像鷓鴣抱不是自己下的蛋，到了中年，那財都必離開他，他終究成為愚頑人。
&我們的聖所是榮耀的寶座，從太初安置在高處。耶和華以色列的盼望啊，凡離棄你的，必致蒙羞。\\\hline
\end{tabular}
\end{table}


\subsubsection{留意當以安息日為聖 17:19-27}
\begin{table}[H]
\centering
\begin{tabular}{p{7.75cm}|p{9.50cm}}\hline\hline
\multicolumn{1}{c|}{责眾人不受教訓}&\multicolumn{1}{c}{聽命的福和不聽的禍}\\\hline
耶和華對我如此說：「你去站在平民的門口，就是猶大君王出入的門，又站在耶路撒冷的各門口，對他們說：『你們這猶大君王和猶大眾人並耶路撒冷的一切居民，凡從這些門進入的，都當聽耶和華的話！耶和華如此說：你們要謹慎，不要在安息日擔什麼擔子進入耶路撒冷的各門，也不要在安息日從家中擔出擔子去。無論何工都不可做，只要以安息日為聖日，正如我所吩咐你們列祖的。他們卻不聽從，不側耳而聽，竟硬著頸項不聽，不受教訓。
&『耶和華說：你們若留意聽從我，在安息日不擔什麼擔子進入這城的各門，只以安息日為聖日，在那日無論何工都不做，那時就有坐大衛寶座的君王和首領，他們與猶大人並耶路撒冷的居民，或坐車或騎馬進入這城的各門，而且這城必存到永遠。也必有人從猶大城邑和耶路撒冷四圍的各處，從便雅憫地、高原、山地並南地而來，都帶燔祭、平安祭、素祭和乳香並感謝祭到耶和華的殿去。你們若不聽從我，不以安息日為聖日，仍在安息日擔擔子進入耶路撒冷的各門，我必在各門中點火，這火也必燒毀耶路撒冷的宮殿，不能熄滅。』\\\hline
\end{tabular}
\end{table}


%\subsection{后悔的神在等人悔改 18}
\subsubsection{以陶匠製器喻離惡行惡的不同结局 18:1-17}
\begin{table}[H]
\centering
\begin{tabular}{p{6.cm}|p{11.0cm}}\hline\hline
\multicolumn{1}{c|}{陶匠製器}&\multicolumn{1}{c}{離惡行惡的不同结局}\\\hline
耶和華的話臨到耶利米說：「你起來，下到窯匠的家裡去，我在那裡要使你聽我的話。」我就下到窯匠的家裡去，正遇他轉輪做器皿。窯匠用泥做的器皿在他手中做壞了，他又用這泥另做別的器皿。窯匠看怎樣好，就怎樣做。
&耶和華的話就臨到我說：「耶和華說：以色列家啊，我待你們，豈不能照這窯匠弄泥嗎？以色列家啊，泥在窯匠的手中怎樣，你們在我的手中也怎樣。我何時論到一邦或一國說要拔出、拆毀、毀壞，我所說的那一邦，若是轉意離開他們的惡，我就必後悔，不將我想要施行的災禍降於他們。我何時論到一邦或一國說要建立、栽植，他們若行我眼中看為惡的事，不聽從我的話，我就必後悔，不將我所說的福氣賜給他們。\\\hline
\end{tabular}
\end{table}




%\subsubsection{警告猶大人必因罪遭仇敵攻擊 18:11-17}
現在你要對猶大人和耶路撒冷的居民說：『耶和華如此說：「我造出災禍攻擊你們，定意刑罰你們。你們各人當回頭離開所行的惡道，改正你們的行動作為。」 12 他們卻說：「這是枉然！我們要照自己的計謀去行，各人隨自己頑梗的惡心做事。」

13 『所以耶和華如此說：你們且往各國訪問，有誰聽見這樣的事？以色列民（處女）行了一件極可憎惡的事。 14 黎巴嫩的雪從田野的磐石上豈能斷絕呢？從遠處流下的涼水豈能乾涸呢？ 15 我的百姓竟忘記我，向假神燒香，使他們在所行的路上、在古道上絆跌，使他們行沒有修築的斜路， 16 以致他們的地令人驚駭，常常嗤笑，凡經過這地的，必驚駭搖頭。 17 我必在仇敵面前分散他們，好像用東風吹散一樣。遭難的日子，我必以背向他們，不以面向他們。』」

%\subsection{先知因說預言被譏辱害命，自叹生不如死 18:18-20:18}



\subsubsection{以毀瓶為喻明示猶大國必滅 19:1-15}
\begin{table}[H]
\centering
\begin{tabular}{p{6.75cm}|p{5.35cm}|p{5.15cm}}\hline\hline
\multicolumn{1}{c|}{向別神燒香,築巴力丘壇}&\multicolumn{1}{c|}{陀斐特和欣嫩子谷稱為殺戮谷}&\multicolumn{1}{c}{耶路撒冷如陀斐特}\\\hline
耶和華如此說:「你去買窯匠的瓦瓶,又帶百姓中的長老和祭司中的長老,出去到欣嫩子谷,哈珥西(瓦片)的門口那裡,宣告我所吩咐你的話說:猶大君王和耶路撒冷的居民哪,當聽耶和華的話!萬軍之耶和華以色列的神如此說:我必使災禍臨到這地方,凡聽見的人都必耳鳴。因為他們和他們列祖並猶大君王離棄我,將這地方看為平常,在這裡向素不認識的別神燒香,又使這地方滿了無辜人的血。又建築巴力的丘壇,好在火中焚燒自己的兒子,作為燔祭獻給巴力。這不是我所吩咐的,不是我所提說的,也不是我心所起的意。 
&耶和華說:因此日子將到,這地方不再稱為陀斐特和欣嫩子谷,反倒稱為殺戮谷。我必在這地方使猶大和耶路撒冷的計謀落空,也必使他們在仇敵面前倒於刀下,並尋索其命的人手下。他們的屍首,我必給空中的飛鳥和地上的野獸做食物。我必使這城令人驚駭,嗤笑,凡經過的人必因這城所遭的災驚駭,嗤笑。我必使他們在圍困窘迫之中,就是仇敵和尋索其命的人窘迫他們的時候,各人吃自己兒女的肉和朋友的肉。
&你要在同去的人眼前打碎那瓶,對他們說:『萬軍之耶和華如此說：我要照樣打碎這民和這城,正如人打碎窯匠的瓦器,以致不能再囫圇。並且人要在陀斐特葬埋屍首,甚至無處可葬。耶和華說:我必向這地方和其中的居民如此行,使這城與陀斐特一樣。耶路撒冷的房屋和猶大君王的宮殿是已經被玷汙的,就是他們在其上向天上的萬象燒香,向別神澆奠祭的宮殿房屋,都必與陀斐特一樣。』
\\\hline
\end{tabular}
\end{table}







\section{沙龍（約哈斯[王下23:30]）死在被擄之地 22:10-12}
%\subsubsection{沙龍（約哈斯）死在被擄之地 22:10-12}
不要為死人哭號,不要為他悲傷,卻要為離家出外的人大大哭號，因為他不得再回來,也不得再見他的本國。因為耶和華論到從這地方出去的猶大王約西亞的兒子沙龍(王下23:30約哈斯),就是接續他父親約西亞做王的,這樣說:他必不得再回到這裡來,卻要死在被擄去的地方,必不得再見這地。


\section{約雅敬 22:13-23/25-26/35-36/45}
\subsection{約雅敬登基時 26}
\subsubsection{以應許與警戒勸殿中眾民及民約反應 26:1-19}
\begin{table}[H]
\centering
\begin{tabular}{p{8.25cm}|p{8.50cm}}\hline\hline
\multicolumn{1}{c|}{以應許與警戒勸殿中眾民}&\multicolumn{1}{c}{祭司、先知、眾民與首領的反應}\\\hline
\multirow{2}{8.35cm}{猶大王約西亞的兒子約雅敬登基的時候,有這話從耶和華臨到耶利米說:「耶和華如此說:你站在耶和華殿的院內,對猶大眾城邑的人,就是到耶和華殿來禮拜的,說我所吩咐你的一切話,一字不可刪減。或者他們肯聽從,各人回頭離開惡道,使我後悔,不將我因他們所行的惡想要施行的災禍降於他們。你要對他們說:『耶和華如此說：你們若不聽從我,不遵行我設立在你們面前的律法,不聽我從早起來差遣到你們那裡去我僕人眾先知的話(你們還是沒有聽從),我就必使這殿如示羅,使這城為地上萬國所咒詛的。』」}
&耶利米在耶和華殿中說的這些話,祭司、先知與眾民都聽見了。耶利米說完了耶和華所吩咐他對眾人說的一切話，\textcolor{red}{祭司、先知與眾民}都來抓住他，說：「你必要死！你為何託耶和華的名預言說『這殿必如示羅，這城必變為荒場無人居住』呢？」於是，眾民都在耶和華的殿中聚集到耶利米那裡。\\\cline{2-2}
&\textcolor{blue}{猶大的首領}聽見這事，就從王宮上到耶和華的殿，坐在耶和華殿的新門口。\textcolor{red}{祭司、先知}對\textcolor{blue}{首領和眾民}說：「這人是該死的！因為他說預言攻擊這城，正如你們親耳所聽見的。」\\\hline
耶利米就對\textcolor{blue}{眾首領和眾民}說：「耶和華差遣我預言攻擊這殿和這城，說你們所聽見的這一切話。 13 現在要改正你們的行動作為，聽從耶和華你們神的話，他就必後悔，不將所說的災禍降於你們。 14 至於我，我在你們手中，你們眼看何為善何為正，就那樣待我吧！ 15 但你們要確實地知道：若把我治死，就使無辜人的血歸到你們和這城並其中的居民了，因為耶和華實在差遣我到你們這裡來，將這一切話傳於你們耳中。」
&\textcolor{blue}{首領和眾民}就對\textcolor{red}{祭司、先知}說:「這人是不該死的,因為他是奉耶和華我們神的名向我們說話。」 國中的長老就有幾個人起來,對聚會的眾民說:「當猶大王希西家的日子,有摩利沙人彌迦對猶大眾人預言說:『萬軍之耶和華如此說:錫安必被耕種像一塊田,耶路撒冷必變為亂堆,這殿的山必像叢林的高處。』 猶大王希西家和猶大眾人豈是把他治死呢?希西家豈不是敬畏耶和華,懇求他的恩嗎？耶和華就後悔,不把自己所說的災禍降於他們。若治死這人,我們就做了大惡自害己命！」\\\hline

\end{tabular}
\end{table}

\subsubsection{烏利亞說預言因懼怕逃埃及被約雅敬王殺,耶利米被人保護 26:20-24}
又有一個人奉耶和華的名說預言，是基列耶琳人示瑪雅的兒子烏利亞。他照耶利米的一切話，說預言攻擊這城和這地。 21 約雅敬王和他眾勇士、眾首領聽見了烏利亞的話，王就想要把他治死。烏利亞聽見就懼怕，逃往埃及去了。 22 約雅敬王便打發亞革波的兒子以利拿單，帶領幾個人往埃及去。 23 他們就從埃及將烏利亞帶出來，送到約雅敬王那裡。王用刀殺了他，把他的屍首拋在平民的墳地中。

然而，沙番的兒子亞希甘保護耶利米，不交在百姓的手中治死他。


\subsection{巴比倫第一次围城 35}
\subsubsection{約雅敬時神差耶利米給利甲族喝酒遭拒，因其克守先祖遺訓 35:1-11}
\begin{table}[H]
\centering
\begin{tabular}{p{7.75cm}|p{9.50cm}}\hline\hline
\multicolumn{1}{c|}{神差耶利米給利甲族喝酒1-5}&\multicolumn{1}{c}{克守先祖不喝酒遺訓}\\\hline
當猶大王約西亞之子約雅敬的時候，耶和華的話臨到耶利米說：「你去見利甲族的人，和他們說話，領他們進入耶和華殿的一間屋子，給他們酒喝。」
&\multirow{2}{9.7cm}{他們卻說：「我們不喝酒，因為我們先祖利甲的兒子約拿達曾吩咐我們說：『你們與你們的子孫永不可喝酒，也不可蓋房、撒種、栽種葡萄園，但一生的年日要住帳篷，使你們的日子在寄居之地得以延長。』凡我們先祖利甲的兒子約拿達所吩咐我們的話，我們都聽從了。我們和我們的妻子、兒女一生的年日都不喝酒，也不蓋房居住，也沒有葡萄園、田地和種子，但住帳篷，聽從我們先祖約拿達的話，照他所吩咐我們的去行。巴比倫王尼布甲尼撒上此地來，我們因怕迦勒底的軍隊和亞蘭的軍隊，就說：『來吧，我們到耶路撒冷去。』這樣，我們才住在耶路撒冷。」}\\\cline{1-1}
我就將哈巴洗尼雅的孫子、雅利米雅的兒子雅撒尼亞和他弟兄，並他眾子以及利甲全族的人，領到耶和華的殿，進入神人伊基大利的兒子哈難眾子的屋子。那屋子在首領的屋子旁邊，在沙龍之子把門的瑪西雅屋子以上。於是我在利甲族人面前設擺盛滿酒的碗和杯，對他們說：「請你們喝酒。」&\\\hline
\end{tabular}
\end{table}

\subsubsection{主責民不守神命致禍，利甲族因守先祖約拿達之命致福 35:12-19}
\begin{table}[H]
\centering
\begin{tabular}{p{9cm}|p{8.5cm}}\hline\hline
\multicolumn{1}{c|}{主責民不守神命致禍}&\multicolumn{1}{c}{利甲族永不缺人侍立在神面前}\\\hline
耶和華的話臨到耶利米說：「萬軍之耶和華以色列的神如此說：你去對猶大人和耶路撒冷的居民說：『耶和華說：你們不受教訓，不聽從我的話嗎？&利甲的兒子約拿達所吩咐他子孫不可喝酒的話，他們已經遵守，直到今日也不喝酒，因為他們聽從先祖的吩咐。\\\hline
我從早起來警戒你們,你們卻不聽從我。我從早起來,差遣我的僕人眾先知去，說：「你們各人當回頭離開惡道，改正行為，不隨從侍奉別神，就必住在我所賜給你們和你們列祖的地上。」只是你們沒有聽從我，也沒有側耳而聽。&利甲的兒子約拿達的子孫能遵守先人所吩咐他們的命，這百姓卻沒有聽從我。\\\hline
因此耶和華萬軍之神,以色列的神如此說：我要使我所說的一切災禍臨到猶大人和耶路撒冷的一切居民。因為我對他們說話，他們沒有聽從；我呼喚他們，他們沒有答應。
&耶利米對利甲族的人說:「萬軍之耶和華以色列的神如此說:因你們聽從你們先祖約拿達的吩咐,謹守他的一切誡命,照他所吩咐你們的去行,所以萬軍之耶和華以色列的神如此說:利甲的兒子約拿達必永不缺人侍立在我面前」\\\hline
\end{tabular}
\end{table}



\subsection{約雅敬第四年 25/36:1-7/45}
\subsubsection{責猶大23年之久不聽警教,被擄70年後刑罰巴比倫 25:1-14}
\begin{table}[H]
\centering
\begin{tabular}{p{8.cm}|p{9.cm}}\hline\hline
\multicolumn{1}{c|}{責猶大23年之久不聽警教1-7}&\multicolumn{1}{c}{被擄七十年後刑罰巴比倫8-14}\\\hline
\multirow{2}{8.25cm}{猶大王約西亞的兒子約雅敬第四年，就是巴比倫王尼布甲尼撒的元年，耶和華論猶大眾民的話臨到耶利米。先知耶利米就將這話對猶大眾人和耶路撒冷的一切居民說：從猶大王亞們的兒子約西亞十三年直到今日，這二十三年之內，常有耶和華的話臨到我，我也對你們傳說，就是從早起來傳說，只是你們沒有聽從。耶和華也從早起來，差遣他的僕人眾先知到你們這裡來（只是你們沒有聽從，也沒有側耳而聽），說：「你們各人當回頭離開惡道和所做的惡，便可居住耶和華古時所賜給你們和你們列祖之地，直到永遠。不可隨從別神，侍奉敬拜，以你們手所做的惹我發怒，這樣我就不加害於你們。然而，你們沒有聽從我，竟以手所做的惹我發怒，陷害自己。」這是耶和華說的。}
&所以萬軍之耶和華如此說：「因為你們沒有聽從我的話，我必召北方的眾族和我僕人巴比倫王尼布甲尼撒來攻擊這地和這地的居民，並四圍一切的國民。我要將他們盡行滅絕，以致他們令人驚駭、嗤笑，並且永久荒涼。」這是耶和華說的。「我又要使歡喜和快樂的聲音、新郎和新婦的聲音、推磨的聲音和燈的亮光從他們中間止息。這全地必然荒涼，令人驚駭，\textcolor{blue}{這些國民要服侍巴比倫王七十年}。
\\\cline{2-2}
&「\textcolor{blue}{七十年滿了以後，我必刑罰巴比倫王和那國民}，並迦勒底人之地，因他們的罪孽使那地永遠荒涼。」這是耶和華說的。「我也必使我向那地所說的話，就是記在這書上的話，是耶利米向這些國民說的預言，都臨到那地。因為有多國和大君王必使迦勒底人做奴僕，我也必照他們的行為，按他們手所做的報應他們。」\\\hline
\end{tabular}
\end{table}

\subsubsection{給列國民喝憤怒的酒喻列國之災 25:15-31}
\begin{table}[H]
\centering
\begin{tabular}{p{9.25cm}|p{8.cm}}\hline\hline
\multicolumn{1}{c|}{列國民的范围 15-26}&\multicolumn{1}{c}{酒喻列國之災的内容 27-31}\\\hline
\multirow{3}{9.35cm}{耶和華以色列的神對我如此說：「你從我手中接這杯憤怒的酒，使我所差遣你去的各國的民喝。 16 他們喝了就要東倒西歪，並要發狂，因我使刀劍臨到他們中間。」 17 我就從耶和華的手中接了這杯，給耶和華所差遣我去的各國的民喝， 18 就是耶路撒冷和猶大的城邑並耶路撒冷的君王與首領，使這城邑荒涼，令人驚駭、嗤笑、咒詛，正如今日一樣。19 又有埃及王法老和他的臣僕、首領以及他的眾民， 20 並雜族的人民和烏斯地的諸王，與非利士地的諸王，亞實基倫，加沙，以革倫，以及亞實突剩下的人， 21 以東，摩押，亞捫人， 22 推羅的諸王，西頓的諸王，海島的諸王， 23 底但，提瑪，布斯和一切剃周圍頭髮的， 24 阿拉伯的諸王，住曠野雜族人民的諸王， 25 心利的諸王，以攔的諸王，瑪代的諸王， 26 北方遠近的諸王以及天下地上的萬國喝了，以後示沙克(就是巴比倫)王也要喝。}
&27 「你要對他們說：『萬軍之耶和華以色列的神如此說：你們要喝，且要喝醉，要嘔吐，且要跌倒不得再起來，都因我使刀劍臨到你們中間。』 。\\\cline{2-2}
&28 他們若不肯從你手接這杯喝，你就要對他們說：『萬軍之耶和華如此說：你們一定要喝！ 29 我既從稱為我名下的城起首施行災禍，你們能盡免刑罰嗎？你們必不能免，因為我要命刀劍臨到地上一切的居民。這是萬軍之耶和華說的。』。\\\cline{2-2}
&所以你要向他們預言這一切的話，攻擊他們，說：『耶和華必從高天吼叫，從聖所發聲，向自己的羊群大聲吼叫。他要向地上一切的居民呐喊，像踹葡萄的一樣。 31 必有響聲達到地極，因為耶和華與列國相爭。凡有血氣的，他必審問；至於惡人，他必交給刀劍。這是耶和華說的。』\\\hline
\end{tabular}
\end{table}

\subsubsection{牧者宜因大災禍哀號哭泣 25:32-37}
萬軍之耶和華如此說：「看哪，必有災禍從這國發到那國，並有大暴風從地極颳起。」 33 到那日，從地這邊直到地那邊，都有耶和華所殺戮的。必無人哀哭，不得收殮，不得葬埋，必在地上成為糞土。 

34 牧人哪，你們當哀號、呼喊！群眾的頭目啊，你們要滾在灰中！因為你們被殺戮分散的日子足足來到，你們要跌碎，好像美器打碎一樣。 35 牧人無路逃跑，群眾的頭目也無法逃脫。 36 聽啊，有牧人呼喊，有群眾頭目哀號的聲音，因為耶和華使他們的草場變為荒場。 

37 耶和華發出猛烈的怒氣，平安的羊圈就都寂靜無聲。 38 他離了隱密處像獅子一樣，他們的地因刀劍凶猛的欺壓，又因他猛烈的怒氣，都成為可驚駭的。

\subsubsection{約雅敬第四年耶利米召巴錄書其預言於卷 36:1-7}
\begin{table}[H]
\centering
\begin{tabular}{p{7.5cm}|p{10.2cm}}\hline\hline
\multicolumn{1}{c|}{願你榮耀你的兒子，使兒子也榮耀你}&\multicolumn{1}{c}{求你使我同你享榮耀}\\\hline
猶大王約西亞的兒子約雅敬第四年,耶和華的話臨到耶利米說"你取一書卷,將我對你說攻擊以色列和猶大並各國的一切話,從我對你說話的那日,就是從約西亞的日子起,直到今日,都寫在其上.或者猶大家聽見我想要降於他們的一切災禍,各人就回頭離開惡道,我好赦免他們的罪孽和罪惡"
&所以耶利米召了尼利亞的兒子巴錄來,巴錄就從耶利米口中,將耶和華對耶利米所說的一切話寫在書卷上.耶利米吩咐巴錄說:「我被拘管,不能進耶和華的殿,所以你要去趁禁食的日子,在耶和華殿中將耶和華的話,就是你從我口中所寫在書卷上的話,念給百姓和一切從猶大城邑出來的人聽.或者他們在耶和華面前懇求,各人回頭離開惡道,因為耶和華向這百姓所說要發的怒氣和憤怒是大的.」 \\\hline
\end{tabular}
\end{table}

\subsubsection{約雅敬第四年，巴錄懷憂耶利米慰以主言 45:1-5}
\begin{table}[H]
\centering
\begin{tabular}{p{8.5cm}|p{8.50cm}}\hline\hline
\multicolumn{1}{c|}{約雅敬第四年巴錄唉哼而困乏}&\multicolumn{1}{c}{耶利米慰以主言}\\\hline
猶大王約西亞的兒子約雅敬第四年，尼利亞的兒子巴錄將先知耶利米口中所說的話寫在書上。耶利米說：巴錄啊，耶和華以色列的神說：巴錄(你)曾說：『哀哉！耶和華將憂愁加在我的痛苦上，我因唉哼而困乏，不得安歇。』
&你要這樣告訴他：『耶和華如此說：我所建立的我必拆毀，我所栽植的我必拔出，在全地我都如此行。你為自己圖謀大事嗎？不要圖謀！我必使災禍臨到凡有血氣的，但你無論往哪裡去，我必使你以自己的命為掠物。』這是耶和華說的。\\\hline
\end{tabular}
\end{table}

\subsection{約雅敬第五年九月 36:8-32}
\subsubsection{約雅敬第五年九月巴錄在殿中誦話於民前,眾首領聽了害怕要告訴王 36:8-19}
 \begin{table}[H]
\centering
\begin{tabular}{p{4.5cm}|p{6.5cm}|p{5.50cm}}\hline\hline
\multicolumn{1}{c|}{巴錄在殿中誦書卷話於民前}&\multicolumn{1}{c|}{眾首領到王宮聚集要聽書上的話}&\multicolumn{1}{c}{眾首領聽了害怕要告訴王}\\\hline
尼利亞的兒子巴錄就照先知耶利米一切所吩咐的去行,在耶和華的殿中從書上念耶和華的話。猶大王約西亞的兒子約雅敬第五年九月,耶路撒冷的眾民和那從猶大城邑來到耶路撒冷的眾民,在耶和華面前宣告禁食的日子。巴錄就在耶和華殿的上院、耶和華殿的新門口、沙番的兒子文士基瑪利雅的屋內,念書上耶利米的話給眾民聽。
&沙番的孫子、基瑪利雅的兒子米該亞聽見書上耶和華的一切話，他就下到王宮，進入文士的屋子。眾首領，就是文士以利沙瑪、示瑪雅的兒子第萊雅、亞革波的兒子以利拿單、沙番的兒子基瑪利雅、哈拿尼雅的兒子西底家和其餘的首領，都坐在那裡。於是米該亞對他們述說他所聽見的一切話，就是巴錄向百姓念那書的時候所聽見的。眾首領就打發古示的曾孫、示利米雅的孫子、尼探雅的兒子猶底到巴錄那裡，對他說：「你將所念給百姓聽的書卷拿在手中，到我們這裡來。」
&尼利亞的兒子巴錄就手拿書卷來到他們那裡。他們對他說：「請你坐下，念給我們聽。」巴錄就念給他們聽。他們聽見這一切話就害怕，面面相觀，對巴錄說：「我們必須將這一切話告訴王。」他們問巴錄說：「請你告訴我們，你怎樣從他口中寫這一切話呢？」巴錄回答說：「他用口向我說這一切話，我就用筆墨寫在書上。」眾首領對巴錄說：「你和耶利米要去藏起來，不可叫人知道你們在哪裡。」\\\hline
\end{tabular}
\end{table}

\subsubsection{約雅敬聞卷中預言割而焚之,復取卷另書前言,焚卷者必受重災 36:20-32}
\begin{table}[H]
\centering
\begin{tabular}{p{8.5cm}|p{9.0cm}}\hline\hline
\multicolumn{1}{c|}{王聞卷中預言割而焚之,差人捉拿巴錄和耶利米}&\multicolumn{1}{c}{復取卷另書前言,焚卷者必受重災}\\\hline
眾首領進院見王，卻先把書卷存在文士以利沙瑪的屋內，以後將這一切話說給王聽。王就打發猶底去拿這書卷來。他便從文士以利沙瑪的屋內取來，念給王和王左右侍立的眾首領聽。那時正是九月，王坐在過冬的房屋裡，王的前面火盆中有燒著的火。猶底念了三四篇(行)，王就用文士的刀將書卷割破，扔在火盆中，直到全卷在火中燒盡了。王和聽見這一切話的臣僕都不懼怕，也不撕裂衣服。以利拿單和第萊雅並基瑪利雅懇求王不要燒這書卷，他卻不聽。
&\multirow{2}{9.4cm}{王燒了書卷（其上有巴錄從耶利米口中所寫的話）以後，耶和華的話臨到耶利米說：「你再取一卷，將猶大王約雅敬所燒第一卷上的一切話寫在其上。論到猶大王約雅敬你要說：『耶和華如此說：你燒了書卷，說：「你為什麼在其上寫著說，巴比倫王必要來毀滅這地，使這地上絕了人民牲畜呢？」所以耶和華論到猶大王約雅敬說：他後裔中必沒有人坐在大衛的寶座上，他的屍首必被拋棄，白日受炎熱，黑夜受寒霜。我必因他和他後裔並他臣僕的罪孽刑罰他們，我要使我所說的一切災禍臨到他們和耶路撒冷的居民並猶大人，只是他們不聽。』」於是，耶利米又取一書卷交給尼利亞的兒子文士巴錄，他就從耶利米的口中寫了猶大王約雅敬所燒前卷上的一切話，另外又添了許多相仿的話。}\\\cline{1-1}
王就吩咐哈米勒的兒子(王的兒子)耶拉篾和亞斯列的兒子西萊雅，並亞伯疊的兒子示利米雅去捉拿文士巴錄和先知耶利米，耶和華卻將他們隱藏。
&\\\hline
\end{tabular}
\end{table}
 
\subsection{約雅敬结局 22:13-23}
\subsubsection{約雅敬蓋房為爭勝貪婪不義行強暴,像驢被扔在耶城門外 22:13-23}
『那行不義蓋房，行不公造樓，白白使用人的手工不給工價的，有禍了！ 14 他說：「我要為自己蓋廣大的房、寬敞的樓，為自己開窗戶，這樓房的護牆板是香柏木的，樓房是丹色油漆的。」 15 難道你做王是在乎造香柏木樓房爭勝嗎？你的父親豈不是也吃也喝，也施行公平和公義嗎？那時他得了福樂。 16 他為困苦和窮乏人申冤，那時就得了福樂，認識我不在乎此嗎？這是耶和華說的。

17 唯有你的眼和你的心專顧貪婪，流無辜人的血，行欺壓和強暴。 18 所以耶和華論到猶大王約西亞的兒子約雅敬如此說：人必不為他舉哀說「哀哉，我的哥哥」，或說「哀哉，我的姐姐」；也不為他舉哀說「哀哉，我的主」，或說「哀哉，我主的榮華」。 19 他被埋葬，好像埋驢一樣，要拉出去扔在耶路撒冷的城門之外。

20 『你要上黎巴嫩哀號，在巴珊揚聲，從亞巴琳哀號，因為你所親愛的都毀滅了。 21 你興盛的時候，我對你說話，你卻說：「我不聽！」你自幼年以來總是這樣，不聽從我的話。 22 你的牧人要被風吞吃，你所親愛的必被擄去，那時你必因你一切的惡抱愧蒙羞。 23 你這住黎巴嫩，在香柏樹上搭窩的，有痛苦臨到你，好像疼痛臨到產難的婦人，那時你何等可憐！



\section{哥尼雅/約雅斤/耶哥尼雅 22:24-30/52:31-34}
\subsection{哥尼雅被擄不得歸回,後裔無人坐大衛寶座 22:24-30}

24 『耶和華說：猶大王約雅敬的兒子哥尼雅(耶哥尼雅)，雖是我右手上帶印的戒指，我憑我的永生起誓，也必將你從其上摘下來。 25 並且我必將你交給尋索你命的人和你所懼怕的人手中，就是巴比倫王尼布甲尼撒和迦勒底人的手中。 26 我也必將你和生你的母親趕到別國，並不是你們生的地方，你們必死在那裡。 27 但心中甚想歸回之地，必不得歸回。 28 哥尼雅這人是被輕看、破壞的器皿嗎？是無人喜愛的器皿嗎？他和他的後裔為何被趕到不認識之地呢？ 29 地啊，地啊，地啊，當聽耶和華的話！ 30 耶和華如此說：要寫明這人算為無子，是平生不得亨通的，因為他後裔中再無一人得亨通，能坐在大衛的寶座上治理猶大。』」

\subsection{被擄後三十七，巴比倫王恩待約雅斤 52:31-34}
猶大王約雅斤被擄後三十七年，巴比倫王以未米羅達元年十二月二十五日，使猶大王約雅斤抬頭，提他出監， 32 又對他說恩言，使他的位高過與他一同在巴比倫眾王的位， 33 給他脫了囚服。他終身在巴比倫王面前吃飯。 34 巴比倫王賜他所需用的食物，日日賜他一份，終身是這樣，直到他死的日子。


\section{西底家 21/22:1-9/23-24/27-34/37-38/39:1-10/52:1-30}
\subsection{耶哥尼雅被擄、西底家登基 24/27}

\subsubsection{殿前的佳劣無花果--佳者得返故土,劣者喻西底家及其臣僕遭禍 24:1-10}
\begin{table}[H]
\centering
\begin{tabular}{p{5.25cm}|p{11.5cm}}\hline\hline
\multicolumn{1}{c|}{殿前的佳劣無花果 1-3}&\multicolumn{1}{c}{佳者得返故土,劣者喻西底家及其臣僕遭禍4-10}\\\hline
\multirow{2}{5.5cm}{巴比倫王尼布甲尼撒將猶大王約雅敬的兒子耶哥尼雅和猶大的首領,並工匠、鐵匠從耶路撒冷擄去,帶到巴比倫。這事以後,耶和華指給我看,有兩筐無花果放在耶和華的殿前:一筐是極好的無花果，好像是初熟的；一筐是極壞的無花果,壞得不可吃。於是耶和華問我說:「耶利米你看見什麼?」我說:「我看見無花果,好的極好,壞的極壞,壞得不可吃。」}
&耶和華的話臨到我說：「耶和華以色列的神如此說：被擄去的猶大人，就是我打發離開這地到迦勒底人之地去的，我必看顧他們如這\textcolor{red}{好無花果}，使他們得好處。我要眷顧他們，使他們得好處，領他們歸回這地。我也要建立他們，必不拆毀；栽植他們，並不拔出。我要\textcolor{blue}{賜他們認識我的心}，知道我是耶和華。他們要做我的子民，我要做他們的神，因為他們要一心歸向我。\\\cline{2-2}
&「耶和華如此說：我必將猶大王西底家和他的首領，以及剩在這地耶路撒冷的餘民，並住在埃及地的猶大人都交出來，好像那\textcolor{red}{極壞，壞得不可吃的無花果}。我必使他們交出來，在天下萬國中拋來拋去,遭遇災禍,在我趕逐他們到的各處成為凌辱、笑談、譏刺、咒詛。我必使\textcolor{blue}{刀劍、饑荒、瘟疫}臨到他們,直到他們從我所賜給他們和他們列祖之地滅絕。」\\\hline
\end{tabular}
\end{table}


\subsubsection{西底家登基，耶利米以索與軛喻列國必侍巴比倫王 27:1-11}
\begin{table}[H]
\centering
\begin{tabular}{p{9cm}|p{7.50cm}}\hline\hline
\multicolumn{1}{c|}{耶利米以索與軛喻列國必侍巴比倫王1-7}&\multicolumn{1}{c}{不侍諸災/服侍之福8-11}\\\hline
\multirow{2}{9.2cm}{猶大王約西亞的兒子約雅敬(西底家27:3)登基的時候，有這話從耶和華臨到耶利米說：耶和華對我如此說：「你做繩索與軛，加在自己的頸項上，藉那些來到耶路撒冷見猶大王西底家的使臣之手，把繩索與軛送到以東王、摩押王、亞捫王、推羅王、西頓王那裡。且囑咐使臣傳於他們的主人說：『萬軍之耶和華以色列的神如此說：我用大能和伸出來的膀臂創造大地和地上的人民、牲畜，我看給誰相宜，就把地給誰。現在我將這些地都交給我僕人巴比倫王尼布甲尼撒的手，我也將田野的走獸給他使用。列國都必服侍他和他的兒孫，直到他本國遭報的日期來到，那時多國和大君王要使他做他們的奴僕。}
&無論哪一邦哪一國，\textcolor{red}{不肯}服侍這巴比倫王尼布甲尼撒，也不把頸項放在巴比倫王的軛下，我必用刀劍、饑荒、瘟疫刑罰那邦，直到我藉巴比倫王的手將他們毀滅。這是耶和華說的。至於你們，不可聽從你們的先知和占卜的、圓夢的、觀兆的以及行邪術的，他們告訴你們說你們不致服侍巴比倫王。他們向你們說假預言，要叫你們遷移，遠離本地，以致我將你們趕出去，使你們滅亡。\\\cline{2-2}
&但哪一邦\textcolor{red}{肯}把頸項放在巴比倫王的軛下服侍他，我必使那邦仍在本地存留，得以耕種居住。這是耶和華說的。\\\hline
\end{tabular}
\end{table}

\subsubsection{勸西底家服侍巴比倫王,對祭司眾民預言聖殿器物必存於巴比倫 27:12-22}
\begin{table}[H]
\centering
\begin{tabular}{p{6.cm}|p{11.50cm}}\hline\hline
\multicolumn{1}{c|}{勸西底家服侍巴比倫王12-15}&\multicolumn{1}{c}{對祭司眾民預言聖殿器物必存於巴比倫16-22}\\\hline
\multirow{2}{6.2cm}{我就照這一切的話對猶大王西底家說:「要把你們的頸項放在巴比倫王的軛下,服侍他和他的百姓,便得存活。你和你的百姓為何要因刀劍、饑荒、瘟疫死亡，正如耶和華論到不服侍巴比倫王的那國說的話呢？不可聽那些先知對你們所說的話，他們說你們不致服侍巴比倫王，其實他們向你們說假預言。耶和華說：我並沒有打發他們，他們卻託我的名說假預言，好使我將你們和向你們說預言的那些先知趕出去，一同滅亡。」}
&我又對祭司和這眾民說:耶和華如此說:你們不可聽那先知對你們所說的預言,他們說『耶和華殿中的器皿快要從巴比倫帶回來』,其實他們向你們說假預言。不可聽從他們,只管服侍巴比倫王便得存活。這城何至變為荒場呢?他們若果是先知,有耶和華的話臨到他們,讓他們祈求萬軍之耶和華,使那在耶和華殿中和猶大王宮內,並耶路撒冷剩下的器皿不被帶到巴比倫去。\\\cline{2-2} 
&因為萬軍之耶和華論到柱子、銅海、盆座並剩在這城裡的器皿，就是巴比倫王尼布甲尼撒擄掠猶大王約雅敬的兒子耶哥尼雅和猶大、耶路撒冷一切貴胄的時候所沒有掠去的器皿—— 論到那在耶和華殿中和猶大王宮內並耶路撒冷剩下的器皿，萬軍之耶和華以色列的神如此說：必被帶到巴比倫存在那裡，直到我眷顧以色列人的日子。那時，我必將這器皿帶回來，交還此地。這是耶和華說的。\\\hline
\end{tabular}
\end{table}

\subsubsection{西底家登基不聽從耶和華藉先知所說的話 37:1-2}
約西亞的兒子西底家代替約雅敬的兒子哥尼雅為王，是巴比倫王尼布甲尼撒立在猶大地做王的。 2 但西底家和他的臣僕並國中的百姓，都不聽從耶和華藉先知耶利米所說的話。

\subsubsection{西底家登基神默示以攔必受懲罰後日必復振興 49:34-39}
猶大王西底家登基的時候，耶和華論以攔的話臨到先知耶利米說： 35 「萬軍之耶和華如此說：我必折斷以攔人的弓，就是他們為首的權力。 36 我要使四風從天的四方颳來臨到以攔人，將他們分散四方 e，這被趕散的人沒有一國不到的。 37 耶和華說：我必使以攔人在仇敵和尋索其命的人面前驚惶。我也必使災禍，就是我的烈怒臨到他們。又必使刀劍追殺他們，直到將他們滅盡。 38 我要在以攔設立我的寶座，從那裡除滅君王和首領。這是耶和華說的。39 「到末後，我還要使被擄的以攔人歸回。這是耶和華說的。」


\subsubsection{西底家登基行耶和華眼中看為惡的事 52:1-3}
西底家登基的時候年二十一歲，在耶路撒冷做王十一年。他母親名叫哈慕她，是立拿人耶利米的女兒。 2 西底家行耶和華眼中看為惡的事，是照約雅敬一切所行的。 3 因此耶和華的怒氣在耶路撒冷和猶大發作，以致將人民從自己的面前趕出。

\subsection{巴比倫王尼布甲尼撒來攻擊西底家 21/22:1-9？/23？}
\subsubsection{西底家因被攻遣使詢主,先知言城必交在巴比倫王手中 21:1-7}
\begin{table}[H]
\centering
\begin{tabular}{p{5.25cm}|p{11.250cm}}\hline\hline
\multicolumn{1}{c|}{西底家因被攻遣使詢主}&\multicolumn{1}{c}{先知言城必交在巴比倫王手中}\\\hline
耶和華的話臨到耶利米。那時西底家王打發瑪基雅的兒子巴施戶珥和瑪西雅的兒子祭司西番雅，去見耶利米，說： 2 「請你為我們求問耶和華，因為巴比倫王尼布甲尼撒來攻擊我們，或者耶和華照他一切奇妙的作為待我們，使巴比倫王離開我們上去。」
&耶利米對他們說：「你們當對西底家這樣說： 4 『耶和華以色列的神如此說：我要使你們手中的兵器，就是你們在城外與巴比倫王和圍困你們的迦勒底人打仗的兵器，翻轉過來，又要使這些都聚集在這城中。 5 並且我要在怒氣、憤怒和大惱恨中，用伸出來的手並大能的膀臂親自攻擊你們。 6 又要擊打這城的居民，連人帶牲畜都必遭遇大瘟疫死亡。 7 以後我要將猶大王西底家和他的臣僕百姓，就是在城內從瘟疫、刀劍、饑荒中剩下的人，都交在巴比倫王尼布甲尼撒的手中和他們仇敵並尋索其命的人手中。巴比倫王必用刀擊殺他們，不顧惜，不可憐，不憐憫。這是耶和華說的。』」\\\hline
\end{tabular}
\end{table}

\subsubsection{勸民降迦勒底、王施行公平，責山上居民骄傲 21:8-14}
\begin{table}[H]
\centering
\begin{tabular}{p{6.5cm}|p{4.25cm}|p{5.25cm}}\hline\hline
\multicolumn{1}{c|}{勸民出降迦勒底人}&\multicolumn{1}{c|}{勸王施行公平}&\multicolumn{1}{c}{責山上居民骄傲}\\\hline
你要對這百姓說:『耶和華如此說:看哪,我將生命的路和死亡的路擺在你們面前。住在這城裡的,必遭刀劍、饑荒、瘟疫而死;但出去歸降圍困你們迦勒底人的,必得存活,要以自己的命為掠物。耶和華說:我向這城變臉,降禍不降福。這城必交在巴比倫王的手中,他必用火焚燒。』
&至於猶大王的家，你們當聽耶和華的話。大衛家啊，耶和華如此說：你們每早晨要施行公平，拯救被搶奪的脫離欺壓人的手，恐怕我的憤怒因你們的惡行發作,如火著起,甚至無人能以熄滅。
&耶和華說:住山谷和平原磐石上的居民,你們說:『誰能下來攻擊我們?誰能進入我們的住處呢?』看哪,我與你們為敵! 耶和華又說:我必按你們做事的結果刑罰你們,我也必使火在耶路撒冷的林中著起,將她四圍所有的盡行燒滅。」\\\hline
\end{tabular}
\end{table}


\subsubsection{勸君王施行公平公義，若不聽城必變為荒場 22:1-9？}
\begin{table}[H]
\centering
\begin{tabular}{p{8.cm}|p{8cm}}\hline\hline
\multicolumn{1}{c|}{勸君王施行公平公義1-4}&\multicolumn{1}{c|}{不聽城必變為荒場5-9}\\\hline
耶和華如此說：「你下到猶大王的宮中,在那裡說這話,說：『坐大衛寶座的猶大王啊,你和你的臣僕並進入城門的百姓,都當聽耶和華的話。 耶和華如此說：你們要施行公平和公義,拯救被搶奪的脫離欺壓人的手,不可虧負寄居的和孤兒寡婦,不可以強暴待他們,在這地方也不可流無辜人的血。你們若認真行這事,就必有坐大衛寶座的君王和他的臣僕百姓,或坐車或騎馬從這城的各門進入。
&你們若不聽這些話,耶和華說:我指著自己起誓,這城必變為荒場!耶和華論到猶大王的家如此說:我看你如基列,如黎巴嫩頂,然而我必使你變為曠野,為無人居住的城邑。我要預備行毀滅的人,各拿器械攻擊你。他們要砍下你佳美的香柏樹,扔在火中。許多國的民要經過這城,各人對鄰舍說:「耶和華為何向這大城如此行呢?」他們必回答說:「是因離棄了耶和華他們神的約,侍奉敬拜別神。」\\\hline
\end{tabular}
\end{table}


\subsubsection{假牧趕散之羊被真牧看顾,大衛苗裔掌王權 23:1-8？}
\begin{table}[H]
\centering
\begin{tabular}{p{8.5cm}|p{8.5cm}}\hline\hline
\multicolumn{1}{c|}{假牧趕散之羊被真牧看顾}&\multicolumn{1}{c}{大衛苗裔掌王權}\\\hline
耶和華說：「那些殘害、趕散我草場之羊的牧人有禍了！」 2 耶和華以色列的神斥責那些牧養他百姓的牧人，如此說：「你們趕散我的羊群，並沒有看顧他們，我必討你們這行惡的罪。」這是耶和華說的。
&\multirow{2}{8.76cm}{耶和華說：「日子將到，我要給大衛興起一個公義的苗裔。他必掌王權，行事有智慧，在地上施行公平和公義。 6 在他的日子，猶大必得救，以色列也安然居住，他的名必稱為『耶和華我們的義』。」 7 耶和華說：「日子將到，人必不再指著那領以色列人從埃及地上來永生的耶和華起誓， 8 卻要指著那領以色列家的後裔從北方和趕他們到的各國中上來永生的耶和華起誓。他們必住在本地。」}\\\cline{1-1}
3 「我要將我羊群中所餘剩的，從我趕他們到的各國內招聚出來，領他們歸回本圈，他們也必生養眾多。 4 我必設立照管他們的牧人牧養他們，他們不再懼怕，不再驚惶，也不缺少一個。」這是耶和華說的。&\\\hline
\end{tabular}
\end{table}

\subsubsection{先知帶祭司行恶褻瀆神，災禍必臨到他們 23:9-15？}
論到那些先知，我心在我裡面憂傷，我骨頭都發顫。因耶和華和他的聖言，我像醉酒的人，像被酒所勝的人。 10 地滿了行淫的人。因妄自賭咒，地就悲哀，曠野的草場都枯乾了。他們所行的道乃是惡的，他們的勇力使得不正。 11 「連先知帶祭司都是褻瀆的，就是在我殿中我也看見他們的惡。」這是耶和華說的。 12 「因此，他們的道路必像黑暗中的滑地，他們必被追趕在這路中仆倒。因為當追討之年，我必使災禍臨到他們。」這是耶和華說的。

13 「我在\textcolor{red}{撒馬利亞}的先知中曾見愚妄，他們藉巴力說預言，使我的百姓以色列走錯了路。 
我在\textcolor{red}{耶路撒冷}的先知中曾見可憎惡的事，他們行姦淫，做事虛妄，又堅固惡人的手，甚至無人回頭離開他的惡。他們在我面前都像所多瑪，耶路撒冷的居民都像蛾摩拉。」 15 所以萬軍之耶和華論到先知如此說：「我必將茵陳給他們吃，又將苦膽水給他們喝，因為褻瀆的事出於耶路撒冷的先知，流行遍地。」


\subsubsection{偽先知用假預言幻夢使民走錯路 23:16-32？}
16 萬軍之耶和華如此說：「這些先知向你們說預言，你們不要聽他們的話。他們以虛空教訓你們，所說的異象是出於自己的心，不是出於耶和華的口。 17 他們常對藐視我的人說：『耶和華說，你們必享平安。』又對一切按自己頑梗之心而行的人說：『必沒有災禍臨到你們。』」 18 有誰站在耶和華的會中得以聽見並會悟他的話呢？有誰留心聽他的話呢？ 19 看哪，耶和華的憤怒好像暴風，已經發出，是暴烈的旋風，必轉到惡人的頭上。 20 耶和華的怒氣必不轉消，直到他心中所擬定的成就了，末後的日子你們要全然明白。 21 「我沒有打發那些先知，他們竟自奔跑；我沒有對他們說話，他們竟自預言。 22 他們若是站在我的會中，就必使我的百姓聽我的話，又使他們回頭離開惡道和他們所行的惡。」
%\subsubsection{主論假先知的各樣惡言 23:23-32}
23 耶和華說：「我豈為近處的神呢？不也為遠處的神嗎？」 
24 耶和華說：「人豈能在隱密處藏身，使我看不見他呢？」
耶和華說：「我豈不充滿天地嗎？ 25 我已聽見那些先知所說的，就是託我名說的假預言，他們說：『我做了夢！我做了夢！』 26 說假預言的先知，就是預言本心詭詐的先知，他們這樣存心要到幾時呢？ 27 他們各人將所做的夢對鄰舍述說，想要使我的百姓忘記我的名，正如他們列祖因巴力忘記我的名一樣。 28 得夢的先知可以述說那夢，得我話的人可以誠實講說我的話！糠秕怎能與麥子比較呢？」這是耶和華說的。 
29 耶和華說：「我的話豈不像火，又像能打碎磐石的大錘嗎？」 
30 耶和華說：「那些先知各從鄰舍偷竊我的言語，因此我必與他們反對。」 
31 耶和華說：「那些先知用舌頭說『是耶和華說的』，我必與他們反對。」 
32 耶和華說：「那些以幻夢為預言，又述說這夢，以謊言和矜誇使我百姓走錯了路的，我必與他們反對。我沒有打發他們，也沒有吩咐他們，他們與這百姓毫無益處。」這是耶和華說的。

\subsubsection{譏真先知者必受凌辱和長久的羞恥 23:33-40？}
「無論是百姓，是先知，是祭司，問你說：『耶和華有什麼默示呢？』你就對他們說：『什麼默示啊？耶和華說：我要撇棄你們！』 34 無論是先知，是祭司，是百姓，說『耶和華的默示』，我必刑罰那人和他的家。 35 你們各人要對鄰舍，各人要對弟兄如此說：『耶和華回答什麼？耶和華說了什麼呢？』 36 『耶和華的默示』你們不可再提，各人所說的話必做自己的重擔(默示)，因為你們謬用永生神萬軍之耶和華我們神的言語。 37 你們要對先知如此說：『耶和華回答你什麼？耶和華說了什麼呢？』 38 你們若說『耶和華的默示』，耶和華就如此說：因你們說『耶和華的默示』這句話，我也打發人到你們那裡去，告訴你們不可說『耶和華的默示』， 39 所以我必全然忘記你們，將你們和我所賜給你們並你們列祖的城撇棄了。 40 又必使永遠的凌辱和長久的羞恥臨到你們，是不能忘記的。」


\subsection{西底家4年5月 28、29-31？}
\subsubsection{西底家4年5月,耶利米言木軛換鐵軛，哈拿尼雅7月亡 28:1-17}
\begin{table}[H]
\centering
\begin{tabular}{p{8.cm}|p{9.8cm}}\hline\hline
\multicolumn{1}{c|}{哈拿尼雅妄言聖器擄民必速返1-4}&\multicolumn{1}{c}{耶利米謂願如其言厥後必知孰言應驗5-9}\\\hline
當年就是猶大王西底家登基第四年五月,基遍人押朔的兒子先知哈拿尼雅,在耶和華的殿中當著祭司和眾民對我說:「萬軍之耶和華以色列的神如此說:我已經折斷巴比倫王的軛。二年之內,我要將巴比倫王尼布甲尼撒從這地掠到巴比倫的器皿,就是耶和華殿中的一切器皿,都帶回此地。我又要將猶大王約雅敬的兒子耶哥尼雅和被擄到巴比倫去的一切猶大人，帶回此地,因為我要折斷巴比倫王的軛。這是耶和華說的。」
&先知耶利米當著祭司和站在耶和華殿裡的眾民，對先知哈拿尼雅說：「阿們！願耶和華如此行，願耶和華成就你所預言的話，將耶和華殿中的器皿和一切被擄去的人從巴比倫帶回此地。然而我向你和眾民耳中所要說的話，你應當聽。從古以來，在你我以前的先知向多國和大邦說預言，論到爭戰、災禍、瘟疫的事。先知預言的平安，到話語成就的時候，人便知道他真是耶和華所差來的。」\\\hline

\multirow{2}{8.cm}{於是，先知哈拿尼雅將先知耶利米頸項上的軛取下來，折斷了。哈拿尼雅又當著眾民說：「耶和華如此說：二年之內，我必照樣從列國人的頸項上折斷巴比倫王尼布甲尼撒的軛。」}
&於是先知耶利米就走了。先知哈拿尼雅把先知耶利米頸項上的軛折斷以後,耶和華的話臨到耶利米說:你去告訴哈拿尼雅說:『耶和華如此說:你折斷木軛,卻換了鐵軛。因為萬軍之耶和華以色列的神如此說:我已將鐵軛加在這些國的頸項上,使他們服侍巴比倫王尼布甲尼撒。他們總要服侍他,我也把田野的走獸給了他。』\\\cline{2-2}
&於是先知耶利米對先知哈拿尼雅說"哈拿尼雅啊,你應當聽!耶和華並沒有差遣你,你竟使這百姓倚靠謊言。所以耶和華如此說:看哪,我要叫你去世,你今年必死,因為你向耶和華說了叛逆的話"\\\hline
這樣，先知哈拿尼雅當年七月間就死了。&\\\hline
\end{tabular}
\end{table}

\subsubsection{投巴比倫遭災之書投於幼發拉底河以示巴比倫必沉淪 51:59-64}
猶大王西底家在位第四年上巴比倫去的時候，瑪西雅的孫子、尼利亞的兒子西萊雅與王同去，西萊雅是王宮的大臣，先知耶利米有話吩咐他。 60 耶利米將一切要臨到巴比倫的災禍，就是論到巴比倫的一切話，寫在書上。 61 耶利米對西萊雅說：「你到了巴比倫，務要念這書上的話。 62 又說：『耶和華啊，你曾論到這地方說要剪除，甚至連人帶牲畜沒有在這裡居住的，必永遠荒涼。』63 「你念完了這書，就把一塊石頭拴在書上，扔在幼發拉底河中， 64 說：『巴比倫因耶和華所要降於她的災禍必如此沉下去，不再興起，人民也必困乏。』」 耶利米的話到此為止。

\subsubsection{致書被擄之民勸其安居毋躁，勿信偽先知夢兆，越七十年必返故土 29:1-14？}
\begin{table}[H]
\centering
\begin{tabular}{p{3.7cm}|p{13.50cm}}\hline\hline
\multicolumn{1}{c|}{致書被擄之民1-3}&\multicolumn{1}{c}{勸其安居毋躁，勿信偽先知夢兆，越七十年必返故土4-14}\\\hline
\multirow{3}{4.cm}{先知耶利米從耶路撒冷寄信於被擄的祭司先知和眾民,並生存的長老,就是尼布甲尼撒從耶路撒冷擄到巴比倫去的(這在耶哥尼雅王和太后太監並猶大,耶路撒冷的首領,以及工匠鐵匠都離了耶路撒冷以後)他藉沙番的兒子以利亞薩和希勒家的兒子基瑪利的手寄去,他們二人是猶大王西底家打發往巴比倫去見尼布甲尼撒王的。} 
&信上說：「萬軍之耶和華以色列的神對一切被擄去的，就是我使他們從耶路撒冷被擄到巴比倫的人，如此說： 5 你們要蓋造房屋住在其中，栽種田園吃其中所產的。 6 娶妻生兒女，為你們的兒子娶妻，使你們的女兒嫁人，生兒養女，在那裡生養眾多，不致減少。 7 我所使你們被擄到的那城，你們要為那城求平安，為那城禱告耶和華，因為那城得平安，你們也隨著得平安。\\\cline{2-2}

&萬軍之耶和華以色列的神如此說:不要被你們中間的先知和占卜的誘惑,也不要聽信自己所做的夢。因為他們託我的名對你們說假預言,我並沒有差遣他們。這是耶和華說的。\\\cline{2-2}
&耶和華如此說：為巴比倫所定的七十年滿了以後，我要眷顧你們，向你們成就我的恩言，使你們仍回此地。 11 耶和華說：我知道我向你們所懷的意念是賜平安的意念，不是降災禍的意念，要叫你們末後有指望。 12 你們要呼求我，禱告我，我就應允你們。 13 你們尋求我，若專心尋求我，就必尋見。 14 耶和華說：我必被你們尋見，我也必使你們被擄的人歸回，將你們從各國中和我所趕你們到的各處招聚了來，又將你們帶回我使你們被擄掠離開的地方。這是耶和華說的。\\\hline
\end{tabular}
\end{table}


\subsubsection{未擄的弟兄像極壞的無花果，必遭刀劍饑荒瘟疫 29:15-23？}
\begin{table}[H]
\centering
\begin{tabular}{p{9.75cm}|p{7.50cm}}\hline\hline
\multicolumn{1}{c|}{未擄的弟兄像極壞的無花果，必遭刀劍饑荒瘟疫15-23}&\multicolumn{1}{c}{假先知\textcolor{red}{亞哈、西底家}行淫必遭惡報21-23}\\\hline
你們說：『耶和華在巴比倫為我們興起先知。』所以耶和華論到坐大衛寶座的王和住在這城裡的一切百姓，就是未曾與你們一同被擄的弟兄，萬軍之耶和華如此說：看哪，我必使刀劍、饑荒、瘟疫臨到他們，使他們像極壞的無花果，壞得不可吃。我必用刀劍、饑荒、瘟疫追趕他們，使他們在天下萬國拋來拋去，在我所趕他們到的各國中，令人咒詛、驚駭、嗤笑、羞辱。耶和華說：這是因為他們沒有聽從我的話，就是我從早起來差遣我僕人眾先知去說的，無奈他們不聽。這是耶和華說的。所以你們一切被擄去的，就是我從耶路撒冷打發到巴比倫去的，當聽耶和華的話。
&萬軍之耶和華以色列的神論到哥賴雅的兒子亞哈並瑪西雅的兒子西底家，如此說：他們是託我名向你們說假預言的。我必將他們交在巴比倫王尼布甲尼撒的手中，他要在你們眼前殺害他們。住巴比倫一切被擄的猶大人必藉這二人賭咒說：『願耶和華使你像巴比倫王在火中燒的西底家和亞哈一樣！』這二人是在以色列中行了醜事，與鄰舍的妻行淫，又假託我名說我未曾吩咐他們的話。知道的是我，做見證的也是我。這是耶和華說的。\\\hline
\end{tabular}
\end{table}

\subsubsection{耶利米寄信預言被擄的示瑪雅和後裔必受惡報 29:24-32?}
\begin{table}[H]
\centering
\begin{tabular}{p{10.cm}|p{7.0cm}}\hline\hline
\multicolumn{1}{c|}{示瑪雅致書給耶路撒冷，命責耶利米24-32}&\multicolumn{1}{c}{耶利米寄信預言\textcolor{red}{示瑪雅}和\textcolor{red}{後裔}必受惡報30-32}\\\hline
論到尼希蘭人示瑪雅，你當說：『萬軍之耶和華以色列的神如此說：你曾用自己的名寄信給耶路撒冷的眾民和祭司瑪西雅的兒子西番雅並眾祭司，說： 「耶和華已經立你西番雅為祭司，代替祭司耶何耶大，使耶和華殿中有官長，好將一切狂妄自稱為先知的人用枷枷住、用鎖鎖住。現在亞拿突人耶利米向你們自稱為先知，你們為何沒有責備他呢？因為他寄信給我們在巴比倫的人說：『被擄的事必長久，你們要蓋造房屋住在其中，栽種田園吃其中所產的。』」祭司西番雅就把這信念給先知耶利米聽。
&『於是耶和華的話臨到耶利米說：你當寄信給一切被擄的人說：「耶和華論到尼希蘭人示瑪雅說：因為示瑪雅向你們說預言，我並沒有差遣他，他使你們倚靠謊言，所以耶和華如此說：我必刑罰尼希蘭人示瑪雅和他的後裔，他必無一人存留住在這民中，也不得見我所要賜予我百姓的福樂，因為他向耶和華說了叛逆的話。」這是耶和華說的。』\\\hline
\end{tabular}
\end{table}



\subsubsection{被擄之民必返故土,雅各遭難的時候必被救出來 30:1-10?}
耶和華的話臨到耶利米說： 2 「耶和華以色列的神如此說：你將我對你說過的一切話都寫在書上。 3 耶和華說：日子將到，我要使我的百姓以色列和猶大被擄的人歸回。我也要使他們回到我所賜給他們列祖之地，他們就得這地為業。這是耶和華說的。」

%\subsubsection{以色列和猶大遭禍之後，到那日必獲拯救 30:4-8}
以下是耶和華論到以色列和猶大所說的話。 5 耶和華如此說：「我們聽見聲音，是戰抖懼怕而不平安的聲音。 6 你們且訪問看看，男人有產難嗎？我怎麼看見人人用手掐腰，像產難的婦人，臉面都變青了呢？ 7 哀哉！那日為大，無日可比。這是雅各遭難的時候，但他必被救出來。」 8 萬軍之耶和華說：「到那日，我必從你頸項上折斷仇敵的軛，扭開他的繩索，外邦人不得再使你做他們的奴僕。 9 你們卻要侍奉耶和華你們的神，和我為你們所要興起的王大衛。」

%\subsubsection{慰藉雅各必使之痊癒 30:10-17}
故此，耶和華說：「我的僕人雅各啊，不要懼怕！以色列啊，不要驚惶！因我要從遠方拯救你，從被擄到之地拯救你的後裔，雅各必回來得享平靖安逸，無人使他害怕。 


\subsubsection{因雅各的罪大惡多被懲治必使之痊癒，虐主民者必受重罰 30:11-24?}
11 因我與你同在，要拯救你。也要將所趕散你到的那些國滅絕淨盡，卻不將你滅絕淨盡。倒要從寬懲治你，萬不能不罰你(以你為無罪)。」這是耶和華說的。耶和華如此說：「你的損傷無法醫治，你的傷痕極其重大。 13 無人為你分訴，使你的傷痕得以纏裹，你沒有醫治的良藥。 14 你所親愛的都忘記你，不來探問 (理會)你。我因你的罪孽甚大、罪惡眾多，曾用仇敵加的傷害傷害你，用殘忍者的懲治懲治你。 15 你為何因損傷哀號呢？你的痛苦無法醫治。我因你的罪孽甚大，罪惡眾多，曾將這些加在你身上。 16 故此，凡吞吃你的必被吞吃，你的敵人個個都被擄去，擄掠你的必成為擄物，搶奪你的必成為掠物。」 17 耶和華說：「我必使你痊癒，醫好你的傷痕，都因人稱你為被趕散的，說『這是錫安，無人來探問 (理會)的』。」

耶和華如此說：「我必使雅各被擄去的帳篷歸回，也必顧惜他的住處，城必建造在原舊的山岡，宮殿也照舊有人居住。 19 必有感謝和歡樂的聲音從其中發出。我要使他們增多，不致減少；使他們尊榮，不致卑微。

%\subsubsection{虐主民者必受重罰 30:20-24}
「他們的兒女要如往日，他們的會眾堅立在我面前，凡欺壓他們的，我必刑罰他。 21 他們的君王必是屬乎他們的，掌權的必從他們中間而出。我要使他就近我，他也要親近我，不然誰有膽量親近我呢？」這是耶和華說的。 22 「你們要做我的子民，我要做你們的神。」

23 看哪，耶和華的憤怒好像暴風，已經發出，是掃滅的暴風，必轉到惡人的頭上。 24 耶和華的烈怒必不轉消，直到他心中所擬定的成就了，末後的日子你們要明白。

\subsubsection{以色列剩下的人必歡樂，從地極被領回招 31:1-9?}
耶和華說：「那時，我必做以色列各家的神，他們必做我的子民。」 2 耶和華如此說：「脫離刀劍的，就是以色列人，我使他享安息的時候，他曾在曠野蒙恩。 3 古時(從遠方)耶和華向以色列 (我)顯現，說：『我以永遠的愛愛你，因此我以慈愛吸引你。』 4 以色列的民(處女)哪，我要再建立你，你就被建立。你必再以擊鼓為美，與歡樂的人一同跳舞而出。 5 又必在撒馬利亞的山上栽種葡萄園，栽種的人要享用所結的果子。 6 日子必到，以法蓮山上守望的人必呼叫說：『起來吧！我們可以上錫安，到耶和華我們的神那裡去。』」

7 耶和華如此說：「你們當為雅各歡樂歌唱，因萬國中為首的歡呼，當傳揚頌讚說：『耶和華啊，求你拯救你的百姓以色列所剩下的人！』 8 我必將他們從北方領來，從地極招聚，同著他們來的有瞎子、瘸子、孕婦、產婦，他們必成為大幫回到這裡來。 9 他們要哭泣而來，我要照他們懇求的引導他們，使他們在河水旁走正直的路，在其上不致絆跌，因為我是以色列的父，以法蓮是我的長子。

\subsubsection{宣揚雅各之得贖，悲哀變為歡喜 31:10-14?}
「列國啊，要聽耶和華的話，傳揚在遠處的海島說：『趕散以色列的必招聚他，又看守他，好像牧人看守羊群。』 11 因耶和華救贖了雅各，救贖他脫離比他更強之人的手。 12 他們要來到錫安的高處歌唱，又流歸耶和華施恩之地，就是有五穀、新酒和油並羊羔、牛犢之地。他們的心必像澆灌的園子，他們也不再有一點愁煩。 13 那時處女必歡樂跳舞，年少的、年老的也必一同歡樂，因為我要使他們的悲哀變為歡喜，並要安慰他們，使他們的愁煩轉為快樂。 14 我必以肥油使祭司的心滿足，我的百姓也要因我的恩惠知足。」這是耶和華說的。

\subsubsection{拉結在拉瑪哭子得主慰藉,主創新事以女衛男 31:15-22?}
耶和華如此說：「在拉瑪聽見號啕痛哭的聲音，是拉結哭她兒女，不肯受安慰，因為他們都不在了。」 16 耶和華如此說：「你禁止聲音不要哀哭，禁止眼目不要流淚，因你所做之工必有賞賜，他們必從敵國歸回。」這是耶和華說的。 17 耶和華說：「你末後必有指望，你的兒女必回到自己的境界。 18 我聽見以法蓮為自己悲嘆說：『你責罰我，我便受責罰，像不慣負軛的牛犢一樣。求你使我回轉，我便回轉，因為你是耶和華我的神。 19 我回轉以後就真正懊悔，受教以後就拍腿嘆息，我因擔當幼年的凌辱就抱愧蒙羞。』」 20 耶和華說：「以法蓮是我的愛子嗎？是可喜悅的孩子嗎？我每逢責備他，仍深顧念他，所以我的心腸戀慕他，我必要憐憫他。」

%\subsubsection{主創新事以女衛男 31:21-22}
以色列民(處女)哪，你當為自己設立指路碑，豎起引路柱，你要留心向大路，就是你所去的原路。你當回轉，回轉到你這些城邑。 22 背道的民(女子)哪，你翻來覆去要到幾時呢？耶和華在地上造了一件新事，就是女子護衛男子。」

\subsubsection{歸回之民被重新建立、栽植,與民立新約，將律法放在他們心上 31:23-34?}
萬軍之耶和華以色列的神如此說：「我使被擄之人歸回的時候，他們在猶大地和其中的城邑必再這樣說：『公義的居所啊，聖山哪，願耶和華賜福給你！』 24 猶大和屬猶大城邑的人，農夫和放羊的人，要一同住在其中。 25 疲乏的人，我使他飽飫；愁煩的人，我使他知足。」 26 先知說：我醒了，覺著睡得香甜。

27 耶和華說：「\textcolor{red}{日子將到}，我要把人的種和牲畜的種播種在以色列家和猶大家。 28 我先前怎樣留意將他們拔出、拆毀、毀壞、傾覆、苦害，也必照樣留意將他們建立、栽植。」這是耶和華說的。 29 「\textcolor{red}{當那些日子}，人不再說：『父親吃了酸葡萄，兒子的牙酸倒了。』 30 但各人必因自己的罪死亡，凡吃酸葡萄的，自己的牙必酸倒。」

%\subsubsection{與民立新約，將律法放在他們心上 31:31-34}
耶和華說：「\textcolor{red}{日子將到},我要與以色列家和猶大家另立新約。不像我拉著他們祖宗的手領他們出埃及地的時候，與他們所立的約,我雖做他們的丈夫,他們卻背了我的約。」這是耶和華說的。耶和華說：「那些\textcolor{red}{日子以後},我與以色列家所立的約乃是這樣：我要將我的律法放在他們裡面,寫在他們心上；我要做他們的神,他們要做我的子民。他們各人不再教導自己的鄰舍和自己的弟兄說『你該認識耶和華』,因為他們從最小的到至大的都必認識我。我要赦免他們的罪孽,不再記念他們的罪惡。」這是耶和華說的。

\subsubsection{永不棄以色列後裔,城必為耶和華建造 31:35-40?}
\begin{wrapfigure}{r}{0.5\linewidth}
\vspace{-2cm}
\hspace{.2cm}
\includegraphics[scale=0.6, ]{ZongJie/DXianZhiFigs/YeLiMi/JiaLiShan}
\end{wrapfigure}
35 那使太陽白日發光，使星月有定例黑夜發亮，又攪動大海，使海中波浪砰訇的，萬軍之耶和華是他的名，他如此說： 36 「這些定例若能在我面前廢掉，以色列的後裔也就在我面前斷絕，永遠不再成國。」這是耶和華說的。 37 耶和華如此說：「若能量度上天，尋察下地的根基，我就因以色列後裔一切所行的棄絕他們。」這是耶和華說的。

耶和華說：「\textcolor{red}{日子將到}，這城必為耶和華建造，從哈楠業樓直到角門。 39 準繩要往外量出，直到迦立山，又轉到歌亞。 40 拋屍的全谷和倒灰之處，並一切田地，直到汲淪溪，又直到東方馬門的拐角，都要歸耶和華為聖。不再拔出，不再傾覆，直到永遠。」


\subsection{西底家被巴比伦围困，埃及出兵退巴比伦 37-38}
\subsubsection{西底家不聽耶和華所說，差人求耶利米祈禱,預言迦勒底軍必再來攻陷斯邑 37:1-10}
約西亞的兒子西底家代替約雅敬的兒子哥尼雅為王，是巴比倫王尼布甲尼撒立在猶大地做王的。 2 但西底家和他的臣僕並國中的百姓，都不聽從耶和華藉先知耶利米所說的話。
3 西底家王打發示利米雅的兒子猶甲和祭司瑪西雅的兒子西番雅，去見先知耶利米，說：「求你為我們禱告耶和華我們的神。」 4 那時耶利米在民中出入，因為他們還沒有把他囚在監裡。

法老的軍隊已經從埃及出來，那圍困耶路撒冷的迦勒底人聽見他們的風聲，就拔營離開耶路撒冷去了。

耶和華的話臨到先知耶利米說： 7 「耶和華以色列的神如此說：猶大王打發你們來求問我，你們要如此對他說：『那出來幫助你們法老的軍隊必回埃及本國去。 8 迦勒底人必再來攻打這城，並要攻取，用火焚燒。 9 耶和華如此說：你們不要自欺說「迦勒底人必定離開我們」，因為他們必不離開。 10 你們即便殺敗了與你們爭戰的迦勒底全軍，但剩下受傷的人也必各人從帳篷裡起來，用火焚燒這城。』」

\subsubsection{伊利雅疑耶利米降敵囚之于獄,西底家私召之問預言并寬待之 37:11-21}
\begin{table}[H]
\centering
\begin{tabular}{p{7.cm}|p{10.0cm}}\hline\hline
\multicolumn{1}{c|}{伊利雅疑耶利米降敵囚之于獄11-15}&\multicolumn{1}{c}{西底家私召之問預言并寬待之16-21}\\\hline
\multirow{2}{7.25cm}{迦勒底的軍隊因怕法老的軍隊拔營離開耶路撒冷的時候，耶利米就雜在民中出離耶路撒冷，要往便雅憫地去，在那裡得自己的地業。他到了便雅憫門那裡，有守門官名叫伊利雅，是哈拿尼亞的孫子、示利米雅的兒子，他就拿住先知耶利米，說：「你是投降迦勒底人哪！」耶利米說：「你這是謊話，我並不是投降迦勒底人。」伊利雅不聽他的話，就拿住他解到首領那裡。首領惱怒耶利米，就打了他，將他囚在文士約拿單的房屋中，因為他們以這房屋當做監牢。}&耶利米來到獄中，進入牢房，在那裡囚了多日。西底家王打發人提出他來，在自己的宮內私下問他說：「從耶和華有什麼話臨到沒有?」耶利米說：「有!」又說：「你必交在巴比倫王手中。」\\\cline{2-2}
&耶利米又對西底家王說：「我在什麼事上得罪你，或你的臣僕，或這百姓，你竟將我囚在監裡呢？對你們預言巴比倫王必不來攻擊你們和這地的先知，現今在哪裡呢？主我的王啊，求你現在垂聽，准我在你面前的懇求，不要使我回到文士約拿單的房屋中，免得我死在那裡。」於是西底家王下令，他們就把耶利米交在護衛兵的院中，每天從餅鋪街取一個餅給他，直到城中的餅用盡了。這樣，耶利米仍在護衛兵的院中。\\\hline
\end{tabular}
\end{table}
 
\subsubsection{首領復陷耶利米於淤泥阱,以伯米勒求王恩救之 38:1-13}
\begin{table}[H]
\centering
\begin{tabular}{p{9cm}|p{8.25cm}}\hline\hline
\multicolumn{1}{c|}{首領復陷耶利米於淤泥阱1-6}&\multicolumn{1}{c}{以伯米勒求王恩救之 7-13}\\\hline
瑪坦的兒子示法提雅,巴施戶珥的兒子基大利,示利米雅的兒子猶甲,瑪基雅的兒子巴示戶珥,聽見耶利米對眾人所說的話說:「耶和華如此說:住在這城裡的必遭刀劍,饑荒,瘟疫而死,但出去歸降迦勒底人的必得存活,就是以自己命為掠物的,必得存活。耶和華如此說:這城必要交在巴比倫王軍隊的手中,他必攻取這城。」於是首領對王說:「求你將這人治死,因他向城裡剩下的兵丁和眾民說這樣的話,使他們的手發軟。這人不是求這百姓得平安,乃是叫他們受災禍。」西底家王說:「他在你們手中,無論何事,王也不能與你們反對。」他們就拿住耶利米,下在哈米勒的兒子(王的兒子)瑪基雅的牢獄裡。那牢獄在護衛兵的院中,他們用繩子將耶利米繫下去。牢獄裡沒有水,只有淤泥,耶利米就陷在淤泥中。
&在王宮的太監古實人以伯米勒,聽見他們將耶利米下了牢獄(那時王坐在便雅憫門口),以伯米勒就從王宮裡出來,對王說:「主我的王啊,這些人向先知耶利米一味地行惡,將他下在牢獄中。他在那裡必因飢餓而死,因為城中再沒有糧食.」王就吩咐古實人以伯米勒說:「你從這裡帶領三十人,趁著先知耶利米未死以前,將他從牢獄中提上來.」於是以伯米勒帶領這些人同去,進入王宮,到庫房以下,從那裡取了些碎布和破爛的衣服,用繩子縋下牢獄去到耶利米那裡。古實人以伯米勒對耶利米說:「你用這些碎布和破爛的衣服放在繩子上,墊你的胳肢窩.」耶利米就照樣行了。這樣,他們用繩子將耶利米從牢獄裡拉上來。耶利米仍在護衛兵的院中。\\\hline
\end{tabular}
\end{table}


\subsubsection{西底家暗詢耶利米,先知勸其投降，王命勿以斯言告眾首領 38:14-28}
\begin{table}[H]
\centering
\begin{tabular}{p{3.75cm}|p{8.250cm}|p{4.750cm}}\hline\hline
\multicolumn{1}{c|}{西底家暗詢先知14-16}&\multicolumn{1}{c|}{耶利米勸其投降迦勒底人17-23}&\multicolumn{1}{c}{王命勿以斯言告眾首領24-28}\\\hline
西底家王打發人帶領先知耶利米，進耶和華殿中第三門裡見王。王就對耶利米說：「我要問你一件事，你絲毫不可向我隱瞞。」耶利米對西底家說：「我若告訴你，你豈不定要殺我嗎？我若勸誡你，你必不聽從我。」西底家王就私下向耶利米說：「我指著那造我們生命之永生的耶和華起誓：我必不殺你，也不將你交在尋索你命的人手中。」
&耶利米對西底家說：「耶和華萬軍之神，以色列的神如此說：你若出去歸降巴比倫王的首領，你的命就必存活，這城也不致被火焚燒，你和你的全家都必存活。你若不出去歸降巴比倫王的首領，這城必交在迦勒底人手中。他們必用火焚燒，你也不得脫離他們的手。」西底家王對耶利米說：「我怕那些投降迦勒底人的猶大人，恐怕迦勒底人將我交在他們手中，他們戲弄我。」耶利米說：「迦勒底人必不將你交出。求你聽從我對你所說耶和華的話，這樣你必得好處，你的命也必存活。你若不肯出去，耶和華指示我的話乃是這樣：猶大王宮裡所剩的婦女必都帶到巴比倫王的首領那裡，這些婦女必說：『你知己的朋友催逼你，勝過你，見你的腳陷入淤泥中，就轉身退後了。』人必將你的后妃和你的兒女帶到迦勒底人那裡，你也不得脫離他們的手，必被巴比倫王的手捉住，你也必使這城被火焚燒。」
&西底家對耶利米說：「不要使人知道這些話，你就不至於死。首領若聽見了我與你說話，就來見你，問你說：『你對王說什麼話,不要向我們隱瞞,我們就不殺你。王向你說什麼話,也要告訴我們。』你就對他們說：『我在王面前懇求，不要叫我回到約拿單的房屋，死在那裡。』」隨後眾首領來見耶利米問他，他就照王所吩咐的一切話回答他們。他們不再與他說話，因為事情沒有洩漏。於是耶利米仍在護衛兵的院中，直到耶路撒冷被攻取的日子。\\\hline
\end{tabular}
\end{table}




\subsection{西底家九——十年尼布甲尼撒率領全軍來圍困耶路撒冷 39:1-14}
猶大王西底家第九年十月，巴比倫王尼布甲尼撒率領全軍來圍困耶路撒冷。 2 西底家十一年四月初九日，城被攻破。 3 耶路撒冷被攻取的時候，巴比倫王的首領尼甲沙利薛、三甲尼波、撒西金、拉撒力、尼甲沙利薛、拉墨並巴比倫王其餘的一切首領，都來坐在中門。


4 猶大王西底家和一切兵丁看見他們，就在夜間從靠近王園兩城中間的門出城逃跑，往亞拉巴逃去。 5 迦勒底的軍隊追趕他們，在耶利哥的平原追上西底家，將他拿住，帶到哈馬地的利比拉巴比倫王尼布甲尼撒那裡，尼布甲尼撒就審判他。 6 巴比倫王在利比拉，西底家眼前殺了他的眾子，又殺了猶大的一切貴胄。 7 並且剜西底家的眼睛，用銅鏈鎖著他，要帶到巴比倫去。



8 迦勒底人用火焚燒王宮和百姓的房屋，又拆毀耶路撒冷的城牆。 9 那時，護衛長尼布撒拉旦將城裡所剩下的百姓和投降他的逃民，以及其餘的民都擄到巴比倫去了。 10 護衛長尼布撒拉旦卻將民中毫無所有的窮人留在猶大地，當時給他們葡萄園和田地。

巴比倫王命善遇耶利米

11 巴比倫王尼布甲尼撒提到耶利米，囑咐護衛長尼布撒拉旦說： 12 「你領他去，好好地看待他，切不可害他。他對你怎麼說，你就向他怎麼行。」 13 護衛長尼布撒拉旦和尼布沙斯班、拉撒力、尼甲沙利薛、拉墨並巴比倫王的一切官長， 14 打發人去，將耶利米從護衛兵院中提出來，交於沙番的孫子、亞希甘的兒子基大利帶回家去。於是耶利米住在民中。





\subsection{西底家第十年，耶利米买地預這地必再有人居住 32-33}
\subsubsection{西底家第十年，王因耶利米預言囚之 32:1-5}
猶大王西底家第十年，就是尼布甲尼撒十八年，耶和華的話臨到耶利米。 2 那時巴比倫王的軍隊圍困耶路撒冷，先知耶利米囚在護衛兵的院內，在猶大王的宮中。 3 因為猶大王西底家已將他囚禁，說：「你為什麼預言說『耶和華如此說：我必將這城交在巴比倫王的手中，他必攻取這城。 4 猶大王西底家必不能逃脫迦勒底人的手，定要交在巴比倫王的手中，要口對口彼此說話，眼對眼彼此相看。 5 巴比倫王必將西底家帶到巴比倫，西底家必住在那裡，直到我眷顧他的時候。你們雖與迦勒底人爭戰，卻不順利。這是耶和華說的』？」

\subsubsection{耶利米購哈拿篾之田,命巴錄將田契存放在瓦器裡 32:6-15}
\begin{table}[H]
\centering
\begin{tabular}{p{6.75cm}|p{10.5cm}}\hline\hline
\multicolumn{1}{c|}{願你榮耀你的兒子，使兒子也榮耀你}&\multicolumn{1}{c}{求你使我同你享榮耀}\\\hline
耶利米說:耶和華的話臨到我說: 「你叔叔沙龍的兒子哈拿篾必來見你,說:『我在亞拿突的那塊地，求你買來,因你買這地是合乎贖回之理。』」我叔叔的兒子哈拿篾果然照耶和華的話來到護衛兵的院內,對我說:「我在便雅憫境內亞拿突的那塊地,求你買來,因你買來是合乎承受之理,是你當贖的。你為自己買來吧。」我耶利米就知道這是耶和華的話。 
&我便向我叔叔的兒子哈拿篾買了亞拿突的那塊地,平了十七舍客勒銀子給他。我在契上畫押,將契封緘,又請見證人來,並用天平將銀子平給他.我便將照例按規所立的買契,就是封緘的那一張和敞著的那一張,當著我叔叔的兒子哈拿篾和畫押做見證的人,並坐在護衛兵院內的一切猶大人眼前,交給瑪西雅的孫子,尼利亞的兒子巴錄。當著他們眾人眼前,我囑咐巴錄說:「萬軍之耶和華以色列的神如此說:要將這封緘的和敞著的兩張契放在瓦器裡,可以存留多日.因為萬軍之耶和華以色列的神如此說:將來在這地必有人再買房屋,田地和葡萄園.」\\\hline
\end{tabular}
\end{table}



\subsubsection{耶利米買地後禱告耶和華 32:16-25}
我將買契交給尼利亞的兒子巴錄以後，便禱告耶和華說： 17 「主耶和華啊，你曾用大能和伸出來的膀臂創造天地，在你沒有難成的事。 18 你施慈愛於千萬人，又將父親的罪孽報應在他後世子孫的懷中，是至大全能的神，萬軍之耶和華是你的名。 19 謀事有大略，行事有大能，注目觀看世人一切的舉動，為要照各人所行的和他做事的結果報應他。 

在埃及地顯神蹟奇事，直到今日在以色列和別人中間也是如此，使自己得了名聲，正如今日一樣。 21 用神蹟奇事和大能的手，並伸出來的膀臂與大可畏的事，領你的百姓以色列出了埃及， 22 將這地賜給他們，就是你向他們列祖起誓應許賜給他們流奶與蜜之地。

他們進入這地得了為業，卻不聽從你的話，也不遵行你的律法。你一切所吩咐他們行的，他們一無所行。因此，你使這一切的災禍臨到他們。 24 看哪，敵人已經來到，築壘要攻取這城，城也因刀劍、饑荒、瘟疫交在攻城的迦勒底人手中。你所說的話都成就了，你也看見了。 25 主耶和華啊，你對我說『要用銀子為自己買那塊地，又請見證人』，其實這城已交在迦勒底人的手中了。」

\subsubsection{神责民擄房燒城毀,因民行惡不受教訓,後必導民復歸故土 32:26-44}
\begin{table}[H]
\centering
\begin{tabular}{p{7.5cm}|p{6.750cm}|p{3.0cm}}\hline\hline
\multicolumn{1}{c|}{在房上向巴力別神燒香}&\multicolumn{1}{c|}{在殿中設立可憎之物}&\multicolumn{1}{c}{欣嫩子谷建壇}\\\hline
耶和華的話臨到耶利米說:「我是耶和華,是凡有血氣者的神,豈有我難成的事嗎?耶和華如此說：我必將這城交付迦勒底人的手和巴比倫王尼布甲尼撒的手,他必攻取這城。攻城的迦勒底人必來放火焚燒這城和其中的房屋,\textcolor{red}{在這房屋上,人曾向巴力燒香,向別神澆奠，惹我發怒}。以色列人和猶大人自從幼年以來，專行我眼中看為惡的事，以色列人盡以手所做的惹我發怒。這是耶和華說的。
&這城自從建造的那日直到今日,常惹我的怒氣和憤怒,使我將這城從我面前除掉.是因以色列人和猶大人一切的邪惡,就是他們和他們的君王,首領,祭司,先知,並猶大的眾人以及耶路撒冷的居民所行的,惹我發怒。他們以背向我,不以面向我。我雖從早起來教訓他們,他們卻不聽從,不受教訓,竟\textcolor{red}{把可憎之物設立在稱為我名下的殿中,汙穢了這殿}
&他們在欣嫩子谷建築巴力的丘壇，好使自己的兒女經火歸摩洛，他們行這可憎的事，使猶大陷在罪裡。這並不是我所吩咐的，也不是我心所起的意。\\\hline
\end{tabular}
\end{table}

%\subsubsection{後必導民復歸故土 32:36-44}
現在論到這城，就是你們所說已經因刀劍、饑荒、瘟疫交在巴比倫王手中的，耶和華以色列的神如此說： 37 我在怒氣、憤怒和大惱恨中將以色列人趕到各國，日後我必從那裡將他們招聚出來，領他們回到此地，使他們安然居住。 38 他們要做我的子民，我要做他們的神。 39 我要使他們彼此同心同道，好叫他們永遠敬畏我，使他們和他們後世的子孫得福樂。 40 又要與他們立永遠的約，必隨著他們施恩，並不離開他們，且使他們有敬畏我的心，不離開我。 41 我必歡喜施恩於他們，要盡心、盡意，誠誠實實將他們栽於此地。 42 因為耶和華如此說：我怎樣使這一切大禍臨到這百姓，我也要照樣使我所應許他們的一切福樂都臨到他們。 43 你們說『這地是荒涼無人民、無牲畜，是交付迦勒底人手之地』，日後在這境內必有人置買田地。 44 在便雅憫地、耶路撒冷四圍的各處、猶大的城邑、山地的城邑、高原的城邑並南地的城邑，人必用銀子買田地，在契上畫押，將契封緘，請出見證人，因為我必使被擄的人歸回。這是耶和華說的。」

\subsubsection{被擄者得返故土,歸回者必喜樂奉感謝祭,所有的城邑百廢俱興 33:1-13?}

耶利米還囚在護衛兵的院內，耶和華的話第二次臨到他說： 2 「成就的是耶和華，造做為要建立的也是耶和華，耶和華是他的名，他如此說： 3 你求告我，我就應允你，並將你所不知道、又大又難的事指示你。 


4 論到這城中的房屋和猶大王的宮室，就是拆毀為擋敵人高壘和刀劍的，耶和華以色列的神如此說： 5 人要與迦勒底人爭戰，正是拿死屍充滿這房屋，就是我在怒氣和憤怒中所殺的人。因他們的一切惡，我就掩面不顧這城。

 6 看哪！我要使這城得以痊癒安舒，使城中的人得醫治，又將豐盛的平安和誠實顯明於他們。 7 我也要使\textcolor{red}{猶大}被擄的和\textcolor{red}{以色列}被擄的歸回，並建立他們和起初一樣。 8 我要除淨他們的一切罪，就是向我所犯的罪；又要赦免他們的一切罪，就是干犯我、違背我的罪。 9 這城要在地上萬國人面前使我得頌讚、得榮耀，名為可喜可樂之城。萬國人因聽見我向這城所賜的福樂、所施的恩惠平安，就懼怕戰兢。

「耶和華如此說：你們論這地方說是荒廢無人民、無牲畜之地，但在這荒涼無人民、無牲畜的猶大城邑和耶路撒冷的街上， 11 必再聽見有歡喜和快樂的聲音、新郎和新婦的聲音，並聽見有人說：『要稱謝萬軍之耶和華，因耶和華本為善，他的慈愛永遠長存！』又有奉感謝祭到耶和華殿中之人的聲音，因為我必使這地被擄的人歸回，和起初一樣。這是耶和華說的。


「萬軍之耶和華如此說：在這荒廢無人民、無牲畜之地並其中所有的城邑，必再有牧人的住處，他們要使羊群躺臥在那裡。 13 在山地的城邑、高原的城邑、南地的城邑、便雅憫地、耶路撒冷四圍的各處和猶大的城邑，必再有羊群從數點的人手下經過。這是耶和華說的。

\subsubsection{大衛裔中必興起義者,居寶座者獻祭之事不斷,大衛之約永存 33:14-26?}
耶和華說：日子將到，我應許\textcolor{blue}{以色列}家和\textcolor{blue}{猶大}家的恩言必然成就。 15 當那日子，那時候，我必使大衛公義的\textcolor{red}{苗裔}長起來，他必在地上施行公平和公義。 16 在那日子，猶大必得救，耶路撒冷必安然居住，他的名必稱為『\textcolor{red}{耶和華我們的義}』。
「因為耶和華如此說：大衛必永不斷人坐在以色列家的寶座上， 18 祭司利未人在我面前也不斷人獻燔祭、燒素祭，時常辦理獻祭的事。」

%\subsubsection{无人能廢棄神所立的約 33:19-26}
19 耶和華的話臨到耶利米說： 20 「耶和華如此說：你們若能廢棄我所立白日黑夜的約，使白日黑夜不按時輪轉， 21 就能廢棄我與我僕人大衛所立的約，使他沒有兒子在他的寶座上為王，並能廢棄我與侍奉我的祭司利未人所立的約。 22 天上的萬象不能數算，海邊的塵沙也不能斗量，我必照樣使我僕人大衛的後裔和侍奉我的利未人多起來。」

23 耶和華的話臨到耶利米說： 24 「你沒有揣摩這百姓的話嗎？他們說：『耶和華所揀選的二族，他已經棄絕了。』他們這樣藐視我的百姓，以為不再成國。 25 耶和華如此說：若是我立白日黑夜的約不能存住，若是我未曾安排天地的定例， 26 我就棄絕雅各的後裔和我僕人大衛的後裔，不使大衛的後裔治理亞伯拉罕、以撒、雅各的後裔。因為我必使他們被擄的人歸回，也必憐憫他們。」


\subsection{亡國之前西底家被围 34}
\subsubsection{西底家被围，預言其被擄耶路撒冷傾陷 34:1-8}
巴比倫王尼布甲尼撒率領他的全軍和地上屬他的各國各邦，攻打耶路撒冷和屬耶路撒冷所有的城邑，那時耶和華的話臨到耶利米說： 2 「耶和華以色列的神說：你去告訴猶大王西底家：『耶和華如此說：我要將這城交付巴比倫王的手，他必用火焚燒。 3 你必不能逃脫他的手，定被拿住交在他的手中。你的眼要見巴比倫王的眼，他要口對口和你說話，你也必到巴比倫去。 4 猶大王西底家啊，你還要聽耶和華的話。耶和華論到你如此說：你必不被刀劍殺死， 5 你必平安而死。人必為你焚燒物件，好像為你列祖，就是在你以前的先王焚燒一般。人必為你舉哀說：「哀哉，我主啊！」耶和華說：這話是我說的。』」

6 於是，先知耶利米在耶路撒冷將這一切話告訴猶大王西底家。 7 那時巴比倫王的軍隊正攻打耶路撒冷，又攻打猶大所剩下的城邑，就是拉吉和亞西加，原來猶大的堅固城只剩下這兩座。

\subsubsection{宣告希伯來僕婢自由反悔，神宣告百姓自由於刀劍/饑荒/瘟疫之下 34:9-22}
\begin{table}[H]
\centering
\begin{tabular}{p{3.75cm}|p{5.0cm}|p{7.70cm}}\hline\hline
\multicolumn{1}{c|}{首領反悔希伯來僕婢自由}&\multicolumn{1}{c|}{责百姓立約反悔}&\multicolumn{1}{c}{宣告百姓自由於刀劍、饑荒、瘟疫之下}\\\hline
\multirow{1}{3.95cm}{西底家王與耶路撒冷的眾民立約，要向他們宣告自由，叫各人任他希伯來的僕人和婢女自由出去，誰也不可使他的一個猶大弟兄做奴僕。（此後，有耶和華的話臨到耶利米。）所有立約的首領和眾民就任他的僕人婢女自由出去，誰也不再叫他們做奴僕。大家都順從，將他們釋放了。後來卻又反悔，叫所任去自由的僕人婢女回來，勉強他們仍為僕婢。}
&因此耶和華的話臨到耶利米說:「耶和華以色列的神如此說：我將你們的列祖從埃及地為奴之家領出來的時候,與他們立約說:『你的一個希伯來弟兄若賣給你,服侍你六年,到第七年你們各人就要任他自由出去。』只是你們列祖不聽從我,也不側耳而聽。如今你們回轉,行我眼中看為正的事,各人向鄰舍宣告自由,並且在稱為我名下的殿中在我面前立約。你們卻又反悔,褻瀆我的名,各人叫所任去隨意自由的僕人婢女回來,勉強他們仍為僕婢。
&所以耶和華如此說:你們沒有聽從我，各人向弟兄、鄰舍宣告自由，看哪，我向你們宣告一樣自由，就是使你們自由於\textcolor{red}{刀劍、饑荒、瘟疫}之下，並且使你們在天下萬國中拋來拋去。這是耶和華說的。猶大的首領，耶路撒冷的首領，太監、祭司和國中的眾民曾將牛犢劈開，分成兩半，從其中經過，在我面前立約。後來又違背我的約，不遵行這約上的話。我必將他們交在仇敵和尋索其命的人手中，他們的屍首必給空中的飛鳥和地上的野獸做食物。並且我必將猶大王西底家和他的首領，交在他們仇敵和尋索其命的人與那暫離你們而去巴比倫王軍隊的手中。耶和華說：我必吩咐他們回到這城，攻打這城，將城攻取，用火焚燒。我也要使猶大的城邑變為荒場，無人居住。\\\hline
\end{tabular}
\end{table}

\subsection{圣城被毁 39:1-10/52:1-30}
\subsubsection{9年10月圍城,11年4月9日城破,剜西底家之目,擄到巴比倫 39:1-10}
\begin{table}[H]
\centering
\begin{tabular}{p{5.25cm}|p{7.0cm}|p{5.0cm}}\hline\hline
\multicolumn{1}{c|}{9年10月圍,11年4月9日破}&\multicolumn{1}{c|}{殺其眾子剜西底家之目}&\multicolumn{1}{c}{逃民擄到巴比倫窮人留猶大}\\\hline
猶大王西底家第九年十月，巴比倫王尼布甲尼撒率領全軍來圍困耶路撒冷。西底家十一年四月初九日，城被攻破。耶路撒冷被攻取的時候，巴比倫王的首領尼甲沙利薛、三甲尼波、撒西金、拉撒力、尼甲沙利薛、拉墨並巴比倫王其餘的一切首領，都來坐在中門。
&猶大王西底家和一切兵丁看見他們，就在夜間從靠近王園兩城中間的門出城逃跑，往亞拉巴逃去。迦勒底的軍隊追趕他們，在耶利哥的平原追上西底家，將他拿住，帶到哈馬地的利比拉巴比倫王尼布甲尼撒那裡，尼布甲尼撒就審判他。巴比倫王在利比拉，西底家眼前殺了他的眾子，又殺了猶大的一切貴胄。並且剜西底家的眼睛，用銅鏈鎖著他，要帶到巴比倫去。
&迦勒底人用火焚燒王宮和百姓的房屋，又拆毀耶路撒冷的城牆。那時，護衛長尼布撒拉旦將城裡所剩下的百姓和投降他的逃民，以及其餘的民都擄到巴比倫去了。護衛長尼布撒拉旦卻將民中毫無所有的窮人留在猶大地，當時給他們葡萄園和田地。\\\hline
\end{tabular}
\end{table}

\subsubsection{西底家做王十一年 52:1-3}
西底家登基的時候年二十一歲，在耶路撒冷做王十一年。他母親名叫哈慕她，是立拿人耶利米的女兒。 2 西底家行耶和華眼中看為惡的事，是照約雅敬一切所行的。 3 因此耶和華的怒氣在耶路撒冷和猶大發作，以致將人民從自己的面前趕出。

\subsubsection{西底家王十一年四月初九日，耶路撒冷淪陷 52:4-9}
西底家背叛巴比倫王。他做王第九年十月初十日，巴比倫王尼布甲尼撒率領全軍來攻擊耶路撒冷，對城安營，四圍築壘攻城。 5 於是城被圍困，直到西底家王十一年。 6 四月初九日，城裡有大饑荒，甚至百姓都沒有糧食。 7 城被攻破，一切兵丁就在夜間從靠近王園兩城中間的門出城逃跑。迦勒底人正在四圍攻城，他們就往亞拉巴逃去。 8 迦勒底的軍隊追趕西底家王，在耶利哥的平原追上他；他的全軍都離開他四散了。 9 迦勒底人就拿住王，帶他到哈馬地的利比拉巴比倫王那裡。巴比倫王便審判他。

\subsubsection{巴比倫王剜西底家之目戮其眾子 52:10-11}
巴比倫王在西底家眼前殺了他的眾子，又在利比拉殺了猶大的一切首領。 11 並且剜了西底家的眼睛，用銅鏈鎖著他，帶到巴比倫去，將他囚在監裡，直到他死的日子。

\subsubsection{五月初十焚毀城垣，留民中最窮的修理葡萄園種田 52:12-16}
巴比倫王尼布甲尼撒十九年五月初十日，在巴比倫王面前侍立的護衛長尼布撒拉旦進入耶路撒冷， 13 用火焚燒耶和華的殿和王宮，又焚燒耶路撒冷的房屋，就是各大戶家的房屋。 14 跟從護衛長迦勒底的全軍就拆毀耶路撒冷四圍的城牆。 15 那時，護衛長尼布撒拉旦將民中最窮的和城裡所剩下的百姓，並已經投降巴比倫王的人，以及大眾所剩下的人，都擄去了。 16 但護衛長尼布撒拉旦留下些民中最窮的，使他們修理葡萄園，耕種田地。

\subsubsection{聖殿被掠 52:17-23}
耶和華殿的銅柱並殿內的盆座和銅海，迦勒底人都打碎了，將那銅運到巴比倫去了。 18 又帶去鍋、鏟子、蠟剪、盤子、調羹，並所用的一切銅器。 19 杯、火鼎、碗、盆、燈臺、調羹、爵，無論金的銀的，護衛長也都帶去了。 20 所羅門為耶和華殿所造的兩根銅柱、一個銅海並座下的十二隻銅牛，這一切的銅，多得無法可稱。 21 這一根柱子高十八肘，厚四指，是空的，圍十二肘。 22 柱上有銅頂，高五肘，銅頂的周圍有網子和石榴，都是銅的。那一根柱子照此一樣，也有石榴。 23 柱子四面有九十六個石榴，在網子周圍共有一百石榴。

\subsubsection{祭司與官吏遭戮 52:24-30}
護衛長拿住大祭司西萊雅、副祭司西番亞和三個把門的， 25 又從城中拿住一個管理兵丁的官(太監)，並在城裡所遇常見王面的七個人和檢點國民軍長的書記，以及城裡所遇見的國民六十個人。 26 護衛長尼布撒拉旦將這些人帶到利比拉巴比倫王那裡。 27 巴比倫王就把他們擊殺在哈馬地的利比拉。這樣，猶大人被擄去離開本地。

28 尼布甲尼撒所擄的民數記在下面：在他第七年擄去猶大人三千零二十三名， 29 尼布甲尼撒十八年從耶路撒冷擄去八百三十二人， 30 尼布甲尼撒二十三年護衛長尼布撒拉旦擄去猶大人七百四十五名，共有四千六百人。



\section{耶路撒冷被攻破以后的领袖 39:11-18/40-44}
\subsection{基大利 39:11-18/40/41:1-15}
\subsubsection{巴比倫王命善遇耶利米，耶和華許搭救以伯米勒 39:11-18}
\begin{table}[H]
\centering
\begin{tabular}{p{8.75cm}|p{9.0cm}}\hline\hline
\multicolumn{1}{c|}{巴比倫王命善遇耶利米}&\multicolumn{1}{c}{耶和華許搭救以伯米勒}\\\hline
巴比倫王尼布甲尼撒提到耶利米,囑咐護衛長尼布撒拉旦說:「你領他去,好好地看待他,切不可害他。他對你怎麼說,你就向他怎麼行。」護衛長尼布撒拉旦和尼布沙斯班,拉撒力,尼甲沙利薛,拉墨並巴比倫王的一切官長,打發人去,將耶利米從護衛兵院中提出來,交於沙番的孫子,亞希甘的兒子基大利帶回家去。於是耶利米住在民中。
&耶利米還囚在護衛兵院中的時候，耶和華的話臨到他說：「你去告訴古實人以伯米勒說：『萬軍之耶和華以色列的神如此說：我說降禍不降福的話必臨到這城，到那時必在你面前成就了。耶和華說：到那日，我必拯救你，你必不致交在你所怕的人手中。我定要搭救你，你不致倒在刀下，卻要以自己的命為掠物，因你倚靠我。這是耶和華說的。」\\\hline
\end{tabular}
\end{table}



\subsubsection{護衛長厚待耶利米,諸軍長眾人聚見基大利 40:1-12}
\begin{table}[H]
\centering
\begin{tabular}{p{8.5cm}|p{8.8cm}}\hline\hline
\multicolumn{1}{c|}{護衛長厚待耶利米}&\multicolumn{1}{c}{諸軍長眾人聚見基大利}\\\hline
耶利米鎖在耶路撒冷和猶大被擄到巴比倫的人中，護衛長尼布撒拉旦將他從拉瑪釋放以後，耶和華的話臨到耶利米。 2 護衛長將耶利米叫來，對他說：「耶和華你的神曾說要降這禍於此地， 3 耶和華使這禍臨到，照他所說的行了。因為你們得罪耶和華，沒有聽從他的話，所以這事臨到你們。 4 現在，我解開你手上的鏈子。你若看與我同往巴比倫去好，就可以去，我必厚待你。你若看與我同往巴比倫去不好，就不必去。看哪，全地在你面前，你以為哪裡美好，哪裡合宜，只管上那裡去吧。」
&\multirow{2}{9cm}{在田野的一切軍長和屬他們的人，聽見巴比倫王立了亞希甘的兒子基大利做境內的省長，並將沒有擄到巴比倫的男人、婦女、孩童和境內極窮的人全交給他， 8 於是軍長尼探雅的兒子以實瑪利、加利亞的兩個兒子約哈難和約拿單、單戶篾的兒子西萊雅，並尼陀法人以斐的眾子、瑪迦人的兒子耶撒尼亞和屬他們的人，都到米斯巴見基大利。 9 沙番的孫子、亞希甘的兒子基大利向他們和屬他們的人起誓說：「不要怕服侍迦勒底人，只管住在這地服侍巴比倫王，就可以得福。 10 至於我，我要住在米斯巴，伺候那到我們這裡來的迦勒底人；只是你們當積蓄酒、油和夏天的果子，收在器皿裡，住在你們所占的城邑中。」 11 在摩押地和亞捫人中，在以東地和各國的一切猶大人，聽見巴比倫王留下些猶大人，並立沙番的孫子、亞希甘的兒子基大利管理他們， 12 這一切猶大人就從所趕到的各處回來，到猶大地的米斯巴基大利那裡，又積蓄了許多的酒並夏天的果子。}\\\cline{1-1}
耶利米還沒有回去，護衛長說：「你可以回到沙番的孫子、亞希甘的兒子基大利那裡去，現在巴比倫王立他做猶大城邑的省長，你可以在他那裡住在民中。不然，你看哪裡合宜，就可以上哪裡去。」於是護衛長送他糧食和禮物，釋放他去了。 6 耶利米就到米斯巴見亞希甘的兒子基大利，在他那裡住在境內剩下的民中。&\\\hline
\end{tabular}
\end{table}


 
\subsubsection{約哈難報以實瑪利之謀，基大利闻而不信被殺 40:13-16/41:1-4}
加利亞的兒子約哈難和在田野的一切軍長來到米斯巴見基大利， 14 對他說：「亞捫人的王巴利斯打發尼探雅的兒子以實瑪利來要你的命，你知道嗎？」亞希甘的兒子基大利卻不信他們的話。 15 加利亞的兒子約哈難在米斯巴私下對基大利說：「求你容我去殺尼探雅的兒子以實瑪利，必無人知道。何必讓他要你的命，使聚集到你這裡來的猶大人都分散，以致猶大剩下的人都滅亡呢？」 16 亞希甘的兒子基大利對加利亞的兒子約哈難說：「你不可行這事，你所論以實瑪利的話是假的。」

七月間，王的大臣宗室以利沙瑪的孫子、尼探雅的兒子以實瑪利帶著十個人，來到米斯巴見亞希甘的兒子基大利，他們在米斯巴一同吃飯。 2 尼探雅的兒子以實瑪利和同他來的那十個人起來，用刀殺了沙番的孫子、亞希甘的兒子基大利，就是巴比倫王所立為全地省長的。 3 以實瑪利又殺了在米斯巴基大利那裡的一切猶大人和所遇見的迦勒底兵丁。4 他殺了基大利，無人知道。

\subsubsection{以實瑪利殺示劍/示羅/撒馬利亞來的70(共80)人逃往亞捫,約哈難救擄民 41:5-15}
 5 第二天，有八十人從示劍和示羅並撒馬利亞來，鬍鬚剃去，衣服撕裂，身體劃破，手拿素祭和乳香，要奉到耶和華的殿。 6 尼探雅的兒子以實瑪利出米斯巴迎接他們，隨走隨哭，遇見了他們，就對他們說：「你們可以來見亞希甘的兒子基大利。」 7 他們到了城中，尼探雅的兒子以實瑪利和同著他的人就將他們殺了，拋在坑中。 8 只是他們中間有十個人對以實瑪利說：「不要殺我們，因為我們有許多大麥、小麥、油、蜜藏在田間。」於是他住了手，沒有將他們殺在弟兄中間。 9 以實瑪利將所殺之人的屍首都拋在坑裡，基大利的旁邊。這坑是從前亞撒王因怕以色列王巴沙所挖的，尼探雅的兒子以實瑪利將那些被殺的人填滿了坑。 

%\subsubsection{以實瑪利擄百姓要逃往亞捫 41:10}
10 以實瑪利將米斯巴剩下的人，就是眾公主和仍住在米斯巴所有的百姓（原是護衛長尼布撒拉旦交給亞希甘的兒子基大利的）都擄去了。尼探雅的兒子以實瑪利擄了他們，要往亞捫人那裡去。

%\subsubsection{約哈難救回所擄之民，以實瑪利和八人逃往亞捫 41:11-15}
加利亞的兒子約哈難和同著他的眾軍長聽見尼探雅的兒子以實瑪利所行的一切惡， 12 就帶領眾人前往，要和尼探雅的兒子以實瑪利爭戰，在基遍的大水旁(大水池旁)遇見他。 13 以實瑪利那裡的眾人看見加利亞的兒子約哈難和同著他的眾軍長，就都歡喜。 14 這樣，以實瑪利從米斯巴所擄去的眾人都轉身歸加利亞的兒子約哈難去了。 15 尼探雅的兒子以實瑪利和八個人脫離約哈難的手，逃往亞捫人那裡去了。

\subsection{約哈難 41:16-18/42-44}
\subsubsection{約哈難居金罕要下埃及,求耶利米諮諏神 41:16-18/42:1-6}


%\subsubsection{約哈難求耶利米諮諏神 42:1-6}
\begin{table}[H]
\centering
\begin{tabular}{p{10.5cm}|p{7.cm}}\hline\hline
\multicolumn{1}{c|}{約哈難並眾百姓求耶利米諮諏神41:16-18//42:1-3}&\multicolumn{1}{c}{百姓願照耶和華的話行}\\\hline
\multirow{2}{10.8cm}{尼探雅的兒子以實瑪利殺了亞希甘的兒子基大利,從米斯巴將剩下的一切百姓兵丁婦女孩童太監擄到基遍之後,加利亞的兒子約哈難和同著他的眾軍長將他們都奪回來帶到靠近伯利恆的金罕寓(基羅特金罕)住下,要進入埃及去。因為尼探雅的兒子以實瑪利殺了巴比倫王所立為省長的亞希甘的兒子基大利,約哈難懼怕迦勒底人。眾軍長和加利亞的兒子約哈難,並何沙雅的兒子耶撒尼亞(亞撒利雅43:2),以及眾百姓從最小的到至大的都進前來,對先知耶利米說:求你准我們在你面前祈求,為我們這剩下的人禱告耶和華你的神。我們本來眾多,現在剩下的極少,這是你親眼所見的。願耶和華你的神指示我們所當走的路,所當做的事。}
&先知耶利米對他們說"我已經聽見你們了,我必照著你們的話禱告耶和華你們的神。耶和華無論回答什麼,我必都告訴你們,毫不隱瞞"\\\cline{2-2}
&於是他們對耶利米說:「我們若不照耶和華你的神差遣你來說的一切話行,願耶和華在我們中間做真實誠信的見證。我們現在請你到耶和華我們的神面前,他說的無論是好是歹,我們都必聽從。我們聽從耶和華我們神的話,就可以得福。」\\\hline
\end{tabular}
\end{table} 


\subsubsection{過十天示眾百姓居猶大得安逸,往埃及必遭諸災 42:7-22}
\begin{table}[H]
\centering
\begin{tabular}{p{4.8cm}|p{12.50cm}}\hline\hline
\multicolumn{1}{c|}{居猶大得享安逸}&\multicolumn{1}{c}{責其不聽警教,往埃及必遭諸災無一存留}\\\hline
\multirow{2}{5.cm}{過了十天，耶和華的話臨到耶利米。他就將加利亞的兒子約哈難和同著他的眾軍長，並眾百姓從最小的到至大的，都叫了來，對他們說：「耶和華以色列的神，就是你們請我在他面前為你們祈求的主，如此說：你們若仍住在這地，我就建立你們必不拆毀，栽植你們並不拔出，因我為降於你們的災禍後悔了。不要怕你們所怕的巴比倫王。耶和華說：不要怕他，因為我與你們同在，要拯救你們脫離他的手。我也要使他發憐憫，好憐憫你們，叫你們歸回本地。}
&倘若你們說『我們不住在這地』，以致不聽從耶和華你們神的話，說：『我們不住這地，卻要進入埃及地，在那裡看不見爭戰，聽不見角聲，也不致無食飢餓。我們必住在那裡。』你們所剩下的猶大人哪，現在要聽耶和華的話！萬軍之耶和華以色列的神如此說：你們若定意要進入埃及，在那裡寄居，你們所懼怕的刀劍在埃及地必追上你們，你們所懼怕的饑荒在埃及要緊緊地跟隨你們，你們必死在那裡。凡定意要進入埃及在那裡寄居的，必遭刀劍、饑荒、瘟疫而死，無一人存留，逃脫我所降於他們的災禍。\\\cline{2-2}
&萬軍之耶和華以色列的神如此說：我怎樣將我的怒氣和憤怒傾在耶路撒冷的居民身上，你們進入埃及的時候，我也必照樣將我的憤怒傾在你們身上，以致你們令人辱罵、驚駭、咒詛、羞辱，你們不得再見這地方。所剩下的猶大人哪，耶和華論到你們說：不要進入埃及去！你們要確實地知道，我今日警教你們了。你們用詭詐自害，因為你們請我到耶和華你們的神那裡，說：『求你為我們禱告耶和華我們神，照耶和華我們的神一切所說的告訴我們，我們就必遵行。』我今日將這話告訴你們，耶和華你們的神為你們的事差遣我到你們那裡說的，你們卻一樣沒有聽從。現在你們要確實地知道：你們在所要去寄居之地必遭刀劍、饑荒、瘟疫而死。\\\hline
\end{tabular}
\end{table}
 
\subsubsection{約哈難不信耶利米預言,攜耶利米與民入埃及地答比匿 43:1-7}
\begin{table}[H]
\centering
\begin{tabular}{p{8.5cm}|p{9.0cm}}\hline\hline
\multicolumn{1}{c|}{約哈難不信耶利米預言1-3}&\multicolumn{1}{c}{攜耶利米與民入埃及地答比匿4-7}\\\hline
耶利米向眾百姓說完了耶和華他們神的一切話，就是耶和華他們神差遣他去所說的一切話， 2 何沙雅的兒子亞撒利雅和加利亞的兒子約哈難，並一切狂傲的人就對耶利米說：「你說謊言！耶和華我們的神並沒有差遣你來說：『你們不可進入埃及，在那裡寄居。』 3 這是尼利亞的兒子巴錄挑唆你害我們，要將我們交在迦勒底人的手中，使我們有被殺的，有被擄到巴比倫去的。」
&於是加利亞的兒子約哈難和一切軍長並眾百姓，不聽從耶和華的話住在猶大地。加利亞的兒子約哈難和一切軍長卻將所剩下的猶大人，就是從被趕到各國回來，在猶大地寄居的男人、婦女、孩童和眾公主，並護衛長尼布撒拉旦所留在沙番的孫子、亞希甘的兒子基大利那裡的眾人，與先知耶利米以及尼利亞的兒子巴錄， 6  7 都帶入埃及地，到了答比匿。這是因他們不聽從耶和華的話。\\\hline
\end{tabular}
\end{table}

\subsubsection{預言埃及必為巴比倫王攻擊 43:8-13}
\begin{table}[H]
\centering
\begin{tabular}{p{8.85cm}|p{8.65cm}}\hline\hline
\multicolumn{1}{c|}{在所藏的石頭上安置巴比倫王寶座8-10}&\multicolumn{1}{c}{巴比倫王來攻擊埃及擄掠刀殺火11-13}\\\hline
在答比匿,耶和華的話臨到耶利米說：「你在猶大人眼前要用手拿幾塊大石頭藏在砌磚的灰泥中,就是在答比匿法老的宮門那裡,對他們說：『萬軍之耶和華以色列的神如此說：我必召我的僕人巴比倫王尼布甲尼撒來,在所藏的石頭上,我要安置他的寶座,他必將光華的寶帳支搭在其上。
&他要來攻擊埃及地，定為死亡的必致死亡，定為擄掠的必被擄掠，定為刀殺的必交刀殺。我要在埃及神的廟中使火著起，巴比倫王要將廟宇焚燒，神像擄去。他要得(披上)埃及地，好像牧人披上外衣，從那裡安然而去。他必打碎埃及地伯示麥的柱像，用火焚燒埃及神的廟宇。』」\\\hline
\end{tabular}
\end{table}

\subsubsection{在猶大降災因民不聽先知燒香奉別神,責民在埃及仍侍他神必受災禍 44:1-14}
\begin{table}[H]
\centering
\begin{tabular}{p{6cm}|p{11.250cm}}\hline\hline
\multicolumn{1}{c|}{在猶大降災因民不聽先知燒香奉別神1-6}&\multicolumn{1}{c}{責民在埃及仍侍他神必受災禍7-14}\\\hline
\multirow{2}{6.26cm}{有臨到耶利米的話，論及一切住在埃及地的猶大人，就是住在密奪、答比匿、挪弗、巴忒羅境內的猶大人，說： 2 「萬軍之耶和華以色列的神如此說：我所降於耶路撒冷和猶大各城的一切災禍，你們都看見了。那些城邑今日荒涼，無人居住。 3 這是\textcolor{red}{因}居民所行的惡，去\textcolor{blue}{燒香侍奉別神}，就是他們和你們並你們列祖所不認識的神，惹我發怒。 4 我從早起來差遣我的僕人眾先知去說：『你們切不要行我所厭惡這可憎之事！』 5 他們卻不聽從，不側耳而聽，不轉離惡事，\textcolor{red}{仍}\textcolor{blue}{向別神燒香}。 6 因此我的怒氣和憤怒都倒出來，在猶大城邑中和耶路撒冷的街市上如火著起，以致都荒廢淒涼，正如今日一樣。}
&現在耶和華萬軍之神以色列的神如此說：你們為何做這大惡自害己命，使你們的男人、婦女、嬰孩和吃奶的都從猶大中剪除不留一人呢？ 8 就是因你們手所做的，在所去寄居的埃及地向別神燒香，惹我發怒，使你們被剪除，在天下萬國中令人咒詛、羞辱。你們列祖的惡行，猶大列王和他們后妃的惡行，你們自己和你們妻子的惡行，就是在猶大地、耶路撒冷街上所行的，你們都忘了嗎？ 10 到如今還沒有懊悔，沒有懼怕，沒有遵行我在你們和你們列祖面前所設立的法度、律例。\\\cline{2-2}
&所以萬軍之耶和華以色列的神如此說：我必向你們變臉降災，以致剪除猶大眾人。 12 那定意進入埃及地在那裡寄居的，就是所剩下的猶大人，我必使他們盡都滅絕。必在埃及地仆倒，必因刀劍、饑荒滅絕，從最小的到至大的都必遭刀劍、饑荒而死，以致令人辱罵、驚駭、咒詛、羞辱。 13 我怎樣用刀劍、饑荒、瘟疫刑罰耶路撒冷，也必照樣刑罰那些住在埃及地的猶大人。 14 甚至那進入埃及地寄居的，就是所剩下的猶大人，都不得逃脫，也不得存留歸回猶大地——他們心中甚想歸回居住之地。除了逃脫的以外，一個都不能歸回。\\\hline
\end{tabular}
\end{table}

\subsubsection{民不聽預言執意仍拜天后 44:15-19}
那些住在埃及地巴忒羅，知道自己妻子向別神燒香的，與旁邊站立的眾婦女，聚集成群，回答耶利米說： 16 「論到你奉耶和華的名向我們所說的話，我們必不聽從！ 17 我們定要成就我們口中所出的一切話，向天后燒香，澆奠祭，按著我們與我們列祖、君王、首領在猶大的城邑中和耶路撒冷的街市上素常所行的一樣。因為那時我們吃飽飯，享福樂，並不見災禍。 18 自從我們停止向天后燒香，澆奠祭，我們倒缺乏一切，又因刀劍、饑荒滅絕。」

婦女說：「我們向天后燒香，澆奠祭，做天后像的餅供奉她，向她澆奠祭，是外乎我們的丈夫嗎？」

\subsubsection{耶利米警民聽耶和華的話,否则必遭諸災埃及必見敗於巴比倫 44:20-30}
耶利米對一切那樣回答他的男人、婦女說： 21 「你們與你們列祖、君王、首領並國內的百姓，在猶大城邑中和耶路撒冷街市上所燒的香，耶和華豈不記念，心中豈不思想嗎？ 22 耶和華因你們所做的惡、所行可憎的事，不能再容忍，所以你們的地荒涼，令人驚駭、咒詛，無人居住，正如今日一樣。 23 你們燒香得罪耶和華，沒有聽從他的話，沒有遵行他的律法、條例、法度，所以你們遭遇這災禍，正如今日一樣。」

24 耶利米又對眾民和眾婦女說：「你們在埃及地的一切猶大人，當聽耶和華的話！\textcolor{red}{萬軍之耶和華以色列的神如此說}：你們和你們的妻都口中說、手裡做，說：『我們定要償還所許的願，向天后燒香，澆奠祭。』現在你們只管堅定所許的願而償還吧！ 26 所以你們住在埃及地的一切猶大人，當聽耶和華的話！

\textcolor{red}{耶和華說}：我指著我的大名起誓：在埃及全地，我的名不再被猶大一個人的口稱呼說『我指著主永生的耶和華起誓』。 27 我向他們留意降禍不降福，在埃及地的一切猶大人必因刀劍、饑荒所滅，直到滅盡。 28 脫離刀劍從埃及地歸回猶大地的，人數很少。那進入埃及地要在那裡寄居的，就是所剩下的猶大人，必知道是誰的話立得住，是我的話呢，是他們的話呢。

\textcolor{red}{耶和華說}：我在這地方刑罰你們必有預兆，使你們知道我降禍於你們的話必要立得住。 30 耶和華如此說：我必將埃及王法老合弗拉交在他仇敵和尋索其命的人手中，像我將猶大王西底家交在他仇敵和尋索其命的巴比倫王尼布甲尼撒手中一樣。



\section{神對列國的審判 46-51}
\subsection{預言法老 46}
\subsubsection{預言法老必敗於幼發拉底河濱 46:1-12}
耶和華論列國的話臨到先知耶利米。

2 論到關乎埃及王法老尼哥的軍隊，這軍隊安營在幼發拉底河邊的迦基米施，是巴比倫王尼布甲尼撒在猶大王約西亞的兒子約雅敬第四年所打敗的：

3 「你們要預備大小盾牌，往前上陣！ 4 你們套上車，騎上馬！頂盔站立，磨槍貫甲！ 5 我為何看見他們驚惶轉身退後呢？他們的勇士打敗了，急忙逃跑，並不回頭，驚嚇四圍都有。」這是耶和華說的。 6 「不要容快跑的逃避，不要容勇士逃脫(快跑的不能逃避，勇士不能逃脫)，他們在北方幼發拉底河邊絆跌仆倒。 7 像尼羅河漲發，像江河之水翻騰的，是誰呢？ 8 埃及像尼羅河漲發，像江河的水翻騰，他說：『我要漲發遮蓋遍地，我要毀滅城邑和其中的居民。』 9 馬匹上去吧！車輛急行吧！勇士，就是手拿盾牌的古實人和弗人(呂彼亞人)，並拉弓的路德族，都出去吧！ 10 那日是主萬軍之耶和華報仇的日子，要向敵人報仇。刀劍必吞吃得飽，飲血飲足，因為主萬軍之耶和華在北方幼發拉底河邊有獻祭的事。 11 埃及的民(處女)哪，可以上基列取乳香去！你雖多服良藥，總是徒然，不得治好。 12 列國聽見你的羞辱，遍地滿了你的哀聲，勇士與勇士彼此相碰，一齊跌倒。」

\subsubsection{預言巴比倫王必擊埃及 46:13-26}
耶和華對先知耶利米所說的話，論到巴比倫王尼布甲尼撒要來攻擊埃及地：

14 「你們要傳揚在埃及，宣告在密奪，報告在挪弗、答比匿，說：『要站起出隊，自做準備，因為刀劍在你四圍施行吞滅的事。』 15 你的壯士為何被沖去呢？他們站立不住，因為耶和華驅逐他們。 16 使多人絆跌，他們也彼此撞倒，說：『起來吧！我們再往本民本地去，好躲避欺壓的刀劍。』 17 他們在那裡喊叫說：『埃及王法老不過是個聲音(已經敗亡)，他已錯過所定的時候了。』」 18 君王，名為萬軍之耶和華的說：「我指著我的永生起誓：尼布甲尼撒(他)來的勢派必像他泊在眾山之中，像迦密在海邊一樣。 19 住在埃及的民(女子)哪，要預備擄去時所用的物件！因為挪弗必成為荒場，且被燒毀無人居住。 20 埃及是肥美的母牛犢，但出於北方的毀滅(牛虻)來到了，來到了！ 21 其中的雇勇好像圈裡的肥牛犢，他們轉身退後，一齊逃跑，站立不住，因為他們遭難的日子、追討的時候已經臨到。 22 其中的聲音好像蛇行一樣，敵人要成隊而來，如砍伐樹木的手拿斧子攻擊她。」 23 耶和華說：「埃及的樹林雖然不能尋察(穿不過)，敵人卻要砍伐，因他們多於蝗蟲，不可勝數。 24 埃及的民(女子)必然蒙羞，必交在北方人的手中。」 25 萬軍之耶和華以色列的神說：「我必刑罰挪的亞捫（埃及尊大之神）和法老，並埃及與埃及的神以及君王，也必刑罰法老和倚靠他的人。 26 我要將他們交付尋索其命之人的手和巴比倫王尼布甲尼撒與他臣僕的手，以後埃及必再有人居住，與從前一樣。」這是耶和華說的。

\subsubsection{言明雅各受罰之後必復蒙恩 46:27-28}
「我的僕人雅各啊，不要懼怕；以色列啊，不要驚惶！因我要從遠方拯救你，從被擄到之地拯救你的後裔。雅各必回來，得享平靖安逸，無人使他害怕。 28 我的僕人雅各啊，不要懼怕，因我與你同在。我要將我所趕你到的那些國滅絕淨盡，卻不將你滅絕淨盡。倒要從寬懲治你，萬不能不罰你(以你為無罪)。」這是耶和華說的。

%\subsection{預言非利士人 46}
\subsection{預言非利士人必滅 47}
法老攻擊加沙之先，有耶和華論非利士人的話臨到先知耶利米。

2 耶和華如此說：「有水從北方發起，成為漲溢的河，要漲過遍地和其中所有的，並城和其中所住的。人必呼喊，境內的居民都必哀號。 3 聽見敵人壯馬蹄跳的響聲和戰車隆隆、車輪轟轟，為父的手就發軟，不回頭看顧兒女。 4 因為日子將到，要毀滅一切非利士人，剪除幫助推羅、西頓所剩下的人。原來耶和華必毀滅非利士人，就是迦斐托海島餘剩的人。 5 加沙成了光禿，平原中所剩的亞實基倫歸於無有。你用刀劃身，要到幾時呢？ 6 耶和華的刀劍哪，你到幾時才止息呢？你要入鞘，安靖不動！ 7 耶和華既吩咐你攻擊亞實基倫和海邊之地，他已經派定你，焉能止息呢？」

\subsection{預言摩押 48}
\subsubsection{預言摩押必受刀劍懲罰 48:1-6}
論摩押。萬軍之耶和華以色列的神如此說：「尼波有禍了，因變為荒場。基列亭蒙羞被攻取，米斯迦蒙羞被毀壞。 2 摩押不再被稱讚，有人在希實本設計謀害她說：『來吧！我們將她剪除，不再成國。』瑪得緬哪，你也必默默無聲，刀劍必追趕你。 3 從何羅念有喊荒涼大毀滅的哀聲。 4 摩押毀滅了，她的孩童(家童)發哀聲，使人聽見。 5 人上魯希坡隨走隨哭，因為在何羅念的下坡聽見毀滅的哀聲。 6 你們要奔逃，自救性命，獨自居住，好像曠野的杜松。

\subsubsection{因自恃財貨 48:7-10}
「你因倚靠自己所做的和自己的財寶，必被攻取，基抹和屬他的祭司、首領也要一同被擄去。 8 行毀滅的必來到各城，並無一城得免。山谷必致敗落，平原必被毀壞，正如耶和華所說的。 9 要將翅膀給摩押，使她可以飛去。她的城邑必致荒涼，無人居住。 10 懶惰為耶和華行事的，必受咒詛；禁止刀劍不經血的，必受咒詛。

\subsubsection{因妄賴勇力 48:11-25}

「摩押自幼年以來常享安逸，如酒在渣滓上澄清，沒有從這器皿倒在那器皿裡，也未曾被擄去。因此她的原味尚存，香氣未變。」 12 耶和華說：「日子將到，我必打發倒酒的往她那裡去，將她倒出來，倒空她的器皿，打碎她的罈子。 13 摩押必因基抹羞愧，像以色列家從前倚靠伯特利的神羞愧一樣。 14 你們怎麼說『我們是勇士，是有勇力打仗的』呢？ 15 摩押變為荒場，敵人上去進了她的城邑，她所特選的少年人下去遭了殺戮。」這是君王，名為萬軍之耶和華說的。 16 「摩押的災殃臨近，她的苦難速速來到。 17 凡在她四圍的和認識她名的，你們都要為她悲傷，說：『那結實的杖和那美好的棍何竟折斷了呢！』 18 住在底本的民(女子)哪，要從你榮耀的位上下來，坐受乾渴，因毀滅摩押的上來攻擊你，毀壞了你的保障。 19 住亞羅珥的啊，要站在道旁觀望，問逃避的男人和逃脫的女人說：『是什麼事呢？』 20 摩押因毀壞蒙羞，你們要哀號呼喊，要在亞嫩旁報告說：『摩押變為荒場！』 21 刑罰臨到平原之地的何倫、雅雜、米法押、 22 底本、尼波、伯低比拉太音、 23 基列亭、伯迦末、伯米恩、 24 加略、波斯拉，和摩押地遠近所有的城邑。 25 摩押的角砍斷了，摩押的膀臂折斷了。」這是耶和華說的。

\subsubsection{因悖逆神輕侮主民 48:26-35}
「你們要使摩押沉醉，因她向耶和華誇大。她要在自己所吐之中打滾，又要被人嗤笑。 27 摩押啊，你不曾嗤笑以色列嗎？她豈是在賊中查出來的呢？你每逢提到她便搖頭。 28 摩押的居民哪，要離開城邑，住在山崖裡，像鴿子在深淵口上搭窩。 29 我們聽說摩押人驕傲，是極其驕傲；聽說她自高自傲，並且狂妄，居心自大。」 30 耶和華說：「我知道她的憤怒是虛空的，她誇大的話一無所成。 31 因此我要為摩押哀號，為摩押全地呼喊，人必為吉珥哈列設人嘆息。 32 西比瑪的葡萄樹啊，我為你哀哭，甚於雅謝人哀哭。你的枝子蔓延過海，直長到雅謝海，那行毀滅的已經臨到你夏天的果子和你所摘的葡萄。 33 肥田和摩押地的歡喜快樂都被奪去，我使酒榨的酒絕流，無人踹酒歡呼，那歡呼卻變為仇敵的呐喊(那歡呼卻不是歡呼)。 34 希實本人發的哀聲達到以利亞利，直達到雅雜，從瑣珥達到何羅念，直到伊基拉施利施亞，因為寧林的水必然乾涸。」 35 耶和華說：「我必在摩押地使那在丘壇獻祭的和那向她的神燒香的，都斷絕了。

\subsubsection{仇敵必攻擊摩押，摩押必被毀滅不再成國 48:36-44}
「我心腹為摩押哀鳴如簫，我心腸為吉珥哈列設人也是如此，因摩押人所得的財物都滅沒了。 37 各人頭上光禿，鬍鬚剪短，手有劃傷，腰束麻布。 38 在摩押的各房頂上和街市上處處有人哀哭，因我打碎摩押，好像打碎無人喜悅的器皿。」這是耶和華說的。 39 「摩押何等毀壞，何等哀號，何等羞愧轉背！這樣，摩押必令四圍的人嗤笑、驚駭。」 40 耶和華如此說：「仇敵必如大鷹飛起，展開翅膀攻擊摩押。 41 加略被攻取，保障也被占據，到那日，摩押的勇士心中疼痛如臨產的婦人。 42 摩押必被毀滅，不再成國，因她向耶和華誇大。」 43 耶和華說：「摩押的居民哪，恐懼、陷坑、網羅都臨近你。 44 躲避恐懼的必墜入陷坑，從陷坑上來的必被網羅纏住，因我必使追討之年臨到摩押。」這是耶和華說的。

\subsubsection{到末後被擄的摩押人歸回 48:45-47}
「躲避的人無力站在希實本的影下，因為有火從希實本發出，有火焰出於西宏的城，燒盡摩押的角和鬨嚷人的頭頂。 46 摩押啊，你有禍了！屬基抹的民滅亡了！因你的眾子都被擄去，你的眾女也被擄去。」 47 耶和華說：「到末後，我還要使被擄的摩押人歸回。」摩押受審判的話到此為止。

\subsection{預言亞捫被火焚燒被趕出被趕出後歸回 49:1-6}
%\subsubsection{預言亞捫人必受懲罰 49:1-5}
論亞捫人。耶和華如此說：「以色列沒有兒子嗎？沒有後嗣嗎？瑪勒堪為何得迦得之地為業呢？屬他的民為何住其中的城邑呢？」 2 耶和華說：「日子將到，我必使人聽見打仗的喊聲，是攻擊亞捫人拉巴的喊聲。拉巴要成為亂堆，屬她的鄉村(女子)要被火焚燒。先前得以色列地為業的，此時以色列倒要得他們的地為業。」這是耶和華說的。 3 「希實本哪，你要哀號！因為愛地變為荒場。拉巴的居民 (女子)哪，要呼喊！以麻布束腰，要哭號，在籬笆中跑來跑去，因瑪勒堪和屬他的祭司、首領要一同被擄去。 4 背道的民(女子)哪，你們為何因有山谷，就是水流的山谷誇張呢？為何倚靠財寶說『誰能來到我們這裡』呢？」 5 主萬軍之耶和華說：「我要使恐嚇從四圍的人中臨到你們，你們必被趕出，各人一直前往，沒有人收聚逃民。
「後來，我還要使被擄的亞捫人歸回。」這是耶和華說的

\subsection{以東因心中的狂傲自欺必受懲罰變為永遠的荒場 49:7-22}
\begin{wrapfigure}{r}{0.6\textwidth}
\centering
\includegraphics[width=0.5\textwidth]{ZongJie/DXianZhiFigs/YeLiMi/Teman}
\end{wrapfigure}
%\subsubsection{以東必受懲罰，變為永遠的荒場 49:7-13}
論以東。萬軍之耶和華如此說：「提幔中再沒有智慧嗎？明哲人不再有謀略嗎？他們的智慧盡歸無有嗎？ 8 底但的居民哪，要轉身逃跑，住在深密處！因為我向以掃追討的時候，必使災殃臨到他。 9 摘葡萄的若來到他那裡，豈不剩下些葡萄呢？盜賊若夜間而來，豈不毀壞直到夠了呢？ 10 我卻使以掃赤露，顯出他的隱密處，他不能自藏。他的後裔、弟兄、鄰舍盡都滅絕，他也歸於無有。 11 你撇下孤兒，我必保全他們的命，你的寡婦可以倚靠我。」 12 耶和華如此說：「原不該喝那杯的一定要喝，你能盡免刑罰嗎？你必不能免，一定要喝。」 13 耶和華說：「我指著自己起誓：波斯拉必令人驚駭、羞辱、咒詛，並且荒涼，她的一切城邑必變為永遠的荒場。」

%\subsubsection{因心中的狂傲自欺，以東必傾覆像所多瑪、蛾摩拉 49:14-22}
我從耶和華那裡聽見信息，並有使者被差往列國去說：「你們聚集來攻擊以東，要起來爭戰！」 15 「我使你在列國中為最小，在世人中被藐視。 16 住在山穴中據守山頂的啊，論到你的威嚇，你因心中的狂傲自欺。你雖如大鷹高高搭窩，我卻從那裡拉下你來。」這是耶和華說的。 17 「以東必令人驚駭，凡經過的人就受驚駭，又因她一切的災禍嗤笑。」 18 耶和華說：「必無人住在那裡，也無人在其中寄居，要像所多瑪、蛾摩拉和鄰近的城邑傾覆的時候一樣。 19 仇敵必像獅子從約旦河邊的叢林上來，攻擊堅固的居所。轉眼之間，我要使以東人逃跑，離開這地。誰蒙揀選，我就派誰治理這地。誰能比我呢？誰能給我定規日期呢？有何牧人能在我面前站立得住呢？」

20 你們要聽耶和華攻擊以東所說的謀略和他攻擊提幔居民所定的旨意：仇敵定要將他們群眾微弱的拉去，定要使他們的居所荒涼。 21 因他們仆倒的聲音，地就震動，人在紅海那裡必聽見呼喊的聲音。 22 仇敵必如大鷹飛起，展開翅膀攻擊波斯拉。到那日，以東的勇士心中疼痛，如臨產的婦人。


\subsection{大馬士革必受懲罰被火燒滅 49:23-27}
\begin{wrapfigure}{r}{0.6\textwidth}
\centering
\includegraphics[width=0.6\textwidth]{ZongJie/DXianZhiFigs/YeLiMi/HaMa}
\end{wrapfigure}
論大馬士革。「哈馬和亞珥拔蒙羞，因他們聽見凶惡的信息就消化了。海上有憂愁，不得平靜。 24 大馬士革發軟，轉身逃跑。戰兢將她捉住，痛苦憂愁將她抓住如產難的婦人一樣。 25 我所喜樂、可稱讚的城，為何被撇棄了呢？ 26 她的少年人必仆倒在街上，當那日一切兵丁必默默無聲。」這是萬軍之耶和華說的。 27 「我必在大馬士革城中使火著起，燒滅便哈達的宮殿。」

\subsection{基達必受懲罰帳篷羊群器皿被掠 49:28-29}
論巴比倫王尼布甲尼撒所攻打的基達和夏瑣的諸國。耶和華如此說：「迦勒底人哪，起來上基達去，毀滅東方人！ 29 他們的帳篷和羊群都要奪去，將幔子和一切器皿並駱駝為自己掠去。人向他們喊著說：『四圍都有驚嚇！』」

\subsection{夏瑣必受懲罰牲畜被擄人被分散住處永遠淒涼 49:30-33}
耶和華說：「夏瑣的居民哪，要逃奔遠方，住在深密處，因為巴比倫王尼布甲尼撒設計謀害你們，起意攻擊你們。」 31 耶和華說：「迦勒底人哪，起來，上安逸無慮的居民那裡去！他們是無門無閂，獨自居住的。 32 他們的駱駝必成為掠物，他們眾多的牲畜必成為擄物。我必將剃周圍頭髮的人分散四方(風)，使災殃從四圍臨到他們。」這是耶和華說的。 33 「夏瑣必成為野狗的住處，永遠淒涼。必無人住在那裡，也無人在其中寄居。」

\subsection{以攔人被刀劍追殺分散末後歸回 49:34-39}
猶大王西底家登基的時候，耶和華論以攔的話臨到先知耶利米說： 35 「萬軍之耶和華如此說：我必折斷以攔人的弓，就是他們為首的權力。 36 我要使四風從天的四方颳來臨到以攔人，將他們分散四方(風)，這被趕散的人沒有一國不到的。 37 耶和華說：我必使以攔人在仇敵和尋索其命的人面前驚惶。我也必使災禍，就是我的烈怒臨到他們。又必使刀劍追殺他們，直到將他們滅盡。 38 我要在以攔設立我的寶座，從那裡除滅君王和首領。這是耶和華說的。

「到末後，我還要使被擄的以攔人歸回。這是耶和華說的。」

\subsection{預言巴比倫 50-51}
\subsubsection{一國從北方上來攻擊巴比倫，以色列猶大返錫安尋主 50:1-8}
\begin{table}[H]
\centering
\begin{tabular}{p{9.45cm}|p{8.0cm}}\hline\hline
\multicolumn{1}{c|}{巴比倫和迦勒底人之地}&\multicolumn{1}{c}{以色列人和猶大人}\\\hline
耶和華藉先知耶利米論巴比倫和迦勒底人之地所說的話。「你們要在萬國中傳揚報告，豎立大旗，要報告，不可隱瞞，說：『巴比倫被攻取，彼勒蒙羞，米羅達驚惶。巴比倫的神像都蒙羞，她的偶像都驚惶。』因有一國從北方上來攻擊她，使她的地荒涼，無人居住，連人帶牲畜都逃走了。」
&耶和華說：「當那日子，那時候，以色列人要和猶大人同來，隨走隨哭，尋求耶和華他們的神。他們必訪問錫安，又面向這裡說：『來吧，你們要與耶和華聯合，為永遠不忘的約。』我的百姓做了迷失的羊，牧人使他們走差路，使他們轉到山上。他們從大山走到小山，竟忘了安歇之處。 \\\hline

凡遇見他們的就把他們吞滅，敵人說：『我們沒有罪，因他們得罪那做公義居所的耶和華，就是他們列祖所仰望的耶和華。』 
&我民哪，你們要從巴比倫中逃走，從迦勒底人之地出去，要像羊群前面走的公山羊。\\\hline
\end{tabular}
\end{table}


\subsubsection{申言聯合的大國從北方上來攻擊巴比倫 50:8-20}
\begin{table}[H]
\centering
\begin{tabular}{p{5.75cm}|p{5.50cm}|p{5.50cm}}\hline\hline
\multicolumn{1}{c|}{聯合的大國從北方上來}&\multicolumn{1}{c|}{巴比倫}&\multicolumn{1}{c}{以色列如羊歸牧}\\\hline
因我必激動聯合的大國從北方上來攻擊巴比倫，他們要擺陣攻擊她，她必從那裡被攻取。他們的箭好像善射之勇士的箭，一枝也不徒然返回。 
&迦勒底必成為掠物，凡擄掠她的都必心滿意足。」這是耶和華說的。「搶奪我產業的啊，你們因歡喜快樂，且像踹穀撒歡的母牛犢，又像發嘶聲的壯馬， 12 你們的母巴比倫就極其抱愧，生你們的必然蒙羞。她要列在諸國之末，成為曠野、旱地、沙漠。因耶和華的憤怒，必無人居住，要全然荒涼。凡經過巴比倫的要受驚駭，又因她所遭的災殃嗤笑。
&\\\hline

所有拉弓的，你們要在巴比倫的四圍擺陣，射箭攻擊她，不要愛惜箭枝，因她得罪了耶和華。你們要在她四圍呐喊，她已經投降，外郭坍塌了，城牆拆毀了，因為這是耶和華報仇的事。你們要向巴比倫報仇，她怎樣待人，也要怎樣待她。你們要將巴比倫撒種的，和收割時拿鐮刀的都剪除了。
&他們各人因怕欺壓的刀劍，必歸回本族，逃到本土。
&以色列是打散的羊，是被獅子趕出的，首先是亞述王將他吞滅，末後是巴比倫王尼布甲尼撒將他的骨頭折斷。\\\hline

所以萬軍之耶和華以色列的神如此說：「我必罰巴比倫王和他的地，像我從前罰亞述王一樣。&
&我必再領以色列回他的草場，他必在迦密和巴珊吃草，又在以法蓮山上和基列境內得以飽足。耶和華說：「當那日子，那時候，雖尋以色列的罪孽，一無所有；雖尋猶大的罪惡，也無所見。因為我所留下的人，我必赦免。」\\\hline
\end{tabular}
\end{table}


\subsubsection{再言攻擊米拉大翁，比割,申以色列人和猶大人的冤 50:21-34}
\begin{table}[H]
\centering
\begin{tabular}{p{12.cm}|p{5.750cm}}\hline\hline
\multicolumn{1}{c|}{巴比倫之地}&\multicolumn{1}{c}{以色列和猶大}\\\hline
耶和華說：「上去攻擊米拉大翁之地，又攻擊比割的居民，要追殺滅盡，照我一切所吩咐你的去行。境內有打仗和大毀滅的響聲。全地的大錘何竟砍斷破壞！巴比倫在列國中何竟荒涼！巴比倫哪，我為你設下網羅，你不知不覺被纏住。你被尋著，也被捉住，因為你與耶和華爭競。耶和華已經開了武庫，拿出他惱恨的兵器，因為主萬軍之耶和華在迦勒底人之地有當做的事。你們要從極遠的邊界來攻擊她，開她的倉廩，將她堆如高堆，毀滅淨盡，絲毫不留。要殺她的一切牛犢，使他們下去遭遇殺戮。他們有禍了！因為追討他們的日子已經來到。 
&有從巴比倫之地逃避出來的人，在錫安揚聲報告：『耶和華我們的神報仇，就是為他的殿報仇。』\\\hline
招集一切弓箭手來攻擊巴比倫，要在巴比倫四圍安營，不要容一人逃脫，照著她所做的報應她。她怎樣待人，也要怎樣待她，因為她向耶和華以色列的聖者發了狂傲。所以她的少年人必仆倒在街上，當那日一切兵丁必默默無聲。」這是耶和華說的。主萬軍之耶和華說：「你這狂傲的啊，我與你反對，因為我追討你的日子已經來到。狂傲的必絆跌仆倒，無人扶起。我也必使火在她的城邑中著起來，將她四圍所有的盡行燒滅。
&萬軍之耶和華如此說：「以色列人和猶大人一同受欺壓，凡擄掠他們的都緊緊抓住他們，不肯釋放。他們的救贖主大有能力，萬軍之耶和華是他的名，他必申清他們的冤，好使全地得平安，並攪擾巴比倫的居民。」\\\hline
\end{tabular}
\end{table}



\subsubsection{復言巴比倫遭北方大國許多君王從地極,約旦河叢林來的刀劍,弓,槍報 50:35-46}
耶和華說：「有刀劍臨到迦勒底人和巴比倫的居民，並她的首領與智慧人。 36 有刀劍臨到矜誇的人，他們就成為愚昧；有刀劍臨到她的勇士，他們就驚惶。 37 有刀劍臨到她的馬匹、車輛和其中雜族的人民，他們必像婦女一樣；有刀劍臨到她的寶物，就被搶奪。 38 有乾旱臨到她的眾水，就必乾涸，因為這是有雕刻偶像之地，人因偶像而顛狂。 39 所以，曠野的走獸和豺狼必住在那裡，鴕鳥也住在其中，永無人煙，世世代代無人居住。」 40 耶和華說：「必無人住在那裡，也無人在其中寄居，要像我傾覆所多瑪、蛾摩拉和鄰近的城邑一樣。

41 「看哪，有一種民從北方而來，並有一大國和許多君王被激動，從地極來到。 42 他們拿弓和槍，性情殘忍，不施憐憫，他們的聲音像海浪砰訇。巴比倫城(女子)啊，他們騎馬，都擺隊伍，如上戰場的人，要攻擊你。 43 巴比倫王聽見他們的風聲，手就發軟，痛苦將他抓住，疼痛彷彿產難的婦人。 44 仇敵必像獅子從約旦河邊的叢林上來，攻擊堅固的居所。轉眼之間，我要使他們逃跑，離開這地。誰蒙揀選，我就派誰治理這地。誰能比我呢？誰能給我定規日期呢？有何牧人能在我們面前站立得住呢？」 45 你們要聽耶和華攻擊巴比倫所說的謀略，和他攻擊迦勒底人之地所定的旨意：仇敵定要將他們群眾微弱的拉去，定要使他們的居所荒涼。 46 因巴比倫被取的聲音，地就震動，人在列邦都聽見呼喊的聲音。

\subsubsection{耶和華必降重災於巴比倫以復其虐待以色列人之仇 51:1-14}

\begin{table}[H]
\begin{tabular}{p{9.75cm}|p{7.50cm}}\hline\hline
\multicolumn{1}{c|}{門徒}&\multicolumn{1}{c}{耶穌}\\\hline
耶和華如此說：「我必使毀滅的風颳起，攻擊巴比倫和住在立加米\footnotemark{}的人。 2 我要打發外邦人來到巴比倫，簸揚她，使她的地空虛，在她遭禍的日子他們要周圍攻擊她。 3 拉弓的要向拉弓的和貫甲挺身的射箭，不要憐惜她的少年人，要滅盡她的全軍。 4 他們必在迦勒底人之地被殺仆倒，在巴比倫的街上被刺透。」
&以色列和猶大雖然境內充滿違背以色列聖者的罪，卻沒有被他的神萬軍之耶和華丟棄。 6 你們要從巴比倫中逃奔，各救自己的性命！不要陷在她的罪孽中一同滅亡，因為這是耶和華報仇的時候，他必向巴比倫施行報應。 \\\hline
巴比倫素來是耶和華手中的金杯，使天下沉醉，萬國喝了她的酒就顛狂了。 8 巴比倫忽然傾覆毀壞，要為她哀號，為止她的疼痛，拿乳香或者可以治好。 9 我們想醫治巴比倫，她卻沒有治好。離開她吧！
&我們各人歸回本國，因為她受的審判通於上天，達到穹蒼。 10 耶和華已經彰顯我們的公義，來吧，我們可以在錫安報告耶和華我們神的作為。\\\hline
\end{tabular}
\end{table}
\footnotetext{迦勒底כשדים，此係希伯來人對迦勒底所稱之暗語，是把巴比倫文的迦勒底一字的字母 ，再換成希伯文的倒數字母，而拼成的新字。 קָמָ֑י לֵ֣ב} 


11 你們要磨尖了箭頭，抓住盾牌。耶和華定意攻擊巴比倫，將她毀滅，所以激動了瑪代君王的心，因這是耶和華報仇，就是為自己的殿報仇。 12 你們要豎立大旗攻擊巴比倫的城牆，要堅固瞭望臺，派定守望的設下埋伏。因為耶和華指著巴比倫居民所說的話、所定的意，他已經做成。 13 住在眾水之上多有財寶的啊，你的結局到了，你貪婪之量滿了！ 14 萬軍之耶和華指著自己起誓說：「我必使敵人充滿你，像螞蚱一樣，他們必呐喊攻擊你。」


\subsubsection{偶像虛無必被毀滅 51:15-18}
耶和華用能力創造大地，用智慧建立世界，用聰明鋪張穹蒼。 16 他一發聲，空中便有多水激動，他使雲霧從地極上騰。他造電隨雨而閃，從他府庫中帶出風來。 17 各人都成了畜類，毫無知識，各銀匠都因他的偶像羞愧。他所鑄的偶像本是虛假的，其中並無氣息， 18 都是虛無的，是迷惑人的工作，到追討的時候，必被除滅。 

\subsubsection{雅各是神爭戰的兵器打碎列國 51:19-24}
19 雅各的份不像這些，因他是造做萬有的主，以色列也是他產業的支派，萬軍之耶和華是他的名。20 「你是我爭戰的斧子和打仗的兵器。我要用你打碎列國，用你毀滅列邦， 21 用你打碎馬和騎馬的，用你打碎戰車和坐在其上的， 22 用你打碎男人和女人，用你打碎老年人和少年人，用你打碎壯丁和處女， 23 用你打碎牧人和他的群畜，用你打碎農夫和他一對牛，用你打碎省長和副省長。」 24 耶和華說：「我必在你們眼前，報復巴比倫人和迦勒底居民在錫安所行的諸惡。」

\subsubsection{巴比倫在列國中何竟變為荒場 51:25-44}
25 耶和華說：「你這行毀滅的山哪，就是毀滅天下的山，我與你反對。我必向你伸手，將你從山巖滾下去，使你成為燒毀的山。 26 人必不從你那裡取石頭為房角石，也不取石頭為根基石，你必永遠荒涼。」這是耶和華說的。

要在境內豎立大旗，在各國中吹角，使列國預備攻擊巴比倫，將亞拉臘、米尼、亞實基拿各國招來攻擊她。又派軍長來攻擊她，使馬匹上來如螞蚱。 28 使列國和瑪代君王，與省長和副省長，並他們所管全地之人都預備攻擊她。 29 地必震動而瘠苦，因耶和華向巴比倫所定的旨意成立了，使巴比倫之地荒涼，無人居住。 30 巴比倫的勇士止息爭戰，藏在堅壘之中，他們的勇力衰盡，好像婦女一樣。巴比倫的住處有火著起，門閂都折斷了。 31 跑報的要彼此相遇，送信的要互相迎接，報告巴比倫王說：「城的四方被攻取了， 32 渡口被占據了，葦塘被火燒了，兵丁也驚慌了。」

33 萬軍之耶和華以色列的神如此說：「巴比倫城(女子)好像踹穀的禾場，再過片時，收割她的時候就到了。」 34 以色列人說：「巴比倫王尼布甲尼撒吞滅我，壓碎我，使我成為空虛的器皿。他像大魚將我吞下，用我的美物充滿他的肚腹，又將我趕出去。」 35 錫安的居民要說：「巴比倫以強暴待我，損害我的身體，願這罪歸給她！」耶路撒冷人要說：「願流我們血的罪歸到迦勒底的居民！」 36 所以，耶和華如此說：「我必為你申冤，為你報仇。我必使巴比倫的海枯竭，使她的泉源乾涸。 37 巴比倫必成為亂堆，為野狗的住處，令人驚駭、嗤笑，並且無人居住。 38 他們要像少壯獅子咆哮，像小獅子吼叫。 39 他們火熱的時候，我必為他們設擺酒席，使他們沉醉，好叫他們快樂，睡了長覺永不醒起。」這是耶和華說的。 40 「我必使他們像羊羔，像公綿羊和公山羊下到宰殺之地。

41 「示沙克(巴比倫)何竟被攻取！天下所稱讚的何竟被占據！巴比倫在列國中何竟變為荒場！ 42 海水漲起漫過巴比倫，她被許多海浪遮蓋。 43 她的城邑變為荒場、旱地、沙漠，無人居住、無人經過之地。 44 我必刑罰巴比倫的彼勒，使他吐出所吞的。萬民必不再流歸他那裡，巴比倫的城牆也必坍塌了。


\subsubsection{警民從其中出去，躲避耶和華的烈怒 51:45-51}
45 「我的民哪，你們要從其中出去，各人拯救自己，躲避耶和華的烈怒。 46 你們不要心驚膽怯，也不要因境內所聽見的風聲懼怕。因為這年有風聲傳來，那年也有風聲傳來，境內有強暴的事，官長攻擊官長。 47 日子將到，我必刑罰巴比倫雕刻的偶像，她全地必然抱愧，她被殺的人必在其中仆倒。 48 那時天地和其中所有的必因巴比倫歡呼，因為行毀滅的要從北方來到她那裡。」這是耶和華說的。 49 巴比倫怎樣使以色列被殺的人仆倒，照樣，她全地被殺的人也必在巴比倫仆倒。

你們躲避刀劍的要快走，不要站住。要在遠方記念耶和華，心中追想耶路撒冷。 51 我們聽見辱罵就蒙羞，滿面慚愧，因為外邦人進入耶和華殿的聖所。 

\subsubsection{大毀滅的仇敵響聲從迦勒底發出，巴比倫變為荒場 51:52-58}
52 耶和華說：「日子將到，我必刑罰巴比倫雕刻的偶像，通國受傷的人必唉哼。 53 巴比倫雖升到天上，雖使她堅固的高處更堅固，還有行毀滅的從我這裡到她那裡。」這是耶和華說的。

54 有哀號的聲音從巴比倫出來，有大毀滅的響聲從迦勒底人之地發出。 55 因耶和華使巴比倫變為荒場，使其中的大聲滅絕。仇敵彷彿眾水，波浪砰訇，響聲已經發出。 56 這是行毀滅的臨到巴比倫，巴比倫的勇士被捉住，他們的弓折斷了，因為耶和華是施行報應的神，必定施行報應。 57 君王，名為萬軍之耶和華的說：「我必使巴比倫的首領、智慧人、省長、副省長和勇士都沉醉，使他們睡了長覺永不醒起。」 58 萬軍之耶和華如此說：「巴比倫寬闊的城牆必全然傾倒，她高大的城門必被火焚燒。眾民所勞碌的必致虛空，列國所勞碌的被火焚燒，他們都必困乏。」

\subsubsection{西底家四年上巴比倫，耶利米吩咐念將書之災於巴比倫 51:59-62}
猶大王西底家在位第四年上巴比倫去的時候，瑪西雅的孫子、尼利亞的兒子西萊雅與王同去，西萊雅是王宮的大臣，先知耶利米有話吩咐他。 60 耶利米將一切要臨到巴比倫的災禍，就是論到巴比倫的一切話，寫在書上。 61 耶利米對西萊雅說：「你到了巴比倫，務要念這書上的話。 62 又說：『耶和華啊，你曾論到這地方說要剪除，甚至連人帶牲畜沒有在這裡居住的，必永遠荒涼。』

\subsubsection{命投於幼發拉底河以示巴比倫必沉淪 51:63-54}
「你念完了這書，就把一塊石頭拴在書上，扔在幼發拉底河中， 64 說：『巴比倫因耶和華所要降於她的災禍必如此沉下去，不再興起，人民也必困乏。』」

耶利米的話到此為止。














